% Generated mechanically from frozen input p07/ja/local-cohomology.tex.
% Do not edit this generated file; edit the bound locale source instead.

% OK, start here.
%

\setcounter{chapter}{50}
\chapter{局所コホモロジー}



\phantomsection
\label{local-cohomology-section-phantom}


\section{序論}
\label{local-cohomology-section-introduction}

\noindent
本章では局所コホモロジーの研究を続ける。
参考文献として \cite{SGA2} がある。
局所コホモロジーの定義は Dualizing Complexes, 第
\ref{dualizing-section-local-cohomology} 節にある。
Noether 環について局所コホモロジーを取ることは、適切な捩れ関手を導来することと
同じである。これは Dualizing Complexes, 第
\ref{dualizing-section-local-cohomology-noetherian} 節で示されている。
深さとの関係については Dualizing Complexes, 第
\ref{dualizing-section-depth} 節を参照されたい。

\medskip\noindent
局所コホモロジーの有限性を論じ、それを用いて Grothendieck の有限性定理の
かなり一般的な形を証明する。定理 \ref{local-cohomology-theorem-finiteness} および
補題 \ref{local-cohomology-lemma-finiteness-Rjstar}（開埋め込みに沿う連接加群の高次順像）を
参照されたい。我々の方法には、Faltings の論文に見られるいくつかの非常に
洗練された議論を取り入れる。\cite{Faltings-annulators} および
\cite{Faltings-finiteness} を参照されたい。

\medskip\noindent
応用として Hartshorne--Lichtenbaum 消滅を論じる。また、局所コホモロジーに
対する Frobenius および微分作用素の作用も扱う。




\section{一般事項}
\label{local-cohomology-section-generalities}

\noindent
次の補題は、関手 $R\Gamma_Z$ が台付きコホモロジーと関係することを述べる。

\begin{lemma}
\label{local-cohomology-lemma-local-cohomology-is-local-cohomology}
$A$ を環、$I$ を有限生成イデアルとする。
$Z = V(I) \subset X = \Spec(A)$ とおく。$K \in D(A)$ が
Derived Categories of Schemes, 補題
\ref{perfect-lemma-affine-compare-bounded} を通じて
$\widetilde{K} \in D_\QCoh(\mathcal{O}_X)$ に対応するとき、関手的同型
$$
R\Gamma_Z(K) = R\Gamma_Z(X, \widetilde{K})
$$
がある。左辺は Dualizing Complexes, 式
(\ref{dualizing-equation-local-cohomology}) のものであり、右辺は
Cohomology, 第 \ref{cohomology-section-cohomology-support-bis} 節
の関手である。
\end{lemma}

\begin{proof}
Cohomology, 補題 \ref{cohomology-lemma-triangle-sections-with-support}
により、区別三角形
$$
R\Gamma_Z(X, \widetilde{K})
\to R\Gamma(X, \widetilde{K})
\to R\Gamma(U, \widetilde{K})
\to R\Gamma_Z(X, \widetilde{K})[1]
$$
が存在する。ここで $U = X \setminus Z$ である。Derived Categories of
Schemes, 補題 \ref{perfect-lemma-affine-compare-bounded} により
$R\Gamma(X, \widetilde{K}) = K$ であることが分かっている。
$I = (f_1, \ldots, f_r)$ と書く。このとき有限アフィン開被覆
$\mathcal{U} : U = D(f_1) \cup \ldots \cup D(f_r)$ を得る。
Derived Categories of Schemes, 補題
\ref{perfect-lemma-alternating-cech-complex-complex-computes-cohomology}
により、交代 {\v C}ech 複体
$\text{Tot}(\check{\mathcal{C}}_{alt}^\bullet(\mathcal{U},
\widetilde{K^\bullet}))$
は $R\Gamma(U, \widetilde{K})$ を計算する。ここで $K^\bullet$ は任意の
$A$-加群の複体であって $K$ を表す。定義を展開すると
$$
R\Gamma(U, \widetilde{K}) =
\text{Tot}\left(
K^\bullet \otimes_A
(\prod\nolimits_{i_0} A_{f_{i_0}} \to
\prod\nolimits_{i_0 < i_1} A_{f_{i_0}f_{i_1}} \to
\ldots \to A_{f_1\ldots f_r})\right)
$$
を得る。
$K^\bullet = R\Gamma(X, \widetilde{K^\bullet}) \to
R\Gamma(U, \widetilde{K}^\bullet)$ は、$A$ から $\prod A_{f_i}$ への
対角写像によって誘導されることが明らかである。したがって
$$
R\Gamma_Z(X, \widetilde{K}) =
\text{Tot}\left(
K^\bullet \otimes_A
(A \to \prod\nolimits_{i_0} A_{f_{i_0}} \to
\prod\nolimits_{i_0 < i_1} A_{f_{i_0}f_{i_1}} \to
\ldots \to A_{f_1\ldots f_r})\right)
$$
を得る。Dualizing Complexes, 補題
\ref{dualizing-lemma-local-cohomology-adjoint} により、右辺の複体は
$R\Gamma_Z(K)$ を計算するので、与えられた $K$ に対して補題が成り立つ。

\medskip\noindent
最後に、この同型を $K$ について関手的に選べることを示す。
$A$-加群の複体の写像
$$
M^\bullet = \text{Tot}\left(
K^\bullet \otimes_A
(A \to \prod\nolimits_{i_0} A_{f_{i_0}} \to
\prod\nolimits_{i_0 < i_1} A_{f_{i_0}f_{i_1}} \to
\ldots \to A_{f_1\ldots f_r})\right)
\longrightarrow
K^\bullet
$$
を考える。上の議論により、この矢印は関手
$R\Gamma_Z(X, \widetilde{\ })$ を適用すると同型になり、また
$R\Gamma(U, \widetilde{M}^\bullet) = 0$ である（詳細は省略する）。すると
$$
R\Gamma_Z(X, \widetilde{K}^\bullet)
\leftarrow
R\Gamma_Z(X, \widetilde{M}^\bullet)
\rightarrow
R\Gamma(X, \widetilde{M}^\bullet) = M^\bullet = R\Gamma_Z(K^\bullet)
$$
を得る。すべての写像は $D(A)$ における同型であり、所望どおり
$K^\bullet$ について関手的である。
\end{proof}

\begin{lemma}
\label{local-cohomology-lemma-local-cohomology}
$A$ を環、$I \subset A$ を有限生成イデアルとする。
$X = \Spec(A)$、$Z = V(I)$、$U = X \setminus Z$ とおき、$j : U \to X$
を包含射とする。$\mathcal{F}$ を準連接 $\mathcal{O}_U$-加群とする。このとき
\begin{enumerate}
\item $A$-加群 $M$ が存在し、$\mathcal{F}$ は $\widetilde{M}$ の $U$ への
制限である。
\item $M$ が与えられると完全列
$$
0 \to H^0_Z(M) \to M \to H^0(U, \mathcal{F}) \to H^1_Z(M) \to 0
$$
および同型 $H^p(U, \mathcal{F}) = H^{p + 1}_Z(M)$ が $p \geq 1$
に対してある。
\item $M = H^0(U, \mathcal{F})$ と取ることができ、この場合
$H^0_Z(M) = H^1_Z(M) = 0$ である。
\end{enumerate}
\end{lemma}

\begin{proof}
$M$ の存在は Properties, 補題 \ref{properties-lemma-extend-trivial} と、
$X$ 上の準連接層が $A$-加群に対応するという事実（Schemes, 補題
\ref{schemes-lemma-equivalence-quasi-coherent}）から従う。
次に区別三角形
$$
R\Gamma_Z(X, \widetilde{M}) \to R\Gamma(X, \widetilde{M}) \to
R\Gamma(U, \widetilde{M}|_U) \to R\Gamma_Z(X, \widetilde{M})[1]
$$
を考える。これは Cohomology, 補題
\ref{cohomology-lemma-triangle-sections-with-support} の三角形である。
$X$ はアフィンなので、Cohomology of Schemes, 補題
\ref{coherent-lemma-quasi-coherent-affine-cohomology-zero} により
$R\Gamma(X, \widetilde{M}) = M$ である。$M$ の選び方から
$\mathcal{F} = \widetilde{M}|_U$ であり、したがって完全列
$$
0 \to H^0_Z(X, \widetilde{M}) \to M \to H^0(U, \mathcal{F}) \to
H^1_Z(X, \widetilde{M}) \to 0
$$
および同型 $H^p(U, \mathcal{F}) = H^{p + 1}_Z(X, \widetilde{M})$ を
$p \geq 1$ に対して得る。
補題 \ref{local-cohomology-lemma-local-cohomology-is-local-cohomology} により
$H^i_Z(M) = H^i_Z(X, \widetilde{M})$ であり、これはすべての $i$ に対して
成り立つ。したがって
(1) と (2) が成り立つ。
最後に $M' = H^0(U, \mathcal{F})$ とおくと、$M \to M'$ の核と余核は
$I$-冪捩れである。よって $\widetilde{M}|_U \to \widetilde{M'}|_U$ は
同型であり、(3) のとおり $M'$ を用いることができる。もちろん、このとき
$H^0_Z(M')$ と $H^0_Z(M')$ はともに零となる。
\end{proof}

\begin{lemma}
\label{local-cohomology-lemma-already-torsion}
$I, J \subset A$ を環 $A$ の有限生成イデアルとする。
$M$ が $I$-冪捩れ加群ならば、標準写像
$$
H^i_{V(I) \cap V(J)}(M) \to H^i_{V(J)}(M)
$$
はすべての $i$ に対して同型である。
\end{lemma}

\begin{proof}
Dualizing Complexes, 補題 \ref{dualizing-lemma-local-cohomology-ss} の
スペクトル系列を用いると、主張は $R\Gamma_I(M) = M$ に帰着する。
これは Dualizing Complexes, 第
\ref{dualizing-section-local-cohomology} 節における局所コホモロジーの
構成から直ちに従う。
\end{proof}

\begin{lemma}
\label{local-cohomology-lemma-multiplicative}
$S \subset A$ を環 $A$ の乗法的集合とする。
$M$ を $A$-加群とし、$S^{-1}M = 0$ と仮定する。このとき
$\colim_{f \in S} H^0_{V(f)}(M) = M$ および
$\colim_{f \in S} H^1_{V(f)}(M) = 0$ である。
\end{lemma}

\begin{proof}
$H^0$ に関する主張は定義から直ちに従う。$H^1$ に関する主張については、
$R\Gamma_{V(f)}$ と $H^1_{V(f)}$ が余極限と可換であることに注意する。
したがって、$M$ はある $f \in S$ により零化されると仮定してよい。このとき
$H^1_{V(ff')}(M) = 0$ がすべての $f' \in S$ に対して成り立つ（例えば
補題 \ref{local-cohomology-lemma-already-torsion} による）。
\end{proof}

\begin{lemma}
\label{local-cohomology-lemma-elements-come-from-bigger}
$I \subset A$ を環 $A$ の有限生成イデアルとする。$\mathfrak p$ を素イデアル、
$M$ を $A$-加群とする。整数 $i \geq 0$ に対して写像
$$
\Psi :
\colim_{f \in A, f \not \in \mathfrak p} H^i_{V((I, f))}(M)
\longrightarrow
H^i_{V(I)}(M)
$$
を考える。このとき
\begin{enumerate}
\item $\Im(\Psi)$ は $H^i_{V(I)}(M)_\mathfrak p$ において零に写る元の集合である。
\item $H^{i - 1}_{V(I)}(M)_\mathfrak p = 0$ ならば $\Psi$ は単射である。
\item $H^{i - 1}_{V(I)}(M)_\mathfrak p = H^i_{V(I)}(M)_\mathfrak p = 0$
ならば $\Psi$ は同型である。
\end{enumerate}
\end{lemma}

\begin{proof}
$f \in A$、$f \not \in \mathfrak p$ に対して、Dualizing Complexes,
補題 \ref{dualizing-lemma-local-cohomology-ss} のスペクトル系列は退化し、
短完全列
$$
0 \to H^1_{V(f)}(H^{i - 1}_{V(I)}(M)) \to
H^i_{V((I, f))}(M) \to H^0_{V(f)}(H^i_{V(I)}(M)) \to 0
$$
を与える。これで (1) が証明され、(2) はこのことと
補題 \ref{local-cohomology-lemma-multiplicative} から従う。(3) は形式的な帰結である。
\end{proof}

\begin{lemma}
\label{local-cohomology-lemma-isomorphism}
$I \subset I' \subset A$ を Noether 環 $A$ の有限生成イデアルとする。
$M$ を $A$-加群とし、$i \geq 0$ を整数とする。写像
$$
\Psi : H^i_{V(I')}(M) \to H^i_{V(I)}(M)
$$
を考える。このとき次が成り立つ。
\begin{enumerate}
\item $H^i_{\mathfrak pA_\mathfrak p}(M_\mathfrak p) = 0$ がすべての
$\mathfrak p \in V(I) \setminus V(I')$ に対して成り立つならば、
$\Psi$ は全射である。
\item $H^{i - 1}_{\mathfrak pA_\mathfrak p}(M_\mathfrak p) = 0$ がすべての
$\mathfrak p \in V(I) \setminus V(I')$ に対して成り立つならば、
$\Psi$ は単射である。
\item $H^i_{\mathfrak pA_\mathfrak p}(M_\mathfrak p) =
H^{i - 1}_{\mathfrak pA_\mathfrak p}(M_\mathfrak p) = 0$ がすべての
$\mathfrak p \in V(I) \setminus V(I')$ に対して成り立つならば、
$\Psi$ は同型である。
\end{enumerate}
\end{lemma}

\begin{proof}
(1) を証明する。$\xi \in H^i_{V(I)}(M)$ とする。$A$ は Noether 環なので、
最大のイデアル $I \subset I'' \subset I'$ が存在し、$\xi$ はある
$\xi'' \in H^i_{V(I'')}(M)$ の像である。
$V(I'') = V(I')$ ならば証明は終わる。そうでなければ、生成点
$\mathfrak p \in V(I'')$ で $V(I')$ に属さないものを選ぶ。仮定により
$H^i_{V(I'')}(M)_\mathfrak p =
H^i_{\mathfrak pA_\mathfrak p}(M_\mathfrak p) = 0$ である。
補題 \ref{local-cohomology-lemma-elements-come-from-bigger} により $I''$ を大きくできるが、
これは最大性に反する。

\medskip\noindent
(2) を証明する。$\xi' \in H^i_{V(I')}(M)$ を $\Psi$ の核の元とする。
$A$ は Noether 環なので、最大のイデアル $I \subset I'' \subset I'$ が存在し、
$\xi'$ は $H^i_{V(I'')}(M)$ において零に写る。
$V(I'') = V(I')$ ならば証明は終わる。そうでなければ、生成点
$\mathfrak p  \in V(I'')$ で $V(I')$ に属さないものを選ぶ。仮定により
$H^{i - 1}_{V(I'')}(M)_\mathfrak p =
H^{i - 1}_{\mathfrak pA_\mathfrak p}(M_\mathfrak p) = 0$ である。
補題 \ref{local-cohomology-lemma-elements-come-from-bigger} により $I''$ を大きくできるが、
これは最大性に反する。

\medskip\noindent
(3) は (1) と (2) から形式的に従う。
\end{proof}





\section{Hartshorne の連結性補題}
\label{local-cohomology-section-hartshorne-connectedness}

\noindent
この節の表題は次の結果を指す。

\begin{lemma}
\label{local-cohomology-lemma-depth-2-connected-punctured-spectrum}
\begin{reference}
\cite[命題 2.1]{Hartshorne-connectedness}
\end{reference}
\begin{slogan}
Hartshorne の連結性
\end{slogan}
$A$ を深さ $\geq 2$ の Noether 局所環とする。このとき $A$、$A^h$、
$A^{sh}$ の穴あきスペクトルは連結である。
\end{lemma}

\begin{proof}
$U$ を $A$ の穴あきスペクトルとする。$U$ が非連結ならば、
$\Gamma(U, \mathcal{O}_U)$ は非自明な冪等元をもつ。しかし $A$ は局所環なので
非自明な冪等元をもたない。したがって
$A \to \Gamma(U, \mathcal{O}_U)$ は同型でない。
補題 \ref{local-cohomology-lemma-local-cohomology} により、$H^0_\mathfrak m(A)$ または
$H^1_\mathfrak m(A)$ のいずれかは零でない。ゆえに Dualizing Complexes,
補題 \ref{dualizing-lemma-depth} により $\text{depth}(A) \leq 1$ となる。
$A^h$ と $A^{sh}$ に対する結果には More on Algebra, 補題
\ref{more-algebra-lemma-henselization-depth} を用いればよい。
\end{proof}

\begin{lemma}
\label{local-cohomology-lemma-catenary-S2-equidimensional}
\begin{reference}
\cite[IV, 系 5.10.9]{EGA}
\end{reference}
$A$ を鎖状かつ $(S_2)$ である Noether 局所環とする。このとき
$\Spec(A)$ は等次元である。
\end{lemma}

\begin{proof}
$X = \Spec(A)$ とおき、$d = \dim(A) = \dim(X)$ とする。$X$ の中で、
$X_1$ を次元 $d$ の既約成分の合併、$X_2$ を次元 $< d$ の既約成分の合併
とする。もちろん $X = X_1 \cup X_2$ である。$X_2 = \emptyset$ ならば
補題が成り立つ。そうでなければ $Z = X_1 \cap X_2$ は $X$ の閉点を少なくとも
含むので、$X$ の空でない閉部分集合である。したがって、$z \in Z$ を $Z$ の
ある既約成分の生成点として選べる。$\mathcal{O}_{Z, z}$ のスペクトルは、$X$ の
点で $z$ に特殊化するものの集合であることを思い出そう。$z$ は次元 $d$ の
既約成分と次元 $< d$ の既約成分の両方に含まれるので、非自明な特殊化
$x_1 \leadsto z$ と $x_2 \leadsto z$ を得て、$x_1$ と $x_2$ の閉包の次元は異なる。
$X$ は鎖状なので、これは特殊化 $x_1 \leadsto z$ と $x_2 \leadsto z$ の少なくとも
一方が直近でない場合にしか起こらない。したがって
$\dim(\mathcal{O}_{Z, z}) \geq 2$ である。よって
$\text{depth}(\mathcal{O}_{Z, z}) \geq 2$ である。これは $A$ が $(S_2)$
だからである。しかし $U$ を $\mathcal{O}_{Z, z}$ の穴あきスペクトルとすると
非連結である。実際、閉部分集合
$U \cap X_1$ と $U \cap X_2$ は（$z$ の選び方により）互いに素であり、$U$ を
被覆する。これは 補題 \ref{local-cohomology-lemma-depth-2-connected-punctured-spectrum}
に反し、証明が完了する。
\end{proof}



\section{コホモロジー次元}
\label{local-cohomology-section-cd}

\noindent
コホモロジー次元について簡単に述べる。

\begin{lemma}
\label{local-cohomology-lemma-cd}
$I \subset A$ を環 $A$ の有限生成イデアルとする。
$Y = V(I) \subset X = \Spec(A)$ とおき、$d \geq -1$ を整数とする。
次の条件は同値である。
\begin{enumerate}
\item $H^i_Y(A) = 0$ が $i > d$ に対して成り立つ。
\item $H^i_Y(M) = 0$ が $i > d$ に対して、任意の $A$-加群 $M$ について成り立つ。
\item $d = -1$ ならば $Y = \emptyset$ であり、$d = 0$ ならば
$Y$ は $X$ の開かつ閉な部分集合であり、$d > 0$ ならば
$H^i(X \setminus Y, \mathcal{F}) = 0$ が $i \geq d$ に対して、任意の
準連接 $\mathcal{O}_{X \setminus Y}$-加群 $\mathcal{F}$ について成り立つ。
\end{enumerate}
\end{lemma}

\begin{proof}
$R\Gamma_Y(-)$ のコホモロジー次元が有限であることに注意する。これは例えば
Dualizing Complexes, 補題 \ref{dualizing-lemma-local-cohomology-adjoint}
による。したがって整数 $i_0$ が存在し、$H^i_Y(M) = 0$ がすべての
$A$-加群 $M$ および $i \geq i_0$ に対して成り立つ。

\medskip\noindent
(1) と (2) が同値であることを示す。(2) から (1) が従うことは明らかである。
(1) を仮定する。$i > d$ に関する降下帰納法により、$H^i_Y(M) = 0$ が
すべての $A$-加群 $M$ に対して成り立つことを示す。
$i \geq i_0$ の場合は既に上で示した。帰納段階として $i_0 > i > d$ とする。
任意の $A$-加群 $M$ を取り、これを短完全列
$0 \to N \to F \to M \to 0$ に組み込む。ここで $F$ は自由 $A$-加群である。
$R\Gamma_Y$ は右随伴なので、$H^i_Y(-)$ は直和と可換である。
したがって仮定 (1) により $H^i_Y(F) = 0$ である。実際、$i > d$ である。
すると所望どおり $H^i_Y(M) = H^{i + 1}_Y(N) = 0$ を得る。

\medskip\noindent
$d = -1$ で (2) が成り立つと仮定する。このとき $0 = H^0_Y(A/I) = A/I \Rightarrow A = I
\Rightarrow Y = \emptyset$ である。したがって (3) が成り立つ。逆の証明は省略する。

\medskip\noindent
$d = 0$ で (2) が成り立つと仮定する。
$J = H^0_I(A) = \{x \in A \mid I^nx = 0 \text{ を満たす }n > 0\text{ が存在する}\}$
とおく。このとき
$$
H^1_Y(A) = \Coker(A \to \Gamma(X \setminus Y, \mathcal{O}_{X \setminus Y}))
\quad\text{および}\quad
H^1_Y(I) = \Coker(I \to \Gamma(X \setminus Y, \mathcal{O}_{X \setminus Y}))
$$
であり、最初の写像の核は $J$ に等しい。補題 \ref{local-cohomology-lemma-local-cohomology}
を参照せよ。(2) から $I(A/J) = A/J$ を得る。よって $f \in I$ を選び、
その像が $1$ となるようにできる。ここで像は $A/J$ において取る。
すると $1 - f \in J$ なので、$I^n(1 - f) = 0$ となる $n > 0$ が存在する。
したがって $f^n = f^{n + 1}$ であり、$e = f^n \in I$ は冪等元である。
補元となる冪等元 $e' = 1 - f^n \in J$ を考える。任意の元 $g \in I$ に対して
$g^m e' = 0$ となる $m > 0$ がある。したがって $I$ はイデアル
$(e) \subset I$ の根基に含まれる。これは (3) のとおり
$Y = V(I) = V(e)$ が $X$ において開かつ閉であることを意味する。
逆に $Y = V(I)$ が開かつ閉ならば、関手 $H^0_Y(-)$ は完全であり、
その高次導来関手は消滅する。

\medskip\noindent
$d > 0$ ならば、補題 \ref{local-cohomology-lemma-local-cohomology} から直ちに
(2) と (3) が同値であることが分かる。
\end{proof}

\begin{definition}
\label{local-cohomology-definition-cd}
$I \subset A$ を環 $A$ の有限生成イデアルとする。整数 $d \geq -1$ のうち
補題 \ref{local-cohomology-lemma-cd} の同値な条件を満たす最小のものを、
{\it $I$ の $A$ におけるコホモロジー次元}と呼び、$\text{cd}(A, I)$ と書く。
\end{definition}

\noindent
したがって $\text{cd}(A, I) = -1$ となるのは $I = A$ の場合であり、
$\text{cd}(A, I) = 0$ となるのは $I$ が局所冪零であるか冪等元で生成される場合である。
次の補題により $\text{cd}(A, I)$ が存在することに注意する。

\begin{lemma}
\label{local-cohomology-lemma-bound-cd}
$I \subset A$ を環 $A$ の有限生成イデアルとする。このとき
\begin{enumerate}
\item $\text{cd}(A, I)$ は $I$ の生成元の個数以下である。
\item $\text{cd}(A, I) \leq r$ は、$f_1, \ldots, f_r \in A$ で
$V(f_1, \ldots, f_r) = V(I)$ を満たすものが存在するとき成り立つ。
\item $\text{cd}(A, I) \leq c$ は、$\Spec(A) \setminus V(I)$ が $c$ 個の
アフィン開集合で被覆できるとき成り立つ。
\end{enumerate}
\end{lemma}

\begin{proof}
$R\Gamma_Y(-)$ の明示的な記述が Dualizing Complexes, 補題
\ref{dualizing-lemma-local-cohomology-adjoint} にあり、これから (1) が従う。
(2) は、$R\Gamma_Z$ が閉部分集合 $Z$ のみに依存し、有限生成イデアル
$I \subset A$ で $V(I) = Z$ を満たすものの選択には依存しないことを用いて
(1) から導ける。この事実は、Dualizing Complexes, 第
\ref{dualizing-section-local-cohomology} 節における局所コホモロジーの構成と
More on Algebra, 補題 \ref{more-algebra-lemma-local-cohomology-closed}
を組み合わせるか、補題 \ref{local-cohomology-lemma-local-cohomology-is-local-cohomology} を
用いれば従う。(3) には 補題 \ref{local-cohomology-lemma-cd} と Cohomology of Schemes,
補題 \ref{coherent-lemma-vanishing-nr-affines} の消滅結果を用いる。
\end{proof}

\begin{lemma}
\label{local-cohomology-lemma-cd-sum}
$I, J \subset A$ を環 $A$ の有限生成イデアルとする。このとき
$\text{cd}(A, I + J) \leq \text{cd}(A, I) + \text{cd}(A, J)$ である。
\end{lemma}

\begin{proof}
定義と Dualizing Complexes, 補題 \ref{dualizing-lemma-local-cohomology-ss}
を用いる。
\end{proof}

\begin{lemma}
\label{local-cohomology-lemma-cd-change-rings}
$A \to B$ を環準同型、$I \subset A$ を有限生成イデアルとする。このとき
$\text{cd}(B, IB) \leq \text{cd}(A, I)$ である。$A \to B$ が忠実平坦ならば
等号が成り立つ。
\end{lemma}

\begin{proof}
定義と Dualizing Complexes, 補題 \ref{dualizing-lemma-torsion-change-rings}
を用いる。
\end{proof}

\begin{lemma}
\label{local-cohomology-lemma-cd-local}
$I \subset A$ を環 $A$ の有限生成イデアルとする。このとき
$\text{cd}(A, I) = \max \text{cd}(A_\mathfrak p, I_\mathfrak p)$ である。
\end{lemma}

\begin{proof}
$Y = V(I)$ および $Y' = V(I_\mathfrak p) \subset \Spec(A_\mathfrak p)$ とする。
等式
$R\Gamma_Y(A) \otimes_A A_\mathfrak p = R\Gamma_{Y'}(A_\mathfrak p)$
が Dualizing Complexes, 補題 \ref{dualizing-lemma-torsion-change-rings}
により成り立つことを思い出そう。Algebra, 補題
\ref{algebra-lemma-characterize-zero-local} により結論を得る。
\end{proof}

\begin{lemma}
\label{local-cohomology-lemma-cd-dimension}
$I \subset A$ を環 $A$ の有限生成イデアルとする。$M$ が有限 $A$-加群ならば、
$H^i_{V(I)}(M) = 0$ が $i > \dim(\text{Supp}(M))$ に対して成り立つ。
特に $\text{cd}(A, I) \leq \dim(A)$ である。
\end{lemma}

\begin{proof}
まず第二の主張を証明する。$\dim(A)$ は Krull 次元を表すことを思い出そう。
補題 \ref{local-cohomology-lemma-cd-local} により $A$ は局所環であると仮定してよい。
$V(I) = \emptyset$ ならば結果は成り立つ。$V(I) \not = \emptyset$ ならば、
$\dim(\Spec(A) \setminus V(I)) < \dim(A)$ である。これは閉点を取り除いているからである。
$U = \Spec(A) \setminus V(I)$ はスペクトル空間 $\Spec(A)$ の準コンパクト
開部分集合なので、それ自身スペクトル空間であることに注意する。
Algebra, 補題 \ref{algebra-lemma-spec-spectral} および
Topology, 補題 \ref{topology-lemma-spectral-sub} を参照せよ。
したがって Cohomology, 命題
\ref{cohomology-proposition-cohomological-dimension-spectral} により
$H^i(U, \mathcal{F}) = 0$ が $i \geq \dim(A)$ に対して成り立つ。
補題 \ref{local-cohomology-lemma-cd} により所望の結論が従う。Noether の場合には、読者は
Grothendieck による Cohomology, 命題
\ref{cohomology-proposition-vanishing-Noetherian} を用いてもよい。

\medskip\noindent
第二の主張から第一の主張を導く。$\mathfrak a$ を有限 $A$-加群 $M$ の零化イデアルとし、
$B = A/\mathfrak a$ とおく。$\Spec(B) = \text{Supp}(M)$ であることを思い出そう。
Algebra, 補題 \ref{algebra-lemma-support-closed} を参照せよ。
$J = IB$ とおく。このとき $M$ は $B$-加群であり、
$H^i_{V(I)}(M) = H^i_{V(J)}(M)$ である。Dualizing Complexes, 補題
\ref{dualizing-lemma-local-cohomology-and-restriction} を参照せよ。
$\text{cd}(B, J) \leq \dim(B) = \dim(\text{Supp}(M))$ が前半で示したことから
成り立つので、結論を得る。
\end{proof}

\begin{lemma}
\label{local-cohomology-lemma-cd-is-one}
$I \subset A$ を環 $A$ の有限生成イデアルとする。
$\text{cd}(A, I) = 1$ ならば $\Spec(A) \setminus V(I)$ は空でないアフィンスキームである。
\end{lemma}

\begin{proof}
これは 補題 \ref{local-cohomology-lemma-cd} と Cohomology of Schemes, 補題
\ref{coherent-lemma-quasi-compact-h1-zero-covering} から従う。
\end{proof}

\begin{lemma}
\label{local-cohomology-lemma-cd-maximal}
$(A, \mathfrak m)$ を次元 $d$ の Noether 局所環とする。このとき
$H^d_\mathfrak m(A)$ は零でなく、$\text{cd}(A, \mathfrak m) = d$ である。
\end{lemma}

\begin{proof}
次元の特徴付けの一つにより、$A$ の定義イデアルで $d$ 個の元で生成されるものが
存在する。Algebra, 命題 \ref{algebra-proposition-dimension} を参照せよ。
したがって $\text{cd}(A, \mathfrak m) \leq d$ が 補題 \ref{local-cohomology-lemma-bound-cd}
により成り立つ。よって $H^d_\mathfrak m(A)$ が零でないこと、
$\text{cd}(A, \mathfrak m) = d$ であること、および
$\text{cd}(A, \mathfrak m) \geq d$ であることは互いに同値である。

\medskip\noindent
$A \to A^\wedge$ を $A$ からその完備化への写像とする。
$A^\wedge$ は $A$ と同じ次元の Noether 局所環であり、極大イデアルは
$\mathfrak m A^\wedge$ であることに注意する。Algebra, 補題
\ref{algebra-lemma-completion-Noetherian-Noetherian},
\ref{algebra-lemma-completion-complete}、
\ref{algebra-lemma-completion-faithfully-flat} および
More on Algebra, 補題 \ref{more-algebra-lemma-completion-dimension} を参照せよ。
補題 \ref{local-cohomology-lemma-cd-change-rings} により、$A^\wedge$ に対して補題を示せば十分である。

\medskip\noindent
前段落により $A$ は完備局所環であると仮定してよい。このとき $A$ は
正規化された双対化複体 $\omega_A^\bullet$ をもつ（Dualizing Complexes,
補題 \ref{dualizing-lemma-ubiquity-dualizing}）。局所双対定理
（Dualizing Complexes, 補題 \ref{dualizing-lemma-special-case-local-duality}
の形）により、$H^d_\mathfrak m(A)$ は
$\text{Ext}^{-d}(A, \omega_A^\bullet) = H^{-d}(\omega_A^\bullet)$ の Matlis 双対である。
後者は例えば Dualizing Complexes, 補題
\ref{dualizing-lemma-nonvanishing-generically-local} により零でない。
\end{proof}

\begin{lemma}
\label{local-cohomology-lemma-cd-bound-dim-local}
$(A, \mathfrak m)$ を Noether 局所環とする。$I \subset A$ を真のイデアル、
$\mathfrak p \subset A$ を素イデアルで
$V(\mathfrak p) \cap V(I) = \{\mathfrak m\}$ を満たすものとする。このとき
$\dim(A/\mathfrak p) \leq \text{cd}(A, I)$ である。
\end{lemma}

\begin{proof}
補題 \ref{local-cohomology-lemma-cd-change-rings} により
$\text{cd}(A, I) \geq \text{cd}(A/\mathfrak p, I(A/\mathfrak p))$ である。
$V(I) \cap V(\mathfrak p) = \{\mathfrak m\}$ なので
$\text{cd}(A/\mathfrak p, I(A/\mathfrak p)) =
\text{cd}(A/\mathfrak p, \mathfrak m/\mathfrak p)$ である。
補題 \ref{local-cohomology-lemma-cd-maximal} により、これは $\dim(A/\mathfrak p)$ に等しい。
\end{proof}

\begin{lemma}
\label{local-cohomology-lemma-cd-blowup}
$A$ を Noether 環、$I \subset A$ をイデアルとする。
$b : X' \to X = \Spec(A)$ を $I$ の爆破とする。$b$ のファイバーの次元が
$\leq d - 1$ ならば、$\text{cd}(A, I) \leq d$ である。
\end{lemma}

\begin{proof}
$U = X \setminus V(I)$ とおく。$j : U \to X'$ を標準的な開埋め込みと書く。
Divisors, 第 \ref{divisors-section-blowing-up} 節を参照せよ。例外因子は
有効 Cartier 因子なので（Divisors, 補題
\ref{divisors-lemma-blowing-up-gives-effective-Cartier-divisor}）、$j$ はアフィンである。
Divisors, 補題 \ref{divisors-lemma-complement-locally-principal-closed-subscheme}
を参照せよ。$\mathcal{F}$ を準連接 $\mathcal{O}_U$-加群とする。このとき
$R^pj_*\mathcal{F} = 0$ が $p > 0$ に対して成り立つ。Cohomology of Schemes,
補題 \ref{coherent-lemma-relative-affine-vanishing} を参照せよ。一方、
$R^qb_*(j_*\mathcal{F}) = 0$ が $q \geq d$ に対して、Limits, 補題
\ref{limits-lemma-higher-direct-images-zero-above-dimension-fibre} により成り立つ。
したがって Leray スペクトル系列（Cohomology, 補題
\ref{cohomology-lemma-relative-Leray}）により
$R^n(b \circ j)_*\mathcal{F} = 0$ が $n \geq d$ に対して成り立つ。
ゆえに $H^n(U, \mathcal{F}) = 0$ が $n \geq d$ に対して成り立つ
（Cohomology, 補題 \ref{cohomology-lemma-apply-Leray} による）。
定義により、これは $\text{cd}(A, I) \leq d$ を意味する。
\end{proof}







\section{より一般の台}
\label{local-cohomology-section-supports}

\noindent
$A$ を Noether 環、$M$ を $A$-加群とする。
$T \subset \Spec(A)$ を特殊化で安定な部分集合とする
（Topology, 定義 \ref{topology-definition-specialization}）。
次のように定義する：
$$
H^0_T(M) = \colim_{Z \subset T} H^0_Z(M)
$$
ここで余極限は、閉部分集合 $Z$（空間 $\Spec(A)$ の閉部分集合であって
$T$\footnote{$T$ は特殊化で安定なので、
$T = \bigcup_{Z \subset T} Z$ である。
Topology, 補題 \ref{topology-lemma-stable-specialization} を参照せよ。}
に含まれるもの）からなる有向半順序集合にわたって取る。
言い換えると、元 $m$（加群 $M$ の元）が $H^0_T(M) \subset M$ に属することと、
その台 $V(\text{Ann}_R(m))$（元 $m$ の台）が
$T$ に含まれることは同値である。

\begin{lemma}
\label{local-cohomology-lemma-support}
$A$ を Noether 環、$T \subset \Spec(A)$ を特殊化で安定な部分集合とする。
$A$-加群 $M$ について、次の条件は同値である。
\begin{enumerate}
\item $H^0_T(M) = M$ である。
\item $\text{Supp}(M) \subset T$ である。
\end{enumerate}
このような $A$-加群の圏は、$A$-加群の圏の Serre 部分圏であり、
直和を取る操作で閉じている。
\end{lemma}

\begin{proof}
この同値性は、$M$ の元の台が $M$ の台に含まれ、
逆に $M$ の台がその各元の台の和集合に等しいことによる。
これらの加群の圏は Serre 部分圏
（Homology, 定義 \ref{homology-definition-serre-subcategory}）を
$\text{Mod}_A$ の中でなす。これは
Algebra, 補題 \ref{algebra-lemma-support-quotient} による。
直和についての主張の証明は省略する。
\end{proof}

\noindent
$A$ を Noether 環、$T \subset \Spec(A)$ を特殊化で安定な部分集合とする。
$\text{Mod}_{A, T} \subset \text{Mod}_A$ により
補題 \ref{local-cohomology-lemma-support} の Serre 部分圏を表す。
$D_T(A) \subset D(A)$ は、$D(A)$ の狭義充満かつ飽和した三角部分圏
（Derived Categories, 補題 \ref{derived-lemma-cohomology-in-serre-subcategory}）
であって、$A$-加群の複体のうち、そのコホモロジー加群が
$\text{Mod}_{A, T}$ に属するものからなる圏を表す。このとき関手
$$
D(\text{Mod}_{A, T}) \to D_T(A) \to D(A)
$$
を得る。Derived Categories, 第 \ref{derived-section-triangulated-sub} 節
の議論を参照せよ。
$RH^0_T : D(A) \to D(\text{Mod}_{A, T})$ により
$H^0_T$ の右導来拡張を表す。また
$$
R\Gamma_T : D^+(A) \to D^+_T(A),
$$
により、$RH^0_T : D^+(A) \to D^+(\text{Mod}_{A, T})$ と
$D^+(\text{Mod}_{A, T}) \to D^+_T(A)$ の合成を表す。$A$ の次元が
有限である場合\footnote{$\dim(A) = \infty$ の場合、この構成は
非有界複体に対して予期しない性質をもつことがある。}には、
$$
R\Gamma_T : D(A) \to D_T(A)
$$
により、$RH^0_T$ と
$D(\text{Mod}_{A, T}) \to D_T(A)$ の合成を表す。

\begin{lemma}
\label{local-cohomology-lemma-adjoint}
$A$ を Noether 環、$T \subset \Spec(A)$ を特殊化で安定な部分集合とする。
関手 $RH^0_T$ は、関手
$D(\text{Mod}_{A, T}) \to D(A)$ の右随伴である。
\end{lemma}

\begin{proof}
関手 $H^0_T(-)$ が包含関手
$\text{Mod}_{A, T} \to \text{Mod}_A$ の右随伴であることから従う。
Derived Categories, 補題 \ref{derived-lemma-derived-adjoint-functors} を参照せよ。
\end{proof}

\begin{lemma}
\label{local-cohomology-lemma-adjoint-ext}
$A$ を Noether 環、$T \subset \Spec(A)$ を特殊化で安定な部分集合とする。
任意の対象 $K$（圏 $D(A)$ の対象）に対して、次が成り立つ。
$$
H^i(RH^0_T(K)) = \colim_{Z \subset T\text{ 閉集合}} H^i_Z(K)
$$
\end{lemma}

\begin{proof}
$J^\bullet$ を $K$ を表す K-入射的複体とする。
定義により $RH^0_T$ は複体
$$
H^0_T(J^\bullet) = \colim H^0_Z(J^\bullet)
$$
で表される。等式は $H^0_T$ の定義から従う。
フィルター余極限は完全なので、この複体の次数 $i$ のコホモロジーは
$\colim H^i(H^0_Z(J^\bullet)) = \colim H^i_Z(K)$
であり、所望の結論を得る。
\end{proof}

\begin{lemma}
\label{local-cohomology-lemma-equal-plus}
$A$ を Noether 環、$T \subset \Spec(A)$ を特殊化で安定な部分集合とする。
関手 $D^+(\text{Mod}_{A, T}) \to D^+_T(A)$ は同値である。
\end{lemma}

\begin{proof}
$M$ を $\text{Mod}_{A, T}$ の対象とする。埋め込み
$M \to J$ を取り、その行き先を入射的 $A$-加群とする。
Dualizing Complexes, 命題
\ref{dualizing-proposition-structure-injectives-noetherian}
により、加群 $J$ は剰余体の入射包絡の直和である。
$E$ を素イデアル $\mathfrak p$ の剰余体の入射包絡とする。
$E$ は $\mathfrak p$-冪捩れなので、
$H^0_T(E) = 0$ が $\mathfrak p \not \in T$ の場合に成り立ち、
$H^0_T(E) = E$ が $\mathfrak p \in T$ の場合に成り立つ。
したがって $H^0_T(J)$ は入射包絡の直和として入射的であり
（上の命題による）、埋め込み $M \to H^0_T(J)$ を得る。
よって任意の対象 $M$（圏 $\text{Mod}_{A, T}$ の対象）は入射分解
$M \to J^\bullet$ であって、各 $J^n$ も $\text{Mod}_{A, T}$ に属するものをもつ。
したがって $RH^0_T(M) = M$ である。

\medskip\noindent
次に $K \in D_T^+(A)$ と仮定する。このときスペクトル系列
$$
R^qH^0_T(H^p(K)) \Rightarrow R^{p + q}H^0_T(K)
$$
(Derived Categories, 補題 \ref{derived-lemma-two-ss-complex-functor})
は収束し、上で見たように $q = 0$ の項だけが零でない。
したがって $RH^0_T(K) \to K$ は同型である。
よって関手 $D^+(\text{Mod}_{A, T}) \to D^+_T(A)$ は同値であり、
その擬逆は $RH^0_T$ で与えられる。
\end{proof}

\begin{lemma}
\label{local-cohomology-lemma-equal-full}
$A$ を Noether 環、$T \subset \Spec(A)$ を特殊化で安定な部分集合とする。
$\dim(A) < \infty$ ならば、関手
$D(\text{Mod}_{A, T}) \to D_T(A)$ は同値である。
\end{lemma}

\begin{proof}
$\dim(A) = d$ とおく。このとき $H^i_Z(M) = 0$ が $i > d$ に対して、
任意の閉部分集合 $Z$（空間 $\Spec(A)$ の閉部分集合）について成り立つ。
補題 \ref{local-cohomology-lemma-cd-dimension} を参照せよ。
補題 \ref{local-cohomology-lemma-adjoint-ext} により、$H^0_T$ のコホモロジー次元は有界である。

\medskip\noindent
$K \in D_T(A)$ とする。$RH^0_T(K) \to K$ が同型であることを示す。
$K$ が下に有界の場合は既に分かっている。
補題 \ref{local-cohomology-lemma-equal-plus} を参照せよ。一方、$H^0_T$ のコホモロジー次元は
有界なので、次数 $i$ のコホモロジー（複体
$RH_T^0(K)$ のもの）は $\tau_{\geq -d + i}K$ のみに依存し、結論を得る。
したがって $D(\text{Mod}_{A, T}) \to D_T(A)$ は同値であり、
その擬逆は $RH^0_T$ である。
\end{proof}

\begin{remark}
\label{local-cohomology-remark-upshot}
$A$ を Noether 環、$T \subset \Spec(A)$ を特殊化で安定な部分集合とする。
以上の議論の結論は、
$R\Gamma_T : D^+(A) \to D_T^+(A)$ が包含関手
$D_T^+(A) \to D^+(A)$ の右随伴である、ということである。
$\dim(A) < \infty$ ならば、
$R\Gamma_T : D(A) \to D_T(A)$ は包含関手
$D_T(A) \to D(A)$ の右随伴である。
いずれの場合にも次が成り立つ：
$$
H^i_T(K) = H^i(R\Gamma_T(K)) = R^iH^0_T(K) =
\colim_{Z \subset T\text{ 閉集合}} H^i_Z(K)
$$
これは 補題 \ref{local-cohomology-lemma-adjoint}, \ref{local-cohomology-lemma-adjoint-ext},
\ref{local-cohomology-lemma-equal-plus} および \ref{local-cohomology-lemma-equal-full} を組み合わせれば従う。
\end{remark}

\begin{lemma}
\label{local-cohomology-lemma-torsion-change-rings}
$A \to B$ を Noether 環の平坦準同型とする。
$T \subset \Spec(A)$ を特殊化で安定な部分集合とし、
$T' \subset \Spec(B)$ を $T$ の逆像とする。
このとき標準的な写像
$$
R\Gamma_T(K) \otimes_A^\mathbf{L} B
\longrightarrow
R\Gamma_{T'}(K \otimes_A^\mathbf{L} B)
$$
は $K \in D^+(A)$ に対して同型である。$A$ と $B$ の次元が有限ならば、
これは $K \in D(A)$ に対しても成り立つ。
\end{lemma}

\begin{proof}
写像 $R\Gamma_T(K) \to K$ から写像
$R\Gamma_T(K) \otimes_A^\mathbf{L} B \to K \otimes_A^\mathbf{L} B$.
を得る。$R\Gamma_T(K) \otimes_A^\mathbf{L} B$ のコホモロジー加群は
$T'$ に台をもつので、補題の射を得る。
この射は、$T$ が $\Spec(A)$ の閉部分集合である場合には
Dualizing Complexes, 補題 \ref{dualizing-lemma-torsion-change-rings} により同型である。
$H^i_T(K)$ は $H^i_Z(K)$ の余極限であることを思い出そう。
ここで $Z$ は $T$ の閉部分集合からなる有向集合を動く。
補題 \ref{local-cohomology-lemma-adjoint-ext} を参照せよ。これに対応して
$H^i_{T'}(K \otimes_A^\mathbf{L} B) =
\colim H^i_{Z'}(K \otimes_A^\mathbf{L} B)$ である。ここで $Z'$ は
$Z$ の逆像である。したがって、$\otimes_A B$ がフィルター余極限と可換であり、
高次 Tor が消滅することから結論が従う。
\end{proof}

\begin{lemma}
\label{local-cohomology-lemma-local-cohomology-ss}
$A$ を環とし、$T, T' \subset \Spec(A)$ を特殊化で安定な部分集合とする。
$K \in D^+(A)$ に対してスペクトル系列
$$
E_2^{p, q} = H^p_T(H^p_{T'}(K)) \Rightarrow H^{p + q}_{T \cap T'}(K)
$$
が存在する。これは Derived Categories, 補題
\ref{derived-lemma-grothendieck-spectral-sequence} にある形のものである。
\end{lemma}

\begin{proof}
$E$ を $D_{T \cap T'}(A)$ の対象とする。このとき
$$
\Hom(E, R\Gamma_T(R\Gamma_{T'}(K))) =
\Hom(E, R\Gamma_{T'}(K)) =
\Hom(E, K)
$$
である。最初の等式は $R\Gamma_T$ の随伴性により、
二つ目は $R\Gamma_{T'}$ の随伴性による。
一方、$J^\bullet$ が下に有界な入射的対象の複体であって
$K$ を表すならば、$H^0_{T'}(J^\bullet)$ は
入射的 $A$-加群の複体であり、$R\Gamma_{T'}(K)$ を表す。
したがって $H^0_T(H^0_{T'}(J^\bullet))$ は
$R\Gamma_T(R\Gamma_{T'}(K))$ を表す複体である。
よって $R\Gamma_T(R\Gamma_{T'}(K))$ は
$D^+_{T \cap T'}(A)$ の対象である。この二つの事実を組み合わせると、
$R\Gamma_{T \cap T'} = R\Gamma_T \circ R\Gamma_{T'}$ を得る。
主張で引用した補題により、これからスペクトル系列が得られる。
\end{proof}

\begin{lemma}
\label{local-cohomology-lemma-torsion-tensor-product}
$A$ を Noether 環、$T \subset \Spec(A)$ を特殊化で安定な部分集合とする。
$A$ の次元は有限であると仮定する。このとき
$$
R\Gamma_T(K) = R\Gamma_T(A) \otimes_A^\mathbf{L} K
$$
が $K \in D(A)$ に対して成り立つ。$K, L \in D(A)$ に対しては
$$
R\Gamma_T(K \otimes_A^\mathbf{L} L) =
K \otimes_A^\mathbf{L} R\Gamma_T(L) =
R\Gamma_T(K) \otimes_A^\mathbf{L} L =
R\Gamma_T(K) \otimes_A^\mathbf{L} R\Gamma_T(L)
$$
が成り立つ。$K$ または $L$ が $D_T(A)$ に属すれば、$K \otimes_A^\mathbf{L} L$ もそうである。
\end{lemma}

\begin{proof}
構成により $R\Gamma_T(A)$ は、複体 $J^\bullet$ であって
$\text{Mod}_{A, T}$ に属するものにより表せる。したがって $K$ を
K-平坦複体 $K^\bullet$ で表すと、$R\Gamma_T(A) \otimes_A^\mathbf{L} K$ は
複体 $\text{Tot}(J^\bullet \otimes_A K^\bullet)$（圏
$\text{Mod}_{A, T}$ 内の複体）で表される。
写像 $R\Gamma_T(A) \to A$ を用いて、写像
$R\Gamma_T(A) \otimes_A^\mathbf{L} K\to K$ を得る。
したがって $R\Gamma_T$ の随伴性により、標準的な写像
$$
R\Gamma_T(A) \otimes_A^\mathbf{L} K \longrightarrow R\Gamma_T(K)
$$
を得て、先ほど構成した写像はこれを経由する。$R\Gamma_T$ が
$D(A)$ 内の直和と可換であることに注意する。これは例えば
補題 \ref{local-cohomology-lemma-adjoint-ext}、有向余極限が直和と可換であること、
および通常の局所コホモロジーが直和と可換であること
（例えば Dualizing Complexes, 補題
\ref{dualizing-lemma-local-cohomology-adjoint}）による。
したがって More on Algebra, 注意 \ref{more-algebra-remark-P-resolution}
により、この写像が $K = A[k]$（ただし $k \in \mathbf{Z}$）
に対して同型であることを確かめれば十分であり、これは明らかである。

\medskip\noindent
最後の主張は、導来テンソル積の推移性により、今示した結果から従う。
\end{proof}





\section{局所コホモロジーのフィルトレーション}
\label{local-cohomology-section-filter-local-cohomology}

\noindent
補題 \ref{local-cohomology-lemma-local-cohomology-ss} のスペクトル系列に関連する
いくつかの手法を述べる。

\begin{lemma}
\label{local-cohomology-lemma-filter-local-cohomology}
$A$ を Noether 環、$T \subset \Spec(A)$ を特殊化で安定な部分集合とする。
$T' \subset T$ を $T$ の極小でない素イデアルの集合とする。
このとき $T'$ は $\Spec(A)$ の特殊化で安定な部分集合であり、
任意の $A$-加群 $M$ に対して完全列
$$
0 \to
\colim_{Z, f} H^1_f(H^{i - 1}_Z(M)) \to
H^i_{T'}(M) \to H^i_T(M) \to
\bigoplus\nolimits_{\mathfrak p \in T \setminus T'}
H^i_{\mathfrak p A_\mathfrak p}(M_\mathfrak p)
$$
が存在する。ここで余極限は、閉部分集合 $Z \subset T$
および元 $f \in A$ であって $V(f) \cap Z \subset T'$ を満たすものにわたって取る。
\end{lemma}

\begin{proof}
各 $Z$ と $f$ に対して、
Dualizing Complexes, 補題 \ref{dualizing-lemma-local-cohomology-ss}
のスペクトル系列は退化し、短完全列
$$
0 \to H^1_f(H^{i - 1}_Z(M)) \to
H^i_{Z \cap V(f)}(M) \to H^0_f(H^i_Z(M)) \to 0
$$
を与える。以下ではこの事実を断りなく用いる。

\medskip\noindent
$\xi \in H^i_T(M)$ の直和への像が零であるとする。
まず $\xi$ を、ある $\xi' \in H^i_Z(M)$ の像として表す。
ここで $Z \subset T$ はある閉部分集合である。補題 \ref{local-cohomology-lemma-adjoint-ext} を参照せよ。
このとき $\xi'$ は $H^i_{\mathfrak p A_\mathfrak p}(M_\mathfrak p)$ において零に写る。
これは任意の $\mathfrak p \in Z$、$\mathfrak p \not \in T'$ に対して成り立つ。
このような素イデアルは有限個しかないので、
$f \in A$ をいずれの素イデアルにも属さないように選び、
しかも $f$ が $\xi'$ を零化するようにできる。このとき $\xi'$ は
ある $\xi'' \in H^i_{Z'}(M)$ の像である。
ここで $Z' = Z \cap V(f)$ とおく。$f$ の選び方により
$Z' \subset T'$ であり、最後から二番目の項での完全性を得る。

\medskip\noindent
$\xi \in H^i_{T'}(M)$ が $H^i_T(M)$ において零に写るとする。
閉部分集合 $Z' \subset Z$ を、$Z' \subset T'$
かつ $Z \subset T$ であり、$\xi$ が $\xi' \in H^i_{Z'}(M)$ から来て、
その元が $H^i_Z(M)$ において零に写るように選ぶ。このとき $f \in A$ で
$V(f) \cap Z = Z'$ を満たすものが存在し、結論を得る。
\end{proof}

\begin{lemma}
\label{local-cohomology-lemma-zero}
$A$ を有限次元の Noether 環とする。
$T \subset \Spec(A)$ を特殊化で安定な部分集合とする。
$\{M_n\}_{n \geq 0}$ を $A$-加群の逆系とし、
$i \geq 0$ を整数とする。任意の $m$ に対して
整数 $m'(m) \geq m$ が存在し、すべての
$\mathfrak p \in T$ について誘導される写像
$$
H^i_{\mathfrak p A_\mathfrak p}(M_{k, \mathfrak p})
\longrightarrow
H^i_{\mathfrak p A_\mathfrak p}(M_{m, \mathfrak p})
$$
が $k \geq m'(m)$ のとき零であると仮定する。$m'' : \mathbf{N} \to \mathbf{N}$
を、$2^{\dim(T)}$ 回の自己合成で $m'$ から得られる写像とする。このとき写像
$H^i_T(M_k) \to H^i_T(M_m)$ は、すべての $k \geq m''(m)$ に対して零である。
\end{lemma}

\begin{proof}
まず一般的な注意を述べる。完全列
$$
(A_n) \to (B_n) \to (C_n)
$$
がアーベル群の逆系の間に与えられているとする。任意の
$m$ に対して整数 $m'(m) \geq m$ が存在し、
$$
A_k \to A_m
\quad\text{および}\quad
C_k \to C_m
$$
が $k \geq m'(m)$ のとき零であると仮定する。このとき $k \geq m'(m'(m))$
ならば、写像 $B_k \to B_m$ は零である。

\medskip\noindent
補題を $\dim(T)$ に関する帰納法で証明する。この次元は
$\dim(A)$ が有限なので有限である。$T' \subset T$ を
$T$ の極小でない素イデアルの集合とする。このとき $T'$ は
$\Spec(A)$ の特殊化で安定な部分集合であり、
補題の仮定は $T'$ に対しても成り立つ。
$\dim(T') < \dim(T)$ なので、$T'$ に対して補題が成り立つ。
任意の $A$-加群 $M$ に対し、完全列
$$
H^i_{T'}(M) \to H^i_T(M) \to
\bigoplus\nolimits_{\mathfrak p \in T \setminus T'}
H^i_{\mathfrak p A_\mathfrak p}(M_\mathfrak p)
$$
が補題 \ref{local-cohomology-lemma-filter-local-cohomology} により存在する。
したがって証明冒頭の注意により結論を得る。
\end{proof}

\begin{lemma}
\label{local-cohomology-lemma-essential-image}
$A$ を Noether 環、$T \subset \Spec(A)$ を特殊化で安定な部分集合とする。
$\{M_n\}_{n \geq 0}$ を $A$-加群の逆系とし、$i \geq 0$ を整数とする。
$A$ の次元は有限であり、任意の $m$ に対して整数 $m'(m) \geq m$ が存在し、
すべての $\mathfrak p \in T$ について次が成り立つと仮定する。
\begin{enumerate}
\item $H^{i - 1}_{\mathfrak p A_\mathfrak p}(M_{k, \mathfrak p})
\to H^{i - 1}_{\mathfrak p A_\mathfrak p}(M_{m, \mathfrak p})$
は $k \geq m'(m)$ のとき零である。
\item $ H^i_{\mathfrak p A_\mathfrak p}(M_{k, \mathfrak p}) \to
H^i_{\mathfrak p A_\mathfrak p}(M_{m, \mathfrak p})$
の像 $G(\mathfrak p, m)$ は $k \geq m'(m)$ に依存せず、さらに
$G(\mathfrak p, m)$ から
$H^i_{\mathfrak p A_\mathfrak p}(M_{0, \mathfrak p})$ への写像は単射である。
\end{enumerate}
このとき整数 $m_0$ が存在し、任意の $m \geq m_0$ に対して
整数 $m''(m) \geq m$ が存在して、
$k \geq m''(m)$ ならば $H^i_T(M_k) \to H^i_T(M_m)$ の像から
$H^i_T(M_{m_0})$ への写像は単射である。
\end{lemma}

\begin{proof}
まず一般的な注意を述べる。完全列
$$
(A_n) \to (B_n) \to (C_n) \to (D_n)
$$
がアーベル群の逆系の間に与えられているとする。
整数 $m_0$ が存在し、任意の $m \geq m_0$ に対して
整数 $m'(m) \geq m$ が存在して、写像
$$
\Im(B_k \to B_m) \longrightarrow B_{m_0}
\quad\text{および}\quad
\Im(D_k \to D_m) \longrightarrow D_{m_0}
$$
が $k \geq m'(m)$ のとき単射であり、$A_k \to A_m$ が
$k \geq m'(m)$ のとき零であると仮定する。このとき $m \geq m'(m_0)$ かつ $k \geq m'(m'(m))$
ならば、写像
$$
\Im(C_k \to C_m) \to C_{m'(m_0)}
$$
は単射である。実際、$c_0 \in C_m$ を $c_3 \in C_k$ の像とし、
$c_0$ が $C_{m'(m_0)}$ において零に写るとする。次の図式を考える：
$$
C_k \to C_{m'(m'(m))} \to C_{m'(m)} \to C_m \to C_{m'(m_0)},\quad
c_3 \mapsto c_2 \mapsto c_1 \mapsto c_0 \mapsto 0
$$
$c_0 = 0$ を示せばよい。
像 $d_3$（元 $c_3$ の像）は $C_{m_0}$ において零に写るので、
像 $d_1 \in D_{m'(m)}$ は零である。
したがって $b_1 \in B_{m'(m)}$ を、その像が
$c_1$ となるように選べる。$c_3$ は
$C_{m'(m_0)}$ において零に写るので、元 $a_{-1} \in A_{m'(m_0)}$ で、
像 $b_{-1} \in B_{m'(m_0)}$（元 $b_1$ の像）に写るものが存在する。
$a_{-1}$ は $A_{m_0}$ において零に写るので、
$b_1$ は $B_{m_0}$ において零に写る。したがって像 $b_0 \in B_m$
は零であり、もちろんこれから所望の $c_0 = 0$ が従う。

\medskip\noindent
補題を $\dim(T)$ に関する帰納法で証明する。この次元は
$\dim(A)$ が有限なので有限である。$T' \subset T$ を
$T$ の極小でない素イデアルの集合とする。このとき $T'$ は
$\Spec(A)$ の特殊化で安定な部分集合であり、補題の仮定は
$T'$ に対しても成り立つ。$\dim(T') < \dim(T)$ なので、
$T'$ に対して補題が成り立つ。任意の $A$-加群 $M$ に対して完全列
$$
0 \to \colim_{Z, f} H^1_f(H^{i - 1}_Z(M)) \to
H^i_{T'}(M) \to H^i_T(M) \to
\bigoplus\nolimits_{\mathfrak p \in T \setminus T'}
H^i_{\mathfrak p A_\mathfrak p}(M_\mathfrak p)
$$
が補題 \ref{local-cohomology-lemma-filter-local-cohomology} により存在する。
したがって証明冒頭の注意と、群の系
$$
\left\{\colim_{Z, f} H^1_f(H^{i - 1}_Z(M_n))\right\}_{n \geq 0}
$$
が補題 \ref{local-cohomology-lemma-zero} により pro-零であることから結論を得る。
ここでは、その補題の関数 $m''(m)$（加群 $H^{i - 1}_Z(M)$ に対するもの）が
$Z$ に依存しないことを用いている。
\end{proof}










\section{局所コホモロジーの有限性 I}
\label{local-cohomology-section-finiteness}

\noindent
局所コホモロジー加群の有限性に関する Faltings の方法に従う。
\cite{Faltings-annulators} および \cite{Faltings-finiteness} を参照せよ。
次の補題は、局所コホモロジー加群が有限加群であることを示すには、
適切な零化イデアルをもつことを示せば十分であることを述べる。

\begin{lemma}
\label{local-cohomology-lemma-check-finiteness-local-cohomology-by-annihilator}
\begin{reference}
\cite[補題 3]{Faltings-annulators}
\end{reference}
$A$ を Noether 環、$T \subset \Spec(A)$ を特殊化で安定な部分集合とする。
$M$ を有限 $A$-加群とし、$n \geq 0$ とする。
次の条件は同値である。
\begin{enumerate}
\item $H^i_T(M)$ は $i \leq n$ に対して有限である。
\item イデアル $J \subset A$ で $V(J) \subset T$ を満たすものが存在し、
$J$ は $H^i_T(M)$ を $i \leq n$ に対して零化する。
\end{enumerate}
もし $T = V(I) = Z$ があるイデアル $I \subset A$ に対して成り立つならば、
これらはさらに次の条件とも同値である。
\begin{enumerate}
\item[(3)] ある $e \geq 0$ が存在し、$I^e$ は
$H^i_Z(M)$ を $i \leq n$ に対して零化する。
\end{enumerate}
\end{lemma}

\begin{proof}
(1) と (2) の同値性を $n$ に関する帰納法で証明する。
$n = 0$ のとき $H^0_T(M) \subset M$ は有限なので、(1) が成り立つ。
$H^0_T(M) = \colim H^0_{V(J)}(M)$ であり、ここで $J$ は (2) のようなものを動くので、
(2) も成り立つ。以下、$n > 0$ と仮定する。

\medskip\noindent
(1) が成り立つと仮定する。$H^i_J(M) = H^i_{V(J)}(M)$ であることを思い出そう。
Dualizing Complexes, 補題 \ref{dualizing-lemma-local-cohomology-noetherian} を参照せよ。
したがって $H^i_T(M) = \colim H^i_J(M)$ である。ここで余極限はイデアル
$J \subset A$ で $V(J) \subset T$ を満たすものにわたって取る。
補題 \ref{local-cohomology-lemma-adjoint-ext} を参照せよ。$H^i_T(M)$ は有限生成であり、
これは $i \leq n$ に対して成り立つので、(2) のような $J \subset A$ を選び、
$H^i_J(M) \to H^i_T(M)$ が $i \leq n$ に対して全射となるようにできる。
したがって有限個の生成元はいずれも $J$-冪捩れ元であり、
$J$ をその適当な冪に置き換えると (2) が成り立つ。

\medskip\noindent
(2) のような $J$ が存在すると仮定する。$N = H^0_T(M)$ および $M' = M/N$ とおく。
$R\Gamma_T$ の構成により、
$H^i_T(N) = 0$ が $i > 0$ に対して成り立ち、$H^0_T(N) = N$ である。
注意 \ref{local-cohomology-remark-upshot} を参照せよ。したがって
$H^0_T(M') = 0$ であり、$H^i_T(M') = H^i_T(M)$ が $i > 0$ に対して成り立つ。
ゆえに $M$ を $M'$ に置き換えてよい。
したがって $H^0_T(M) = 0$ と仮定してよい。
これは、$M$ の随伴素イデアルからなる有限集合が
$T$ と交わらないことを意味する。素イデアル回避補題（Algebra, 補題 \ref{algebra-lemma-silly}）
により、$f \in J$ を $M$ のいずれの随伴素イデアルにも属さないように選べる。
このとき短完全列
$$
0 \to M \to M \to M/fM \to 0
$$
に付随する局所コホモロジーの長完全列は、短完全列
$$
0 \to H^i_T(M) \to H^i_T(M/fM) \to H^{i + 1}_T(M) \to 0
$$
に分かれる。ただし $i < n$ である。したがって $J^2$ は $H^i_T(M/fM)$ を
$i < n$ に対して零化する。帰納法の仮定により、$H^i_T(M/fM)$ は
$i < n$ に対して有限である。短完全列をもう一度用いると、
$H^{i + 1}_T(M)$ が $i < n$ に対して有限であることを得る。これが所望の結論である。

\medskip\noindent
$T = V(I)$ の場合の (2) と (3) の同値性の証明は省略する。
\end{proof}

\noindent
次の Faltings の結果により、局所コホモロジーの有限性を
局所環の段階で証明することができる。

\begin{lemma}
\label{local-cohomology-lemma-check-finiteness-local-cohomology-locally}
\begin{reference}
これは \cite[Satz 1]{Faltings-finiteness} の特殊な場合である。
\end{reference}
$A$ を Noether 環、$I \subset A$ をイデアル、$M$ を有限
$A$-加群、$n \geq 0$ を整数とする。$Z = V(I)$ とおく。
次の条件は同値である。
\begin{enumerate}
\item 加群 $H^i_Z(M)$ は $i \leq n$ に対して有限である。
\item すべての $\mathfrak p \in \Spec(A)$ に対して、加群
$H^i_Z(M)_\mathfrak p$（$i \leq n$）は有限 $A_\mathfrak p$-加群である。
\end{enumerate}
\end{lemma}

\begin{proof}
(1) $\Rightarrow$ (2) は直ちに従う。逆を
$n$ に関する帰納法で証明する。$n = 0$ の場合には (1) と (2) が
常に成り立つので明らかである。

\medskip\noindent
$n > 0$ とし、(2) が成り立つと仮定する。$N = H^0_Z(M)$ および $M' = M/N$ とおく。
Dualizing Complexes, 補題 \ref{dualizing-lemma-divide-by-torsion} により
$M$ を $M'$ に置き換えてよい。
したがって $H^0_Z(M) = 0$ と仮定してよい。
これは $\text{depth}_I(M) > 0$ を意味する
（Dualizing Complexes, 補題 \ref{dualizing-lemma-depth}）。
$f \in I$ を $M$ 上の非零因子として選び、短完全列
$$
0 \to M \to M \to M/fM \to 0
$$
を考える。これから長完全列
$$
0 \to H^0_Z(M/fM) \to H^1_Z(M) \to H^1_Z(M) \to H^1_Z(M/fM) \to
H^2_Z(M) \to \ldots
$$
を得て、局所化後にも同様の列を得る。したがって仮定 (2) により、
加群 $H^i_Z(M/fM)_\mathfrak p$ は $i < n$ に対して有限である。よって
帰納法の仮定により $H^i_Z(M/fM)$ は $i < n$ に対して有限である。

\medskip\noindent
$\mathfrak p$ を $A$ の素イデアルであって、
$H^i_Z(M)$ の随伴素イデアルとなるものとする。ここで $i \leq n$ とする。
$\mathfrak p$ は元 $x \in H^i_Z(M)$ の零化イデアルであるとする。
このとき $\mathfrak p \in Z$ なので、
$f \in \mathfrak p$ である。したがって $fx = 0$ であり、$x$ は
$H^{i - 1}_Z(M/fM)$ のある元から、上の長完全列の境界写像 $\delta$ により得られる。
ゆえに $\mathfrak p$ は有限加群 $\Im(\delta)$ の随伴素イデアルである。
したがって $\text{Ass}(H^i_Z(M))$ は $i \leq n$ に対して有限である。
Algebra, 補題 \ref{algebra-lemma-finite-ass} を参照せよ。

\medskip\noindent
次を思い出そう：
$$
H^i_Z(M) \subset
\prod\nolimits_{\mathfrak p \in \text{Ass}(H^i_Z(M))}
H^i_Z(M)_\mathfrak p
$$
これは Algebra, 補題 \ref{algebra-lemma-zero-at-ass-zero} による。
仮定により右辺の加群は有限かつ $I$-冪捩れなので、
整数 $e_{\mathfrak p, i} \geq 0$（$i \leq n$、
$\mathfrak p \in \text{Ass}(H^i_Z(M))$）を選び、
$I^{e_{\mathfrak p, i}}$ が $H^i_Z(M)_\mathfrak p$ を零化するようにできる。
したがって $I^e$（ただし $e = \max\{e_{\mathfrak p, i}\}$）は $H^i_Z(M)$ を
$i \leq n$ に対して零化する。
補題 \ref{local-cohomology-lemma-check-finiteness-local-cohomology-by-annihilator}
により、$H^i_Z(M)$ は $i \leq n$ に対して有限である。
\end{proof}

\begin{lemma}
\label{local-cohomology-lemma-annihilate-local-cohomology}
$A$ を環とし、$J \subset I \subset A$ を有限生成イデアルとする。
$i \geq 0$ を整数とし、$Z = V(I)$ とおく。
$H^i_Z(A)$ が $J^n$ によって、ある $n$ に対して零化されるならば、
$H^i_Z(M)$ は $J^m$ によって、ある $m = m(M)$ に対して零化される。
これは任意の有限表示 $A$-加群 $M$ であって、
$M_f$ が有限局所自由 $A_f$-加群であることがすべての $f \in I$ に対して成り立つものについて成り立つ。
\end{lemma}

\begin{proof}
零化イデアル $\mathfrak a$（加群 $H^i_Z(M)$ の零化イデアル）を考える。
$\mathfrak p \subset A$ を $\mathfrak p \not \in Z$ となる素イデアルとする。
仮定により $f \in I$、$f \not \in \mathfrak p$
および同型 $\varphi : A_f^{\oplus r} \to M_f$ が存在し、
これは $A_f$-加群の同型である。分母を払うと
（さらに $M$ が有限表示であることを用いると）、写像
$$
a : A^{\oplus r} \longrightarrow M
\quad\text{および}\quad
b : M \longrightarrow A^{\oplus r}
$$
であって、$a_f = f^N \varphi$ および $b_f = f^N \varphi^{-1}$ がある $N$ に対して成り立つものを得る。
さらに $a \circ b$ と $b \circ a$ は
$f^{2N}$ 倍写像に等しいと仮定してよい。したがって $H^i_Z(M)$ は
$f^{2N}J^n$ により零化される。すなわち $f^{2N}J^n \subset \mathfrak a$ である。

\medskip\noindent
$U = \Spec(A) \setminus Z$ は準コンパクトなので、有限個の元
$f_1, \ldots, f_t$ と整数 $N_1, \ldots, N_t$ を選び、$U = \bigcup D(f_j)$ かつ
$f_j^{2N_j}J^n \subset \mathfrak a$ とできる。このとき $V(I) = V(f_1, \ldots, f_t)$ であり、
$I$ は有限生成なので、
$I^M \subset (f_1, \ldots, f_t)$ がある $M$ に対して成り立つ。
以上から $J^m \subset \mathfrak a$ が
$m \gg 0$ に対して成り立つ。例えば $m = M (2N_1 + \ldots + 2N_t) n$ と取ればよい。
\end{proof}

\begin{lemma}
\label{local-cohomology-lemma-local-finiteness-for-finite-locally-free}
$A$ を Noether 環、$I \subset A$ をイデアルとし、$Z = V(I)$ とおく。
$n \geq 0$ を整数とする。$H^i_Z(A)$ が $0 \leq i \leq n$ に対して有限ならば、
同じことが $H^i_Z(M)$（$0 \leq i \leq n$）についても成り立つ。
ここで任意の有限 $A$-加群 $M$ であって、$M_f$ が有限局所自由
$A_f$-加群であることがすべての $f \in I$ に対して成り立つものを考える。
\end{lemma}

\begin{proof}
$H^i_Z(A)$ が $0 \leq i \leq n$ に対して有限であるという仮定から、
ある $e \geq 0$ が存在し、$I^e$ が
$H^i_Z(A)$ を $0 \leq i \leq n$ に対して零化することが従う。
補題 \ref{local-cohomology-lemma-check-finiteness-local-cohomology-by-annihilator} を参照せよ。
すると補題 \ref{local-cohomology-lemma-annihilate-local-cohomology} により、
$H^i_Z(M)$（$0 \leq i \leq n$）は
$I^m$ によって、ある $m = m(M, i)$ に対して零化される。同じ $m$ を
すべての $0 \leq i \leq n$ に対して取ってよい。したがって
補題 \ref{local-cohomology-lemma-check-finiteness-local-cohomology-by-annihilator} により、
$H^i_Z(M)$ は $0 \leq i \leq n$ に対して有限であり、
所望の結論を得る。
\end{proof}





\section{順像の有限性 I}
\label{local-cohomology-section-finiteness-pushforward}

\noindent
本節では有限性定理の最も易しい非自明な場合を論じる。
すなわち、第一局所コホモロジーの有限性、あるいは同値なこととして
$j_*\mathcal{F}$ の有限性を扱う。ここで $j : U \to X$ は開浸入、
$X$ は局所 Noether、$\mathcal{F}$ は $U$ 上の連接層である。
Koll\'ar の方法（\cite{Kollar-variants} および
\cite{Kollar-local-global-hulls}）に従い、
必要十分条件を得る。命題 \ref{local-cohomology-proposition-kollar} を参照せよ。
高次順像または高次局所コホモロジー群に関心のある読者は、
第 \ref{local-cohomology-section-finiteness-pushforward-II} 節または
第 \ref{local-cohomology-section-finiteness-II} 節へ進むとよい
（これらは本節の残りとは独立に展開される）。

\begin{lemma}
\label{local-cohomology-lemma-check-finiteness-pushforward-on-associated-points}
$X$ を局所 Noether 概形とし、$j : U \to X$ を、
補集合が $Z$ である開部分概形の包含とする。$x \in U$ に対して、
$i_x : W_x \to U$ を、生成点が $x$ である整閉部分概形とする。
$\mathcal{F}$ を連接 $\mathcal{O}_U$-加群とする。
次の条件は同値である。
\begin{enumerate}
\item すべての $x \in \text{Ass}(\mathcal{F})$ に対して、
$\mathcal{O}_X$-加群 $j_*i_{x, *}\mathcal{O}_{W_x}$ は連接である。
\item $j_*\mathcal{F}$ は連接である。
\end{enumerate}
\end{lemma}

\begin{proof}
まず (1) が (2) を含意することを示す。(1) を仮定する。
主張は $X$ 上局所的なので、$X$ はアフィンであると仮定してよい。
このとき $U$ は準コンパクトであり、したがって
$\text{Ass}(\mathcal{F})$ は有限である
（Divisors, 補題 \ref{divisors-lemma-finite-ass}）。
よって随伴点の個数に関する帰納法を用いてよい。$x \in U$ を
$\mathcal{F}$ の台の既約成分の生成点とする。
Divisors, 補題 \ref{divisors-lemma-finite-ass} により
$x \in \text{Ass}(\mathcal{F})$ である。$x$ の選び方により、
$\dim(\mathcal{F}_x) = 0$ である。ここで次元は
$\mathcal{O}_{X, x}$-加群としての次元である。
したがって $\mathcal{F}_x$ は $\mathcal{O}_{X, x}$-加群として有限長である
（Algebra, 補題 \ref{algebra-lemma-support-point}）。
そこでこの長さに関する帰納法を用いる。

\medskip\noindent
$\mathcal{G} = j_*i_{x, *}\mathcal{O}_{W_x}$ とおく。
これは仮定により連接 $\mathcal{O}_X$-加群である。
$\mathcal{G}_x = \kappa(x)$ が成り立つ。非零写像
$\varphi_x : \mathcal{F}_x \to \kappa(x) = \mathcal{G}_x$ を選ぶ。
Cohomology of Schemes, 補題 \ref{coherent-lemma-map-stalks-local-map}
により、開集合 $x \in V \subset U$ と写像
$\varphi_V : \mathcal{F}|_V \to \mathcal{G}|_V$ が存在し、
その $x$ における茎は $\varphi_x$ である。
$f \in \Gamma(X, \mathcal{O}_X)$ を $x$ で消えず、
$D(f) \subset V$ となるように選ぶ。例えば
Cohomology of Schemes, 補題 \ref{coherent-lemma-homs-over-open} により、
$\varphi_V$ は写像 $f^n\mathcal{F} \to \mathcal{G}|_U$ へ拡張される。
ここで $n$ はある整数である。
$f^n$ 倍写像を前合成すると、写像
$\mathcal{F} \to \mathcal{G}|_U$ で、その $x$ における茎が非零なものを得る。
$\mathcal{F}' \subset \mathcal{F}$ をその核とする。
$\text{Ass}(\mathcal{F}') \subset \text{Ass}(\mathcal{F})$ であることに注意する。
Divisors, 補題 \ref{divisors-lemma-ses-ass} を参照せよ。
$\text{length}_{\mathcal{O}_{X, x}}(\mathcal{F}'_x) =
\text{length}_{\mathcal{O}_{X, x}}(\mathcal{F}_x) - 1$
なので、帰納法の仮定を適用し、$j_*\mathcal{F}'$ が連接であることを得る。
$\mathcal{G} = j_*(\mathcal{G}|_U) = j_*i_{x, *}\mathcal{O}_{W_x}$
は連接なので、完全列
$$
0 \to j_*\mathcal{F}' \to j_*\mathcal{F} \to \mathcal{G}
$$
を考えることができる。Schemes, 補題
\ref{schemes-lemma-push-forward-quasi-coherent} により、
層 $j_*\mathcal{F}$ は準連接である。
したがって $j_*\mathcal{F}$ の $j_*(\mathcal{G}|_U)$ における像は、
Cohomology of Schemes, 補題
\ref{coherent-lemma-coherent-Noetherian-quasi-coherent-sub-quotient}
により連接である。最後に
Cohomology of Schemes, 補題 \ref{coherent-lemma-coherent-abelian-Noetherian}
により、$j_*\mathcal{F}$ は連接である。

\medskip\noindent
(2) が成り立つと仮定する。上とまったく同様にして、
$X$ がアフィンである場合に帰着する。$x \in \text{Ass}(\mathcal{F})$ を選び、
$\mathcal{G} = j_*i_{x, *}\mathcal{O}_{W_x}$ とおく。
次に非零写像
$\varphi_x : \mathcal{G}_x = \kappa(x) \to \mathcal{F}_x$ を選ぶ。
これはまさに $x$ が $\mathcal{F}$ の随伴点であるために存在する。
上とまったく同様に議論して、$\varphi_x$ は
$\mathcal{O}_U$-加群の写像
$\varphi : \mathcal{G}|_U \to \mathcal{F}$ に拡張されると仮定してよい。
このとき $\varphi$ は単射であり（例えば
Divisors, 補題 \ref{divisors-lemma-check-injective-on-ass} による）、
単射
$\mathcal{G} = j_*(\mathcal{G}|_V) \to j_*\mathcal{F}$ を得る。
したがって (1) が成り立つ。
\end{proof}

\begin{lemma}
\label{local-cohomology-lemma-finiteness-pushforwards-and-H1-local}
$A$ を Noether 環、$I \subset A$ をイデアルとする。
$X = \Spec(A)$、$Z = V(I)$、$U = X \setminus Z$ とおき、
$j : U \to X$ を包含射とする。
$\mathcal{F}$ を連接 $\mathcal{O}_U$-加群とする。このとき
\begin{enumerate}
\item 有限 $A$-加群 $M$ が存在し、$\mathcal{F}$ は
$\widetilde{M}$ の $U$ への制限である。
\item $M$ が与えられると完全列
$$
0 \to H^0_Z(M) \to M \to H^0(U, \mathcal{F}) \to H^1_Z(M) \to 0
$$
と、同型 $H^p(U, \mathcal{F}) = H^{p + 1}_Z(M)$ が
$p \geq 1$ に対して存在する。
\item $M$ および $p \geq 0$ が与えられると、次の条件は同値である。
\begin{enumerate}
\item $R^pj_*\mathcal{F}$ は連接である。
\item $H^p(U, \mathcal{F})$ は有限 $A$-加群である。
\item $H^{p + 1}_Z(M)$ は有限 $A$-加群である。
\end{enumerate}
\item (3) の同値な条件が $p = 0$ に対して成り立つならば、
$M = \Gamma(U, \mathcal{F})$ と取ることができ、この場合
$H^0_Z(M) = H^1_Z(M) = 0$ である。
\end{enumerate}
\end{lemma}

\begin{proof}
Properties, 補題 \ref{properties-lemma-lift-finite-presentation}
により、連接 $\mathcal{O}_X$-加群 $\mathcal{F}'$ であって、
その $U$ への制限が $\mathcal{F}$ と同型であるものが存在する。
$\mathcal{F}'$ は (1) の有限 $A$-加群 $M$ に対応するとする。
$R^pj_*\mathcal{F}$ は準連接であり
（Cohomology of Schemes, 補題
\ref{coherent-lemma-quasi-coherence-higher-direct-images}）、
$A$-加群 $H^p(U, \mathcal{F})$ に対応することに注意する。
補題 \ref{local-cohomology-lemma-local-cohomology-is-local-cohomology} と
Cohomology, 第 \ref{cohomology-section-cohomology-support} 節および
第 \ref{cohomology-section-cohomology-support-bis} 節の議論により、
完全列
$$
0 \to H^0_Z(M) \to M \to H^0(U, \mathcal{F}) \to H^1_Z(M) \to 0
$$
と、同型 $H^p(U, \mathcal{F}) = H^{p + 1}_Z(M)$ を
$p \geq 1$ に対して得る。
ここで $H^j(X, \mathcal{F}') = 0$ が $j > 0$ に対して成り立つことを用いた。
実際、$X$ はアフィンかつ $\mathcal{F}'$ は準連接である
（Cohomology of Schemes, 補題
\ref{coherent-lemma-quasi-coherent-affine-cohomology-zero}）。
これで (2) が示された。
(3) と (4) は (2) から直ちに従う。補題
\ref{local-cohomology-lemma-local-cohomology} も参照せよ。
\end{proof}

\begin{lemma}
\label{local-cohomology-lemma-finiteness-pushforward}
$X$ を局所 Noether 概形とする。
$j : U \to X$ を、補集合が $Z$ である開部分概形の包含とする。
$\mathcal{F}$ を連接 $\mathcal{O}_U$-加群とする。次を仮定する。
\begin{enumerate}
\item $X$ は Nagata である。
\item $X$ は普遍鎖状である。
\item $x \in \text{Ass}(\mathcal{F})$ および
$z \in Z \cap \overline{\{x\}}$ に対して
$\dim(\mathcal{O}_{\overline{\{x\}}, z}) \geq 2$ である。
\end{enumerate}
このとき $j_*\mathcal{F}$ は連接である。
\end{lemma}

\begin{proof}
補題 \ref{local-cohomology-lemma-check-finiteness-pushforward-on-associated-points} により、
$j_*i_{x, *}\mathcal{O}_{W_x}$ が連接であることを
$x \in \text{Ass}(\mathcal{F})$ に対して示せば十分である。
$\pi : Y \to X$ を、$X$ の $\Spec(\kappa(x))$ における正規化とする。
Morphisms, 第 \ref{morphisms-section-normalization} 節を参照せよ。
Morphisms, 補題 \ref{morphisms-lemma-nagata-normalization-finite-general}
により、射 $\pi$ は有限である。$\pi$ が有限なので、
$\mathcal{G} = \pi_*\mathcal{O}_Y$ は
Cohomology of Schemes, 補題 \ref{coherent-lemma-finite-pushforward-coherent}
により連接 $\mathcal{O}_X$-加群である。
$W_x = U \cap \pi(Y)$ であることに注意する。したがって
$\pi|_{\pi^{-1}(U)} : \pi^{-1}(U) \to U$ は $i_x : W_x \to U$ を経由し、
標準的な写像
$$
i_{x, *}\mathcal{O}_{W_x}
\longrightarrow
(\pi|_{\pi^{-1}(U)})_*(\mathcal{O}_{\pi^{-1}(U)}) =
(\pi_*\mathcal{O}_Y)|_U = \mathcal{G}|_U
$$
を得る。この写像は単射である（例えば Divisors, 補題
\ref{divisors-lemma-check-injective-on-ass} による）。したがって
$j_*i_{x, *}\mathcal{O}_{W_x} \subset j_*\mathcal{G}|_U$ であり、
$j_*\mathcal{G}|_U$ が連接であることを示せば十分である。

\medskip\noindent
$j_*(\mathcal{G}|_U)$ が連接であることを示せばよい。
Divisors, 補題 \ref{divisors-lemma-depth-2-hartog} を
$$
\mathcal{G} \longrightarrow j_*(\mathcal{G}|_U)
$$
に適用できると主張する。これで証明は終わる。
$\text{depth}(\mathcal{G}_z) \geq 2$ を $z \in Z$ に対して
示せば十分である。
$y_1, \ldots, y_n \in Y$ を $z$ に写る点とする。
Algebra, 補題 \ref{algebra-lemma-depth-goes-down-finite} により、
$\text{depth}(\mathcal{O}_{Y, y_i}) \geq 2$ が $i = 1, \ldots, n$
に対して成り立つことを示せば十分である。もしそうでなければ、
Properties, 補題 \ref{properties-lemma-criterion-normal} により、
$\dim(\mathcal{O}_{Y, y_i}) = 1$ となる $i$ が存在する。
これは $\pi : Y \to \overline{\{x\}}$ に対する次元公式
（Morphisms, 補題 \ref{morphisms-lemma-dimension-formula}）
と仮定 (3) に反する。
\end{proof}

\begin{lemma}
\label{local-cohomology-lemma-sharp-finiteness-pushforward}
$X$ を整局所 Noether 概形とする。$j : U \to X$ を、
補集合が $Z$ である空でない開部分概形の包含とする。
すべての $z \in Z$ と、$\mathfrak p$ を
$\mathcal{O}_{X, z}^\wedge$ の任意の随伴素イデアルとして、
$\dim(\mathcal{O}_{X, z}^\wedge/\mathfrak p) \geq 2$ であると仮定する。
このとき $j_*\mathcal{O}_U$ は連接である。
\end{lemma}

\begin{proof}
$X$ はアフィンであると仮定してよい。
補題 \ref{local-cohomology-lemma-check-finiteness-local-cohomology-locally} および
\ref{local-cohomology-lemma-finiteness-pushforwards-and-H1-local} を用いると、
$X = \Spec(A)$ で、$(A, \mathfrak m)$ は Noether 局所整域、
$\mathfrak m \in Z$ である場合に帰着する。
そこで $d = \dim(A)$ に関する帰納法を用いる。
（基底の場合は $d = 0, 1$ だが、局所環に関する仮定によりこれらは起こらない。）
$V = \Spec(A) \setminus \{\mathfrak m\}$ とおく。
$V$ の局所環の次元は $d$ より真に小さいことに注意する。
$j' : U \to V$ に対して同じ議論を繰り返し、
帰納法を用いると、$j'_*\mathcal{O}_U$ は連接
$\mathcal{O}_V$-加群であることが分かる。
$f \in A$ を非零元として選び、これは $Z$ 上で消えるものとする。
$D(f) \cap V \subset U$ なので、ある $n$ が存在して、
$f^n$ 倍写像が $U$ 上で写像
$f^n : j'_*\mathcal{O}_U \to \mathcal{O}_V$ へ拡張される。
これは $V$ 上の写像である
（例えば Cohomology of Schemes, 補題
\ref{coherent-lemma-homs-over-open} による）。
この写像は単射なので、単射
$$
j_*\mathcal{O}_U = j''_* j'_* \mathcal{O}_U \to j''_*\mathcal{O}_V
$$
を $X$ 上で得る。ここで $j'' : V \to X$ は包含射である。
したがって $j''_*\mathcal{O}_V$ が連接であることを示せば十分である。
言い換えると、$X$ は Noether 局所整域のスペクトルであり、
$Z$ はその閉点のみからなると仮定してよい。

\medskip\noindent
$X = \Spec(A)$、$(A, \mathfrak m)$ は局所環、
$Z = \{\mathfrak m\}$ と仮定する。
$A^\wedge$ を $A$ の完備化とし、
$X^\wedge = \Spec(A^\wedge)$、$Z^\wedge = \{\mathfrak m^\wedge\}$、
$U^\wedge = X^\wedge \setminus Z^\wedge$、
$\mathcal{F}^\wedge = \mathcal{O}_{U^\wedge}$ とおく。
環 $A^\wedge$ は普遍鎖状かつ Nagata である
（Algebra, 注意
\ref{algebra-remark-Noetherian-complete-local-ring-universally-catenary} および
補題 \ref{algebra-lemma-Noetherian-complete-local-Nagata}）。
さらに $X^\wedge, Z^\wedge, U^\wedge, \mathcal{F}^\wedge$ に対する
補題 \ref{local-cohomology-lemma-finiteness-pushforward} の条件 (3) は仮定により成り立つ。
したがって
$(U^\wedge \to X^\wedge)_*\mathcal{O}_{U^\wedge}$
は連接である。射 $c : X^\wedge \to X$ は平坦なので、
$j_*\mathcal{O}_U$ の引き戻しは
$(U^\wedge \to X^\wedge)_*\mathcal{O}_{U^\wedge}$ である
（Cohomology of Schemes, 補題
\ref{coherent-lemma-flat-base-change-cohomology}）。
最後に $c$ は忠実平坦なので、
Descent, 補題 \ref{descent-lemma-finite-type-descends} により
$j_*\mathcal{O}_U$ は連接である。
\end{proof}

\begin{remark}
\label{local-cohomology-remark-closure}
$j : U \to X$ を局所 Noether 概形の開浸入とする。
$x \in U$ とする。$i_x : W_x \to U$ を生成点 $x$ をもつ整閉部分概形とし、
$\overline{\{x\}}$ を $X$ における閉包とする。
このとき可換図式
$$
\xymatrix{
W_x \ar[d]_{i_x} \ar[r]_{j'} & \overline{\{x\}} \ar[d]^i \\
U \ar[r]^j & X
}
$$
を得る。$j_*i_{x, *}\mathcal{O}_{W_x} = i_*j'_*\mathcal{O}_{W_x}$ である。
左の垂直射は閉浸入なので、
$j_*i_{x, *}\mathcal{O}_{W_x}$ が連接であることと
$j'_*\mathcal{O}_{W_x}$ が連接であることは同値である。
\end{remark}

\begin{remark}
\label{local-cohomology-remark-no-finiteness-pushforward}
$X$ を局所 Noether 概形とし、$j : U \to X$ を、
補集合が $Z$ である開部分概形の包含とする。
$\mathcal{F}$ を連接 $\mathcal{O}_U$-加群とする。
$x \in \text{Ass}(\mathcal{F})$ と
$z \in Z \cap \overline{\{x\}}$ が存在し、
$\dim(\mathcal{O}_{\overline{\{x\}}, z}) \leq 1$ ならば、
$j_*\mathcal{F}$ は連接ではない。これを証明するため、
$\mathcal{O}_{X, z}$ のスペクトルへの平坦基底変換を行ってよい。
$X' = \overline{\{x\}}$ とおく。
仮定から $\mathcal{O}_{X' \cap U} \subset \mathcal{F}$ が従う。
したがって $j_*\mathcal{O}_{X' \cap U}$ が連接でないことを示せば十分である。
これは $X' = \{x, z\}$ なので明らかである。実際、
$j_*\mathcal{O}_{X' \cap U}$ は $\kappa(x)$ に対応し、
$\mathcal{O}_{X, z}$-加群としては、$x$ が閉点でないため有限ではあり得ない。

\medskip\noindent
実際、補題 \ref{local-cohomology-lemma-sharp-finiteness-pushforward} の逆も成り立つ。
整 Noether 概形の開浸入 $j : U \to X$ が与えられ、
$z \in X \setminus U$ と随伴素イデアル $\mathfrak p$ が存在し、
後者が完備化 $\mathcal{O}_{X, z}^\wedge$ のものであり、
$\dim(\mathcal{O}_{X, z}^\wedge/\mathfrak p) = 1$ を満たすものが存在するならば、
$j_*\mathcal{O}_U$ は連接ではない。
実際、局所環へ移り、$U$ を穴あきスペクトルまで拡大し、
完備化へ移れば、上の議論が有限でないことを与える。
\end{remark}

\begin{proposition}[Koll\'ar]
\label{local-cohomology-proposition-kollar}
\begin{reference}
特殊な場合については \cite{k-coherent} および
\cite[IV, 命題 7.2.2]{EGA} を参照せよ。
\end{reference}
\begin{slogan}
Hartogs の定理の弱い類似：Noether 概形上で、連接層を、
その層の台の中で補集合が余次元 2 である開集合へ制限したものは連接である。
\end{slogan}
$j : U \to X$ を局所 Noether 概形の開浸入とし、その補集合を $Z$ とする。
$\mathcal{F}$ を連接 $\mathcal{O}_U$-加群とする。
次の条件は同値である。
\begin{enumerate}
\item $j_*\mathcal{F}$ は連接である。
\item $x \in \text{Ass}(\mathcal{F})$、
$z \in Z \cap \overline{\{x\}}$、および任意の随伴素イデアル
$\mathfrak p$ で、完備化
$\mathcal{O}_{\overline{\{x\}}, z}^\wedge$ のものに対して
$\dim(\mathcal{O}_{\overline{\{x\}}, z}^\wedge/\mathfrak p) \geq 2$ である。
\end{enumerate}
\end{proposition}

\begin{proof}
(2) が成り立つならば、補題
\ref{local-cohomology-lemma-check-finiteness-pushforward-on-associated-points}、
注意 \ref{local-cohomology-remark-closure}、および
補題 \ref{local-cohomology-lemma-sharp-finiteness-pushforward} を組み合わせて (1) を得る。
(2) が成り立たないならば、注意
\ref{local-cohomology-remark-no-finiteness-pushforward} の議論
（および注意 \ref{local-cohomology-remark-closure}）により、
$j_*i_{x, *}\mathcal{O}_{W_x}$ は、ある
$x \in \text{Ass}(\mathcal{F})$ に対して有限でない。
したがって補題
\ref{local-cohomology-lemma-check-finiteness-pushforward-on-associated-points} により、
$j_*\mathcal{F}$ は連接でない。
\end{proof}

\begin{lemma}
\label{local-cohomology-lemma-kollar-finiteness-H1-local}
$A$ を Noether 環、$I \subset A$ をイデアルとする。
$Z = V(I)$ とおき、$M$ を有限 $A$-加群とする。
次の条件は同値である。
\begin{enumerate}
\item $H^1_Z(M)$ は有限 $A$-加群である。
\item すべての $\mathfrak p \in \text{Ass}(M)$ で
$\mathfrak p \not \in Z$ であるもの、およびすべての
$\mathfrak q \in V(\mathfrak p + I)$ に対して、
$(A/\mathfrak p)_\mathfrak q$ の完備化は次元 $1$ の随伴素イデアルをもたない。
\end{enumerate}
\end{lemma}

\begin{proof}
命題 \ref{local-cohomology-proposition-kollar} から、補題
\ref{local-cohomology-lemma-finiteness-pushforwards-and-H1-local} を介して直ちに従う。
\end{proof}

\noindent
次の補題の定式化には、条件 (1) と (2) が $X$ 上有限型の概形に
継承されるという利点がある。さらに、これは第
\ref{local-cohomology-section-finiteness-pushforward-II} 節で高次順像へ一般化する有限性の形である。

\begin{lemma}
\label{local-cohomology-lemma-finiteness-pushforward-general}
$X$ を局所 Noether 概形とする。
$j : U \to X$ を、補集合が $Z$ である開部分概形の包含とする。
$\mathcal{F}$ を連接 $\mathcal{O}_U$-加群とする。次を仮定する。
\begin{enumerate}
\item $X$ は普遍鎖状である。
\item すべての $z \in Z$ に対して、$\mathcal{O}_{X, z}$ の形式的ファイバーは
$(S_1)$ を満たす。
\end{enumerate}
この状況で、次の条件は同値である。
\begin{enumerate}
\item[(a)] $x \in \text{Ass}(\mathcal{F})$ および
$z \in Z \cap \overline{\{x\}}$ に対して
$\dim(\mathcal{O}_{\overline{\{x\}}, z}) \geq 2$ である。
\item[(b)] $j_*\mathcal{F}$ は連接である。
\end{enumerate}
\end{lemma}

\begin{proof}
$x \in \text{Ass}(\mathcal{F})$ とする。命題 \ref{local-cohomology-proposition-kollar}
により、$A = \mathcal{O}_{\overline{\{x\}}, z}$ が、その完備化の
随伴素イデアルに関する同命題の条件を満たすことと
$\dim(A) \geq 2$ であることが同値であると示せば十分である。
$A$ は普遍鎖状であり（これは明らかである）、
その形式的ファイバーは $(S_1)$ を満たす。後者は
More on Algebra, 補題 \ref{more-algebra-lemma-formal-fibres-normal} および
命題 \ref{more-algebra-proposition-finite-type-over-P-ring} から従う。
$\mathfrak p' \subset A^\wedge$ を随伴素イデアルとする。
$A \to A^\wedge$ は平坦なので、
Algebra, 補題 \ref{algebra-lemma-bourbaki} により、
$\mathfrak p'$ は $(0) \subset A$ の上にある。
形式的ファイバー $A^\wedge \otimes_A F$ は $(S_1)$ を満たす。
ここで $F$ は $A$ の分数体である。
したがって $\mathfrak p'$ は極小素イデアルである。
Algebra, 補題 \ref{algebra-lemma-criterion-no-embedded-primes} を参照せよ。
$A$ は普遍鎖状なので、More on Algebra, 命題
\ref{more-algebra-proposition-ratliff} により形式的に鎖状である。
したがって $\dim(A^\wedge/\mathfrak p') = \dim(A)$ となり、
同値性が示された。
\end{proof}






\section{深さと次元}
\label{local-cohomology-section-dept-dimension}

\noindent
いくつかの補助補題を与える。

\begin{lemma}
\label{local-cohomology-lemma-ideal-depth-function}
$A$ を Noether 環とし、$I \subset A$ をイデアルとする。
$M$ を有限 $A$-加群とし、$\mathfrak p \in V(I)$ を素イデアルとする。
$e = \text{depth}_{IA_\mathfrak p}(M_\mathfrak p) < \infty$ と仮定する。
このとき空でない開集合 $U \subset V(\mathfrak p)$ が存在して、
$\text{depth}_{IA_\mathfrak q}(M_\mathfrak q) \geq e$ が
すべての $\mathfrak q \in U$ に対して成り立つ。
\end{lemma}

\begin{proof}
深さの定義により $IM_\mathfrak p \not = M_\mathfrak p$ であり、
$M_\mathfrak p$-正則列
$f_1, \ldots, f_e \in IA_\mathfrak p$ が存在する。$A$ を主局所化で
置き換えると、$f_1, \ldots, f_e \in I$ が $M$-正則列をなすと仮定してよい。
Algebra, 補題 \ref{algebra-lemma-regular-sequence-in-neighbourhood} を参照せよ。
加群 $M' = M/IM$ を考える。$\mathfrak p \in \text{Supp}(M')$ であり、
有限加群の台は閉だから、$V(\mathfrak p) \subset \text{Supp}(M')$ を得る。
したがって $\mathfrak q \in V(\mathfrak p)$ に対して
$IM_\mathfrak q \not = M_\mathfrak q$ である。ゆえに局所化が完全であることと
深さの定義から、$\text{depth}_{IA_\mathfrak q}(M_\mathfrak q) \geq e$ が
任意の $\mathfrak q \in V(I)$ に対して成り立つ。
\end{proof}

\begin{lemma}
\label{local-cohomology-lemma-depth-function}
$A$ を Noether 環とし、$M$ を有限 $A$-加群とする。
$\mathfrak p$ を素イデアルとし、
$e = \text{depth}_{A_\mathfrak p}(M_\mathfrak p) < \infty$ と仮定する。
このとき空でない開集合 $U \subset V(\mathfrak p)$ が存在して、
$\text{depth}_{A_\mathfrak q}(M_\mathfrak q) \geq e$ が
すべての $\mathfrak q \in U$ に対して成り立ち、さらに
有限個を除くすべての $\mathfrak q \in U$ に対して
$\text{depth}_{A_\mathfrak q}(M_\mathfrak q) > e$ が成り立つ。
\end{lemma}

\begin{proof}
深さの定義により $\mathfrak p M_\mathfrak p \not = M_\mathfrak p$ であり、
$M_\mathfrak p$-正則列
$f_1, \ldots, f_e \in \mathfrak p A_\mathfrak p$ が存在する。$A$ を
主局所化で置き換えると、$f_1, \ldots, f_e \in \mathfrak p$ が
$M$-正則列をなすと仮定してよい。Algebra, 補題
\ref{algebra-lemma-regular-sequence-in-neighbourhood} を参照せよ。
加群 $M' = M/(f_1, \ldots, f_e)M$ を考える。
$\mathfrak p \in \text{Supp}(M')$ であり、有限加群の台は閉だから、
$V(\mathfrak p) \subset \text{Supp}(M')$ を得る。したがって
$\mathfrak q \in V(\mathfrak p)$ に対して
$\mathfrak q M_\mathfrak q \not = M_\mathfrak q$ である。ゆえに局所化が
完全であることと深さの定義から、
$\text{depth}_{A_\mathfrak q}(M_\mathfrak q) \geq e$ が
任意の $\mathfrak q \in V(I)$ に対して成り立つ。
さらに $\mathfrak q$ が加群 $M'$ の随伴素イデアルでなければ、深さは増加する。
したがって最後の主張は Algebra, 補題 \ref{algebra-lemma-finite-ass} から従う。
\end{proof}

\begin{lemma}
\label{local-cohomology-lemma-finite-nr-points-next-S}
$X$ を双対化複体 $\omega_X^\bullet$ をもつ Noether 概形とする。
$\mathcal{F}$ を連接 $\mathcal{O}_X$-加群とし、$k \geq 0$ を整数とする。
$\mathcal{F}$ が $(S_k)$ を満たすと仮定する。このとき次を満たす点
$x \in X$ は有限個しかない。
$$
\text{depth}(\mathcal{F}_x) = k
\quad\text{かつ}\quad
\dim(\text{Supp}(\mathcal{F}_x)) > k
$$
\end{lemma}

\begin{proof}
この補題を $k$ に関する帰納法で証明する。基底の場合 $k = 0$ は、
$\mathcal{F}$ の埋め込まれた随伴点が有限個であるという主張であり、
これは Divisors, 補題 \ref{divisors-lemma-finite-ass} から従う。

\medskip\noindent
$k > 0$ とし、より小さいすべての $k$ に対して結果が成り立つと仮定する。
$X$ は有限個のアフィン開集合で被覆できるので、
$X = \Spec(A)$ はアフィンであると仮定してよい。このとき $\mathcal{F}$ は、
連接 $\mathcal{O}_X$-加群であり、有限 $A$-加群 $M$ に付随する。
この加群は $(S_k)$ を満たす。以下では Algebra, 補題
\ref{algebra-lemma-one-equation-module} および
\ref{algebra-lemma-depth-drops-by-one} を断りなく用いる。

\medskip\noindent
$f \in A$ を $M$ 上の非零因子とする。このとき $M/fM$ は $(S_{k - 1})$ を満たす。
帰納法により、
素イデアル $\mathfrak p \in V(f)$ で、
$\text{depth}((M/fM)_\mathfrak p) = k - 1$ かつ
$\dim(\text{Supp}((M/fM)_\mathfrak p)) > k - 1$ を満たすものは有限個しかない。
これらはちょうど、素イデアル $\mathfrak p \in V(f)$ で、
$\text{depth}(M_\mathfrak p) = k$ かつ
$\dim(\text{Supp}(M_\mathfrak p)) > k$ を満たすものである。したがって有限性を証明する際に、
$A$ を $A_f$ で、$M$ を $M_f$ で置き換えてよい。

\medskip\noindent
$M$ は $(S_k)$ を満たし、$k > 0$ だから、$M$ は埋め込まれた随伴素イデアルを
もたない（Algebra, 補題 \ref{algebra-lemma-criterion-no-embedded-primes}）。
したがって $\text{Ass}(M)$ は $M$ の台の生成点全体である。
ゆえに Dualizing Complexes, 補題 \ref{dualizing-lemma-CM-open} により、集合
$U = \{\mathfrak q \mid M_\mathfrak q\text{ は Cohen--Macaulay である}\}$
は $\text{Ass}(M)$ を含む開集合である。素イデアル回避
（Algebra, 補題 \ref{algebra-lemma-silly}）により、
$f \in A$ を、$f \not \in \mathfrak p$ が
$\mathfrak p \in \text{Ass}(M)$ に対して成り立ち、かつ
$D(f) \subset U$ となるように選べる。
すると $f$ は $M$ 上の非零因子である
（Algebra, 補題 \ref{algebra-lemma-ass-zero-divisors}）。上で述べたように
$A$ を $A_f$ で、$M$ を $M_f$ で置き換えると、$M$ は Cohen--Macaulay
であることが分かる。したがってすべての $\mathfrak q \subset A$ に対して
$\dim(M_\mathfrak q) = \text{depth}(M_\mathfrak q)$ である。
ゆえに補題で記述した集合は空であり、まして有限である。
\end{proof}

\begin{lemma}
\label{local-cohomology-lemma-sitting-in-degrees}
$(A, \mathfrak m)$ を正規化された双対化複体 $\omega_A^\bullet$ をもつ
Noether 局所環とする。$M$ を有限 $A$-加群とする。
$E^i = \text{Ext}_A^{-i}(M, \omega_A^\bullet)$ とおく。このとき次が成り立つ。
\begin{enumerate}
\item $E^i$ は有限 $A$-加群であり、
$0 \leq i \leq \dim(\text{Supp}(M))$ の場合にのみ零でない。
\item $\dim(\text{Supp}(E^i)) \leq i$,
\item $\text{depth}(M)$ は、整数 $\delta \geq 0$ で
$E^\delta \not = 0$ を満たすもののうち最小である。
\item $\mathfrak p \in \text{Supp}(E^0 \oplus \ldots \oplus E^i)
\Leftrightarrow
\text{depth}_{A_\mathfrak p}(M_\mathfrak p) + \dim(A/\mathfrak p) \leq i$,
\item $E^i$ の零化イデアルは $H^i_\mathfrak m(M)$ の零化イデアルに等しい。
\end{enumerate}
\end{lemma}

\begin{proof}
(1)、(2)、(3) は Dualizing Complexes, 補題
\ref{dualizing-lemma-sitting-in-degrees} の主張そのものである。
素イデアル $\mathfrak p$ を $A$ のものとすると、
$(\omega_A^\bullet)_\mathfrak p[-\dim(A/\mathfrak p)]$ は
$A_\mathfrak p$ の正規化された双対化複体である。Dualizing Complexes, 補題
\ref{dualizing-lemma-dimension-function} を参照せよ。したがって
$$
E^i_\mathfrak p =
\text{Ext}^{-i}_A(M, \omega_A^\bullet)_\mathfrak p =
\text{Ext}^{-i + \dim(A/\mathfrak p)}_{A_\mathfrak p}
(M_\mathfrak p, (\omega_A^\bullet)_\mathfrak p[-\dim(A/\mathfrak p)])
$$
は、$i - \dim(A/\mathfrak p) < \text{depth}_{A_\mathfrak p}(M_\mathfrak p)$
ならば零であり、$i = \dim(A/\mathfrak p) +
\text{depth}_{A_\mathfrak p}(M_\mathfrak p)$ ならば零でない。
これは $A_\mathfrak p$ 上で (3) を用いることによる。これで (4) が示された。
$E$ を $A$ の剰余体の入射包絡とすると、
$$
\Hom_A(H^i_\mathfrak m(M), E) =
\text{Ext}^{-i}_A(M, \omega_A^\bullet)^\wedge =
(E^i)^\wedge =
E^i \otimes_A A^\wedge
$$
を得る。これは局所双対定理（Dualizing Complexes, 補題
\ref{dualizing-lemma-special-case-local-duality} の形）による。
$A \to A^\wedge$ は忠実平坦だから、Matlis 双対性
（Dualizing Complexes, 命題 \ref{dualizing-proposition-matlis}）により (5) が従う。
\end{proof}






\section{局所コホモロジーの零化イデアル I}
\label{local-cohomology-section-annihilators}

\noindent
この節では Faltings による結果を論じる。
\cite{Faltings-annulators} を参照せよ。

\begin{proposition}
\label{local-cohomology-proposition-annihilator}
\begin{reference}
\cite{Faltings-annulators}.
\end{reference}
$A$ を双対化複体をもつ Noether 環とする。
$T \subset T' \subset \Spec(A)$ を特殊化で安定な部分集合とする。
$s \geq 0$ を整数とし、$M$ を有限 $A$-加群とする。次の条件は同値である。
\begin{enumerate}
\item イデアル $J \subset A$ で $V(J) \subset T'$ を満たし、
$J$ が $H^i_T(M)$ を $i \leq s$ に対して零化するものが存在する。
\item すべての $\mathfrak p \not \in T'$ と、
$\mathfrak q \in T$ で $\mathfrak p \subset \mathfrak q$ を満たすものに対して、
$$
\text{depth}_{A_\mathfrak p}(M_\mathfrak p) +
\dim((A/\mathfrak p)_\mathfrak q) > s
$$
\end{enumerate}
\end{proposition}

\begin{proof}
$\omega_A^\bullet$ を双対化複体とし、$\delta$ をその次元関数とする。
Dualizing Complexes, 第 \ref{dualizing-section-dimension-function} 節を参照せよ。
以下では有限 $A$-加群
$$
E^i = \Ext_A^i(M, \omega_A^\bullet)
$$
が重要な役割を果たす。$\mathfrak p \subset A$ に対し、$H^i_\mathfrak p$ により
$A_\mathfrak p$-加群の $\mathfrak pA_\mathfrak p$ に関する
局所コホモロジーを表すことにする。このとき $\mathfrak pA_\mathfrak p$-進完備化を考えると、
$$
(E^i)_\mathfrak p =
\Ext^{\delta(\mathfrak p) + i}_{A_\mathfrak p}(M_\mathfrak p,
(\omega_A^\bullet)_\mathfrak p[-\delta(\mathfrak p)])
$$
は
$$
H^{-\delta(\mathfrak p) - i}_{\mathfrak p}(M_\mathfrak p)
$$
の Matlis 双対である。これは Dualizing Complexes, 補題
\ref{dualizing-lemma-special-case-local-duality} による。
特に次の事実が従う。イデアル $J \subset A$ が
$(E^i)_\mathfrak p$ を零化することと、$J$ が
$H^{-\delta(\mathfrak p) - i}_{\mathfrak p}(M_\mathfrak p)$ を零化することは同値である。

\medskip\noindent
$T_n = \{\mathfrak p \in T \mid \delta(\mathfrak p) \leq n\}$ とおく。
$\delta$ は有界関数だから、$T_a = \emptyset$ が $a \ll 0$ に対して成り立ち、
$T_b = T$ が $b \gg 0$ に対して成り立つ。

\medskip\noindent
(2) を仮定し、(1) のような $J$ の存在を示そう。そのため二重帰納法を用いる。
$i \leq s$ に対して、帰納法の仮定 $IH_i$ を次のように取る。
$H^a_T(M)$ は、ある $J \subset A$ で $V(J) \subset T'$ を満たすものにより、
$0 \leq a \leq i$ に対して零化される。$IH_0$ の場合は明らかである。
実際、$H^0_T(M)$ は $M$ の部分加群だから有限であり、したがって
あるイデアル $J$ で $V(J) \subset T$ を満たすものにより零化される。

\medskip\noindent
帰納段階を考える。$IH_{i - 1}$ がある $0 < i \leq s$ に対して成り立つと仮定する。
$J'$ を、$V(J') \subset T'$ を満たし、$H^a_T(M)$ を
$0 \leq a \leq i - 1$ に対して零化するものとして選ぶ
（帰納法の仮定により可能である）。
$n$ に関する降下帰納法により、イデアル $J$ で $V(J) \subset T'$ を満たし、
$J H^i_T(M)$ の随伴素イデアルが $T_n$ に属するものが存在することを示す。
$n \ll 0$ なら、これは $JH^i_T(M) = 0$ を意味する
（Algebra, 補題 \ref{algebra-lemma-ass-zero}）。したがって $IH_i$ が成り立つ。
基底の場合 $n \gg 0$ は明らかである。この場合 $T = T_n$ であり、
$H^i_T(M)$ のすべての随伴素イデアルは $T$ に属するからである。

\medskip\noindent
そこで $J$ が $n$ に対する性質をもつものとして与えられたと仮定する。
$\mathfrak q \in T_n$ とし、$T_\mathfrak q \subset \Spec(A_\mathfrak q)$ を
$T$ の逆像とする。補題 \ref{local-cohomology-lemma-torsion-change-rings} により
$H^j_T(M)_\mathfrak q = H^j_{T_\mathfrak q}(M_\mathfrak q)$
である。補題 \ref{local-cohomology-lemma-local-cohomology-ss} のスペクトル系列
$$
H_\mathfrak q^p(H^q_{T_\mathfrak q}(M_\mathfrak q))
\Rightarrow
H^{p + q}_\mathfrak q(M_\mathfrak q)
$$
を考える。以下で、イデアル $J'' \subset A$ で $V(J'') \subset T'$ を満たし、
$H^i_\mathfrak q(M_\mathfrak q)$ を $J''$ により、すべての
$\mathfrak q \in T_n \setminus T_{n - 1}$ に対して零化するものを見いだす。
$J (J')^i J''$ が $n - 1$ に対して使えると主張する。
実際、$\mathfrak q \in T_n \setminus T_{n - 1}$ とする。
上のスペクトル系列はフィルトレーション
$$
E_\infty^{0, i} = E_{i + 2}^{0, i} \subset \ldots \subset E_3^{0, i} \subset
E_2^{0, i} = H^0_\mathfrak q(H^i_{T_\mathfrak q}(M_\mathfrak q))
$$
を定める。加群 $E_\infty^{0, i}$ は $J''$ により零化される。
部分商 $E_j^{0, i}/E_{j + 1}^{0, i}$ は $i + 1 \geq j \geq 2$ に対して
$J'$ により零化される。実際、$d_j^{0, i}$ の標的は
$$
H^j_\mathfrak q(H^{i - j + 1}_{T_\mathfrak q}(M_\mathfrak q)) =
H^j_\mathfrak q(H^{i - j + 1}_T(M)_\mathfrak q)
$$
の部分商であり、$H^{i - j + 1}_T(M)_\mathfrak q$ は $J'$ の選び方により
$J'$ によって零化される。最後に $J$ の選び方から
$J H^i_T(M)_\mathfrak q \subset H^0_\mathfrak q(H^i_T(M)_\mathfrak q)$
である。なぜなら $\Spec(A_\mathfrak q)$ の閉点でない点では
$\delta$ の値がより大きいからである。したがって望みどおり、
$\mathfrak q$ は $J(J')^iJ'' H^i_T(M)$ の随伴素イデアルではあり得ない。

\medskip\noindent
冒頭の注意から、$J''$ は
$$
(E^{-\delta(\mathfrak q) - i})_\mathfrak q =
(E^{-n - i})_\mathfrak q
$$
をすべての $\mathfrak q \in T_n \setminus T_{n - 1}$ に対して零化すべきである。
しかし $J''$ が一つの $\mathfrak q$ に対して使えるなら、その $\mathfrak q$ の
ある開近傍内のすべての $\mathfrak q$ に対して使える。これは加群
$E^{-n - i}$ が有限だからである。
$\Spec(A)$ の任意の部分集合は誘導位相に関して Noether である
（Topology, 補題 \ref{topology-lemma-Noetherian}）から、
$J''$ の存在を一つの $\mathfrak q$ に対して示せば十分である。

\medskip\noindent
Ext 加群は有限だから、$J''$ の存在は
$$
\text{Supp}(E^{-n - i}) \cap \Spec(A_\mathfrak q) \subset T'.
$$
と同値である。これは、$E^{-n - i}$ の局所化が、すべての
$\mathfrak p \subset \mathfrak q$、$\mathfrak p \not \in T'$ において零であることと同値である。
$A_\mathfrak p$ 上の局所双対性を用いると、
$$
H^{i + n - \delta(\mathfrak p)}_\mathfrak p(M_\mathfrak p) =
H^{i - \dim((A/\mathfrak p)_\mathfrak q)}_\mathfrak p(M_\mathfrak p)
$$
が零であることを示せばよい（ここで $\delta$ が次元関数であることを用いた）。
これは命題の仮定、$i \leq s$、および Dualizing Complexes, 補題
\ref{dualizing-lemma-depth} により消滅する。

\medskip\noindent
逆向きの含意を示すため、(2) が成り立たないと仮定し、上の議論を逆にたどる。
まず $\mathfrak q \in T$、$\mathfrak p \subset \mathfrak q$、
$\mathfrak p \not \in T'$ を、
$$
i = \text{depth}_{A_\mathfrak p}(M_\mathfrak p) +
\dim((A/\mathfrak p)_\mathfrak q) \leq s
$$
が最小となるように選ぶ。このとき
$H^{i - \dim((A/\mathfrak p)_\mathfrak q)}_\mathfrak p(M_\mathfrak p)$
は零でない。これは Dualizing Complexes, 補題
\ref{dualizing-lemma-depth} の非消滅による。$n = \delta(\mathfrak q)$ とおく。
するとイデアル $J \subset A$ で $V(J) \subset T'$ を満たし、
$J(E^{-n - i})_\mathfrak q = 0$ となるものは存在しない。
したがって $H^i_\mathfrak q(M_\mathfrak q)$ は、イデアル $J \subset A$ で
$V(J) \subset T'$ を満たすものにより零化されない。$i$ の最小性と上に表示した
スペクトル系列から、加群 $H^i_T(M)_\mathfrak q$ は
イデアル $J \subset A$ で $V(J) \subset T'$ を満たすものにより零化されない。
したがって $H^i_T(M)$ はイデアル $J \subset A$ で
$V(J) \subset T'$ を満たすものにより零化されない。これで命題の証明が終わる。
\end{proof}

\begin{lemma}
\label{local-cohomology-lemma-kill-local-cohomology-at-prime}
$I$ を Noether 環 $A$ のイデアルとする。$M$ を有限 $A$-加群、
$\mathfrak p \subset A$ を素イデアル、$s \geq 0$ を整数とする。次を仮定する。
\begin{enumerate}
\item $A$ は双対化複体をもつ。
\item $\mathfrak p \not \in V(I)$ である。
\item すべての素イデアル $\mathfrak p' \subset \mathfrak p$ と、
$\mathfrak q \in V(I)$ で $\mathfrak p' \subset \mathfrak q$ を満たすものに対して、
$$
\text{depth}_{A_{\mathfrak p'}}(M_{\mathfrak p'}) +
\dim((A/\mathfrak p')_\mathfrak q) > s
$$
\end{enumerate}
このとき $f \in A$、$f \not \in \mathfrak p$ であって、
$H^i_{V(I)}(M)$ を $i \leq s$ に対して零化するものが存在する。
\end{lemma}

\begin{proof}
集合
$$
T = V(I)
\quad\text{および}\quad
T' = \bigcup\nolimits_{f \in A, f \not \in \mathfrak p} V(f)
$$
を考える。これらは $\Spec(A)$ の特殊化で安定な部分集合である。
$T \subset T'$ かつ $\mathfrak p \not \in T'$ であることに注意する。
仮定 (3) は命題 \ref{local-cohomology-proposition-annihilator} の仮定 (2) が成り立つことを述べている。
したがって $J \subset A$ で $V(J) \subset T'$ を満たし、
$J H^i_{V(I)}(M) = 0$ が $i \leq s$ に対して成り立つものを取れる。
$f \in A$、$f \not \in \mathfrak p$ を $V(J) \subset V(f)$ となるように選ぶ。
$f$ のある冪は $H^i_{V(I)}(M)$ を $i \leq s$ に対して零化する。
\end{proof}





\section{局所コホモロジーの有限性 II}
\label{local-cohomology-section-finiteness-II}

\noindent
第 \ref{local-cohomology-section-finiteness} 節で始めた局所コホモロジーの有限性の議論を続ける。
Faltings 零化イデアル定理を用いると、次の基本的結果が容易に証明できる。

\begin{proposition}
\label{local-cohomology-proposition-finiteness}
\begin{reference}
\cite{Faltings-annulators}.
\end{reference}
$A$ を双対化複体をもつ Noether 環とする。
$T \subset \Spec(A)$ を特殊化で安定な部分集合とする。
$s \geq 0$ を整数とし、$M$ を有限 $A$-加群とする。次の条件は同値である。
\begin{enumerate}
\item $H^i_T(M)$ は有限 $A$-加群である（$i \leq s$）。
\item すべての $\mathfrak p \not \in T$ と、$\mathfrak q \in T$ で
$\mathfrak p \subset \mathfrak q$ を満たすものに対して、
$$
\text{depth}_{A_\mathfrak p}(M_\mathfrak p) +
\dim((A/\mathfrak p)_\mathfrak q) > s
$$
\end{enumerate}
\end{proposition}

\begin{proof}
命題 \ref{local-cohomology-proposition-annihilator} と補題
\ref{local-cohomology-lemma-check-finiteness-local-cohomology-by-annihilator} の形式的帰結である。
\end{proof}

\noindent
後で用いるいくつかの補題のほか、この節の残りでは、命題
\ref{local-cohomology-proposition-finiteness} における「$A$ が双対化複体をもつ」という条件を
どの程度弱められるかを論じる。大まかに言えば、$A$ の形式的ファイバーが
$(S_n)$ を満たすと仮定する必要があり、ここで $n$ は十分大きい。

\medskip\noindent
$A$ を Noether 環、$I \subset A$ をイデアルとする。
$X = \Spec(A)$ および $Z = V(I) \subset X$ とおき、$M$ を有限 $A$-加群とする。
次のように定義する。
\begin{equation}
\label{local-cohomology-equation-cutoff}
s_{A, I}(M) =
\min \{
\text{depth}_{A_\mathfrak p}(M_\mathfrak p) + \dim((A/\mathfrak p)_\mathfrak q)
\mid
\mathfrak p \in X \setminus Z, \mathfrak q \in Z,
\mathfrak p \subset \mathfrak q
\}
\end{equation}
深さに関する規約では $0$ の深さは $\infty$ であるから、素イデアル
$\mathfrak p$ のうち $M$ の台に属するものだけを考えればよい。$s_{A, I}(M)$ はこの状況の
重要な不変量であることが後に分かる。

\begin{lemma}
\label{local-cohomology-lemma-cutoff}
$A \to B$ を Noether 環の有限準同型とする。$I \subset A$ をイデアルとし、
$J = IB$ とおく。$M$ を有限 $B$-加群とする。$A$ が普遍鎖状ならば
$s_{B, J}(M) = s_{A, I}(M)$ である。
\end{lemma}

\begin{proof}
$\mathfrak p \subset \mathfrak q \subset A$ を、$I \subset \mathfrak q$ かつ
$I \not \subset \mathfrak p$ を満たす素イデアルとする。$A \to B$ は有限なので、
$\mathfrak p_i$ を $\mathfrak p$ の上にある素イデアルとすれば、その個数は有限である。
Algebra, 補題 \ref{algebra-lemma-depth-goes-down-finite} により
$$
\text{depth}(M_\mathfrak p) = \min \text{depth}(M_{\mathfrak p_i})
$$
である。$\mathfrak p_i \subset \mathfrak q_{ij}$ を $\mathfrak q$ の上にある
素イデアルとする。$A \to B$ に対する上昇定理
（Algebra, 補題 \ref{algebra-lemma-integral-going-up}）により、少なくとも一つの
$\mathfrak q_{ij}$ が存在し、これは各 $i$ について成り立つ。このとき
$$
\dim((B/\mathfrak p_i)_{\mathfrak q_{ij}}) =
\dim((A/\mathfrak p)_\mathfrak q)
$$
が次元公式により成り立つ。Algebra, 補題
\ref{algebra-lemma-dimension-formula} を参照せよ。したがって、$s_{B, J}(M)$ の定義に
用いる量を組 $(\mathfrak p_i, \mathfrak q_{ij})$ にわたって最小化した値は、
組 $(\mathfrak p, \mathfrak q)$ に対する量に等しい。これで補題が従う。
\end{proof}

\begin{lemma}
\label{local-cohomology-lemma-change-completion}
$A$ を双対化複体をもつ Noether 環とし、$I \subset A$ をイデアルとする。
$M$ を有限 $A$-加群とする。$A', M'$ を $I$-進完備化とし、それぞれ $A, M$ から得られるものとする。
$\mathfrak p' \subset \mathfrak q'$ を $A'$ の素イデアルで、
$\mathfrak q' \in V(IA')$ を満たし、$\mathfrak p \subset \mathfrak q$ の上にあるものとする。
ここで後者は $A$ の素イデアルである。このとき
$$
\text{depth}_{A_{\mathfrak p'}}(M'_{\mathfrak p'})
\geq
\text{depth}_{A_\mathfrak p}(M_\mathfrak p)
$$
かつ
$$
\text{depth}_{A_{\mathfrak p'}}(M'_{\mathfrak p'}) +
\dim((A'/\mathfrak p')_{\mathfrak q'}) =
\text{depth}_{A_\mathfrak p}(M_\mathfrak p) +
\dim((A/\mathfrak p)_\mathfrak q)
$$
\end{lemma}

\begin{proof}
次が成り立つ。
$$
\text{depth}(M'_{\mathfrak p'}) =
\text{depth}(M_\mathfrak p) +
\text{depth}(A'_{\mathfrak p'}/\mathfrak p A'_{\mathfrak p'})
\geq \text{depth}(M_\mathfrak p)
$$
これは $A \to A'$ の平坦性による。Algebra, 補題
\ref{algebra-lemma-apply-grothendieck-module} を参照せよ。
$A \to A'$ のファイバーは Cohen--Macaulay である
（Dualizing Complexes, 補題
\ref{dualizing-lemma-dualizing-gorenstein-formal-fibres} および
More on Algebra, 第
\ref{more-algebra-section-properties-formal-fibres} 節）から、
$\text{depth}(A'_{\mathfrak p'}/\mathfrak p A'_{\mathfrak p'}) =
\dim(A'_{\mathfrak p'}/\mathfrak p A'_{\mathfrak p'})$。
したがって
\begin{align*}
\text{depth}(M'_{\mathfrak p'}) +
\dim((A'/\mathfrak p')_{\mathfrak q'})
& =
\text{depth}(M_\mathfrak p) +
\dim(A'_{\mathfrak p'}/\mathfrak p A'_{\mathfrak p'}) +
\dim((A'/\mathfrak p')_{\mathfrak q'}) \\
& =
\text{depth}(M_\mathfrak p) +
\dim((A'/\mathfrak p A')_{\mathfrak q'}) \\
& =
\text{depth}(M_\mathfrak p) +
\dim((A/\mathfrak p)_\mathfrak q)
\end{align*}
を得る。第2の等号は $A'$ が鎖状であることによる。第3の等号は
More on Algebra, 補題 \ref{more-algebra-lemma-completion-dimension} による。
実際、$(A/\mathfrak p)_\mathfrak q$ と $(A'/\mathfrak p A')_{\mathfrak q'}$ は
同じ $I$-進完備化をもつ。
\end{proof}





\begin{lemma}
\label{local-cohomology-lemma-cutoff-completion}
$A$ を普遍鎖状 Noether 局所環とする。
$I \subset A$ をイデアルとし、$M$ を有限 $A$-加群とする。このとき
$$
s_{A, I}(M) \geq s_{A^\wedge, I^\wedge}(M^\wedge)
$$
$A$ の形式的ファイバーが $(S_n)$ を満たすならば、
$\min(n + 1, s_{A, I}(M)) \leq s_{A^\wedge, I^\wedge}(M^\wedge)$ である。
\end{lemma}

\begin{proof}
$X = \Spec(A)$、$X^\wedge = \Spec(A^\wedge)$、$Z = V(I) \subset X$ および
$Z^\wedge = V(I^\wedge)$ とおく。
$\mathfrak p' \subset \mathfrak q' \subset A^\wedge$ を
$\mathfrak p' \not \in Z^\wedge$ かつ
$\mathfrak q' \in Z^\wedge$ を満たす素イデアルとする。
$\mathfrak p \subset \mathfrak q$ を対応する $A$ の素イデアルとする。このとき
$\mathfrak p \not \in Z$ かつ $\mathfrak q \in Z$ である。図式は
$$
\xymatrix{
\mathfrak p' \ar[r] & \mathfrak q' \ar[r] & A^\wedge \\
\mathfrak p \ar[r] \ar@{-}[u] &
\mathfrak q \ar[r] \ar@{-}[u] & A \ar[u]
}
$$
次のようにおく。
\begin{align*}
a & = \dim(A/\mathfrak p) = \dim(A^\wedge/\mathfrak pA^\wedge),\\
b & = \dim(A/\mathfrak q) = \dim(A^\wedge/\mathfrak qA^\wedge),\\
a' & = \dim(A^\wedge/\mathfrak p'),\\
b' & = \dim(A^\wedge/\mathfrak q')
\end{align*}
これらの等式は More on Algebra, 補題
\ref{more-algebra-lemma-completion-dimension} による。また
\begin{align*}
p & = \dim(A^\wedge_{\mathfrak p'}/\mathfrak p A^\wedge_{\mathfrak p'}) =
\dim((A^\wedge/\mathfrak p A^\wedge)_{\mathfrak p'}) \\
q & = \dim(A^\wedge_{\mathfrak q'}/\mathfrak p A^\wedge_{\mathfrak q'}) =
\dim((A^\wedge/\mathfrak q A^\wedge)_{\mathfrak q'})
\end{align*}
とおく。$A$ は普遍鎖状であるから、
$A^\wedge/\mathfrak pA^\wedge = (A/\mathfrak p)^\wedge$
は次元 $a$ の等次元環である（More on Algebra, 命題
\ref{more-algebra-proposition-ratliff}）。
したがって $a = a' + p$ である。同様に $b = b' + q$ である。
Algebra, 補題 \ref{algebra-lemma-apply-grothendieck-module} を平坦な局所環準同型
$A_\mathfrak p \to A^\wedge_{\mathfrak p'}$
に適用すると、
$$
\text{depth}(M^\wedge_{\mathfrak p'})
=
\text{depth}(M_\mathfrak p) +
\text{depth}(A^\wedge_{\mathfrak p'} / \mathfrak p A^\wedge_{\mathfrak p'})
$$
を得る。$s_{A, I}(M)$ の定義で最小化している量は
$$
s(\mathfrak p, \mathfrak q) =
\text{depth}(M_\mathfrak p) + \dim((A/\mathfrak p)_\mathfrak q) =
\text{depth}(M_\mathfrak p) + a - b
$$
である（最後の等号は $A$ が鎖状であることによる）。
$s_{A^\wedge, I^\wedge}(M^\wedge)$ の定義で最小化している量は
$$
s(\mathfrak p', \mathfrak q') =
\text{depth}(M^\wedge_{\mathfrak p'}) +
\dim((A^\wedge/\mathfrak p')_{\mathfrak q'}) =
\text{depth}(M^\wedge_{\mathfrak p'}) + a' - b'
$$
である（最後の等号は $A^\wedge$ が鎖状であることによる）。
以上で証明を始めるための記号が整った。

\medskip\noindent
$\mathfrak p \subset \mathfrak q \subset A$ を、
$\mathfrak p \not \in Z$ かつ $\mathfrak q \in Z$ を満たし、
$s_{A, I}(M) = s(\mathfrak p, \mathfrak q)$ となる素イデアルとする。
$\mathfrak q'$ を $\mathfrak q A^\wedge$ の上で極小に選び、
$\mathfrak p' \subset \mathfrak q'$ を $\mathfrak p A^\wedge$ の上で極小に選べる
（$A \to A^\wedge$ に対する下降定理を用いる）。
すると上のような四つの素イデアルが得られ、$p = 0$ かつ $q = 0$ である。
さらに
$\text{depth}(A^\wedge_{\mathfrak p'} / \mathfrak p A^\wedge_{\mathfrak p'})=0$
も $p = 0$ から従う。したがって
$s(\mathfrak p', \mathfrak q') = s(\mathfrak p, \mathfrak q)$。
これにより最初の不等式を得る。

\medskip\noindent
次に、$A$ の形式的ファイバーが $(S_n)$ を満たすと仮定する。このとき
$\text{depth}(A^\wedge_{\mathfrak p'} / \mathfrak p A^\wedge_{\mathfrak p'})
\geq \min(n, p)$。
したがって
$$
s(\mathfrak p', \mathfrak q') \geq
s(\mathfrak p, \mathfrak q) + q + \min(n, p) - p \geq
s_{A, I}(M) + q + \min(n, p) - p
$$
である。問題が生じ得るのは $p > n$ の場合だけである。この場合は
\begin{align*}
s(\mathfrak p', \mathfrak q')
& =
\text{depth}(M^\wedge_{\mathfrak p'}) +
\dim((A^\wedge/\mathfrak p')_{\mathfrak q'}) \\
& =
\text{depth}(M_\mathfrak p) +
\text{depth}(A^\wedge_{\mathfrak p'} / \mathfrak p A^\wedge_{\mathfrak p'}) +
\dim((A^\wedge/\mathfrak p')_{\mathfrak q'}) \\
& \geq
0 + n + 1
\end{align*}
となる。実際、$(A^\wedge/\mathfrak p')_{\mathfrak q'}$ は少なくとも二つの
素イデアルをもつ。これで第二の不等式が証明された。
\end{proof}

\noindent
次の補題の証明法はより一般的な状況にも適用できるが、そこから得られる
より強い結果は、後の定理 \ref{local-cohomology-theorem-finiteness} に含まれる。

\begin{lemma}
\label{local-cohomology-lemma-local-annihilator}
\begin{reference}
これは \cite[Satz 1]{Faltings-annulators} の特殊な場合である。
\end{reference}
$A$ を Gorenstein Noether 局所環とする。$I \subset A$ をイデアルとし、
$Z = V(I) \subset \Spec(A)$ とおく。
$M$ を有限 $A$-加群とする。(\ref{local-cohomology-equation-cutoff}) のように
$s = s_{A, I}(M)$ とおく。このとき $H^i_Z(M)$ は $i < s$ に対して有限であるが、
$H^s_Z(M)$ は有限でない。
\end{lemma}

\begin{proof}
Gorenstein 局所環は双対化複体をもつので、これは命題
\ref{local-cohomology-proposition-finiteness} の特殊な場合である。
後で一般の有限性定理を証明する際に用いるため、この特殊な場合の短い証明が
あれば有用であろう。
\end{proof}

\noindent
次の定理の仮定は、優秀 Noether 環については定義により、
双対化複体をもつ Noether 環については
Dualizing Complexes, 補題 \ref{dualizing-lemma-universally-catenary} および
Dualizing Complexes, 補題
\ref{dualizing-lemma-dualizing-gorenstein-formal-fibres} により、また
正則 Noether 環の商についても満たされることに注意する。

\begin{theorem}
\label{local-cohomology-theorem-finiteness}
\begin{reference}
これは \cite[Satz 2]{Faltings-finiteness} の特殊な場合である。
\end{reference}
$A$ を Noether 環とし、$I \subset A$ をイデアルとする。
$Z = V(I) \subset \Spec(A)$ とおく。$M$ を有限 $A$-加群とする。
(\ref{local-cohomology-equation-cutoff}) のように $s = s_{A, I}(M)$ とおく。次を仮定する。
\begin{enumerate}
\item $A$ は普遍鎖状である。
\item $A$ の局所環の形式的ファイバーは Cohen--Macaulay である。
\end{enumerate}
このとき $H^i_Z(M)$ は $0 \leq i < s$ に対して有限であり、
$H^s_Z(M)$ は有限でない。
\end{theorem}

\begin{proof}
補題 \ref{local-cohomology-lemma-check-finiteness-local-cohomology-locally} により、
$A$ は局所環であると仮定してよい。

\medskip\noindent
$A$ が Noether 完備局所環ならば、Cohen の構造定理
（Algebra, 定理 \ref{algebra-theorem-cohen-structure-theorem}）により、
$A$ を正則完備局所環 $B$ の商として書ける。
補題 \ref{local-cohomology-lemma-cutoff} と Dualizing Complexes, 補題
\ref{dualizing-lemma-local-cohomology-and-restriction} を用いると、
正則局所環の場合に帰着する。この場合は補題
\ref{local-cohomology-lemma-local-annihilator} から従う。実際、正則局所環は Gorenstein である
（Dualizing Complexes, 補題 \ref{dualizing-lemma-regular-gorenstein}）。

\medskip\noindent
$A$ を Noether 局所環とし、$\mathfrak m$ をその極大イデアルとする。
$I \subset \mathfrak m$ と仮定してよい。そうでなければこの主張は自明である。
$A^\wedge$ を $A$ の完備化とし、$Z^\wedge = V(IA^\wedge)$ とおく。また
$M^\wedge = M \otimes_A A^\wedge$ を $M$ の完備化とする
（Algebra, 補題 \ref{algebra-lemma-completion-tensor}）。
このとき $H^i_Z(M) \otimes_A A^\wedge = H^i_{Z^\wedge}(M^\wedge)$ である。
これは Dualizing Complexes, 補題
\ref{dualizing-lemma-torsion-change-rings} と $A \to A^\wedge$ の平坦性
（Algebra, 補題 \ref{algebra-lemma-completion-flat}）による。
したがって、$H^i_{Z^\wedge}(M^\wedge)$ が $i < s$ に対して有限であり、
$i = s$ に対して有限でないことを示せば十分である。Algebra, 補題
\ref{algebra-lemma-descend-properties-modules} を参照せよ。
結論が $A^\wedge$ に対して成り立つことは既知なので、
$s_{A, I}(M) = s_{A^\wedge, I^\wedge}(M^\wedge)$ を示せば十分である。
これは補題 \ref{local-cohomology-lemma-cutoff-completion} から従う。
\end{proof}

\begin{remark}
\label{local-cohomology-remark-astute-reader}
注意深い読者は、$A$ の局所環の形式的ファイバーに課す条件を
わずかに弱められることに気づいたであろう。すなわち、定理
\ref{local-cohomology-theorem-finiteness} の状況で、$A$ は普遍鎖状であると仮定する一方、
形式的ファイバーには何も仮定しないとする。整数 $n$ が与えられ、
$H^i_Z(M)$ が $i \leq n$ に対して有限であることを示したいとする。
まったく同じ証明から、$s_{A, I}(M) > n$ であり、かつ
$A$ の局所環の形式的ファイバーが $(S_n)$ を満たせば十分であると分かる。
一方、$H^s_Z(M)$ が有限でないことを示したいとし、ここで
$s = s_{A, I}(M)$ とする。この場合、形式的ファイバーが
$(S_{s - 1})$ を満たせば、上の議論によって結論が得られる。
\end{remark}







\section{順像の有限性 II}
\label{local-cohomology-section-finiteness-pushforward-II}

\noindent
この節は第 \ref{local-cohomology-section-finiteness-pushforward} 節の続きである。
ここでは、第 \ref{local-cohomology-section-finiteness-II} 節で行った仕事の成果を得る。

\begin{lemma}
\label{local-cohomology-lemma-finiteness-Rjstar}
$X$ を局所 Noether スキームとする。$j : U \to X$ を、補集合が $Z$ である
開部分スキームの包含とする。$\mathcal{F}$ を連接
$\mathcal{O}_U$-加群とし、$n \geq 0$ を整数とする。次を仮定する。
\begin{enumerate}
\item $X$ は普遍鎖状である。
\item すべての $z \in Z$ に対して、$\mathcal{O}_{X, z}$ の形式的ファイバーは
$(S_n)$ を満たす。
\end{enumerate}
このとき、次の条件は同値である。
\begin{enumerate}
\item[(a)] $x \in \text{Supp}(\mathcal{F})$ および
$z \in Z \cap \overline{\{x\}}$ に対して
$\text{depth}_{\mathcal{O}_{X, x}}(\mathcal{F}_x) +
\dim(\mathcal{O}_{\overline{\{x\}}, z}) > n$ である。
\item[(b)] $R^pj_*\mathcal{F}$ は $0 \leq p < n$ に対して連接である。
\end{enumerate}
\end{lemma}

\begin{proof}
この主張は $X$ 上局所的なので、$X$ はアフィンであると仮定してよい。
$X = \Spec(A)$ および $Z = V(I)$ と書く。$M$ を有限 $A$-加群とし、
それに対応する連接 $\mathcal{O}_X$-加群は $\mathcal{F}$ の $U$ からの
延長であるとする。補題
\ref{local-cohomology-lemma-finiteness-pushforwards-and-H1-local} を参照せよ。
この補題から、$R^pj_*\mathcal{F}$ が連接であることと
$H^{p + 1}_Z(M)$ が有限 $A$-加群であることも同値である。
さらに、式
$\text{depth}_{\mathcal{O}_{X, x}}(\mathcal{F}_x) +
\dim(\mathcal{O}_{\overline{\{x\}}, z})$
の最小値は、(\ref{local-cohomology-equation-cutoff}) の数 $s_{A, I}(M)$ であることに注意する。
以上から、定理 \ref{local-cohomology-theorem-finiteness} により補題が従う。
この適用については注意 \ref{local-cohomology-remark-astute-reader} を参照せよ。
\end{proof}

\begin{lemma}
\label{local-cohomology-lemma-finiteness-for-finite-locally-free}
$X$ を局所 Noether スキームとする。$j : U \to X$ を、補集合が $Z$ である
開部分スキームの包含とし、$n \geq 0$ を整数とする。
$R^pj_*\mathcal{O}_U$ が $0 \leq p < n$ に対して連接ならば、
$R^pj_*\mathcal{F}$ も $0 \leq p < n$ に対して連接である。
これは任意の有限局所自由 $\mathcal{O}_U$-加群 $\mathcal{F}$ について成り立つ。
\end{lemma}

\begin{proof}
問題は $X$ 上局所的なので、$X$ はアフィンであると仮定してよい。
$X = \Spec(A)$ および $Z = V(I)$ と書く。補題
\ref{local-cohomology-lemma-finiteness-pushforwards-and-H1-local} を介して、本補題は補題
\ref{local-cohomology-lemma-local-finiteness-for-finite-locally-free} から従う。
\end{proof}

\begin{lemma}
\label{local-cohomology-lemma-annihilate-Hp}
\begin{reference}
\cite[補題 1.9]{Bhatt-local}
\end{reference}
$A$ を環とし、$J \subset I \subset A$ を有限生成イデアルとする。
$p \geq 0$ を整数とし、$U = \Spec(A) \setminus V(I)$ とおく。
$H^p(U, \mathcal{O}_U)$ が $J^n$ で零化されるような $n$ が存在するならば、
$H^p(U, \mathcal{F})$ は $J^m$ で零化される。ここで、そのような
$m = m(\mathcal{F})$ が存在し、
これは任意の有限局所自由 $\mathcal{O}_U$-加群 $\mathcal{F}$ について成り立つ。
\end{lemma}

\begin{proof}
まず、$\mathfrak a$ を $H^p(U, \mathcal{F})$ の零化イデアルとする。
$u \in U$ とする。開近傍 $u \in U' \subset U$ と同型
$\varphi : \mathcal{O}_{U'}^{\oplus r} \to \mathcal{F}|_{U'}$
が存在する。$f \in A$ を $u \in D(f) \subset U'$ となるように選ぶ。
次の射が存在する。
$$
a : \mathcal{O}_U^{\oplus r} \longrightarrow \mathcal{F}
\quad\text{および}\quad
b : \mathcal{F} \longrightarrow \mathcal{O}_U^{\oplus r}
$$
これらの $D(f)$ 上の制限は、それぞれ $f^N \varphi$ および
$f^N \varphi^{-1}$ に等しく、これはある $N$ に対して成り立つ。さらに、
$a \circ b$ と $b \circ a$ はともに $f^{2N}$ 倍写像であると仮定してよい。
これは Properties, 補題
\ref{properties-lemma-section-maps-backwards} から従う。実際、$U$ は
準コンパクトであり（$I$ は有限生成である）、分離的で、かつ
$\mathcal{F}$ と $\mathcal{O}_U^{\oplus r}$ は有限表示である。
したがって、$H^p(U, \mathcal{F})$ は $f^{2N}J^n$ で零化される。すなわち
$f^{2N}J^n \subset \mathfrak a$ である。

\medskip\noindent
$U$ は準コンパクトなので、有限個の $f_1, \ldots, f_t$ と
$N_1, \ldots, N_t$ を、$U = \bigcup D(f_i)$ かつ
$f_i^{2N_i}J^n \subset \mathfrak a$ となるように取れる。
すると $V(I) = V(f_1, \ldots, f_t)$ である。$I$ は有限生成なので、
$I^M \subset (f_1, \ldots, f_t)$ となる $M$ が存在する。
以上より、$J^m \subset \mathfrak a$ が $m \gg 0$ に対して成り立つ。
例えば $m = M (2N_1 + \ldots + 2N_t) n$ とすればよい。
\end{proof}











\section{局所コホモロジーの零化イデアル II}
\label{local-cohomology-section-annihilators-II}

\noindent
第 \ref{local-cohomology-section-annihilators} 節における局所コホモロジーの零化イデアルの議論を、
有限なコホモロジー加群をもつ下に有界な複体へ拡張する。

\begin{definition}
\label{local-cohomology-definition-depth-complex}
$I$ を Noether 環 $A$ のイデアルとし、
$K \in D^+_{\textit{Coh}}(A)$ とする。{\it $I$-深さ}を $K$ に対して定義し、
$\text{depth}_I(K)$ と書く。それは、$m \in \mathbf{Z} \cup \{\infty\}$ のうち、
$H^i_I(K) = 0$ がすべての $i < m$ に対して成り立つ最大のものとする。
$A$ が極大イデアル $\mathfrak m$ をもつ局所環ならば、
$\text{depth}_\mathfrak m(K)$ を単に $K$ の{\it 深さ}と呼ぶ。
\end{definition}

\noindent
この定義は Algebra, 定義 \ref{algebra-definition-depth} と矛盾しない。
これは Dualizing Complexes, 補題 \ref{dualizing-lemma-depth} による。

\begin{proposition}
\label{local-cohomology-proposition-annihilator-complex}
$A$ を双対化複体をもつ Noether 環とする。
$T \subset T' \subset \Spec(A)$ を特殊化で安定な部分集合とする。
$s \in \mathbf{Z}$ とし、$K$ を $D_{\textit{Coh}}^+(A)$ の対象とする。
次の条件は同値である。
\begin{enumerate}
\item イデアル $J \subset A$ が存在し、$V(J) \subset T'$ を満たして、
$J$ は $H^i_T(K)$ を $i \leq s$ に対して零化する。
\item すべての $\mathfrak p \not \in T'$ と、
$\mathfrak q \in T$ で $\mathfrak p \subset \mathfrak q$ を満たすものに対して
$$
\text{depth}_{A_\mathfrak p}(K_\mathfrak p) +
\dim((A/\mathfrak p)_\mathfrak q) > s
$$
である。
\end{enumerate}
\end{proposition}

\begin{proof}
この結果は命題 \ref{local-cohomology-proposition-annihilator} の自然な一般化であり、
読者はまず同命題の証明を読むべきである。
$\omega_A^\bullet$ を双対化複体とし、$\delta$ をその次元関数とする。
Dualizing Complexes, 第 \ref{dualizing-section-dimension-function} 節を参照せよ。
次の有限 $A$-加群が重要な役割を果たす。
$$
E^i = \Ext_A^i(K, \omega_A^\bullet)
$$
$\mathfrak p \subset A$ に対し、$H^i_\mathfrak p$ と書いて、
$D(A_\mathfrak p)$ の対象の $\mathfrak pA_\mathfrak p$ に関する
局所コホモロジーを表す。このとき、次の加群の
$\mathfrak pA_\mathfrak p$-進完備化は
$$
(E^i)_\mathfrak p =
\Ext^{\delta(\mathfrak p) + i}_{A_\mathfrak p}(K_\mathfrak p,
(\omega_A^\bullet)_\mathfrak p[-\delta(\mathfrak p)])
$$
次の加群の Matlis 双対である。
$$
H^{-\delta(\mathfrak p) - i}_{\mathfrak p}(K_\mathfrak p)
$$
これは Dualizing Complexes, 補題
\ref{dualizing-lemma-special-case-local-duality} による。
とくに、次の事実を得る。イデアル $J \subset A$ が
$(E^i)_\mathfrak p$ を零化することと、$J$ が
$H^{-\delta(\mathfrak p) - i}_{\mathfrak p}(K_\mathfrak p)$ を零化することは同値である。

\medskip\noindent
$T_n = \{\mathfrak p \in T \mid \delta(\mathfrak p) \leq n\}$ とおく。
$\delta$ は有界関数なので、$T_a = \emptyset$ は $a \ll 0$ に対して成り立ち、
$T_b = T$ は $b \gg 0$ に対して成り立つ。

\medskip\noindent
(2) を仮定し、(1) のような $J$ の存在を示そう。そのために二重帰納法を用いる。
$i \leq s$ に対して、次の帰納法の仮定 $IH_i$ を考える。
$H^a_T(K)$ は、$J \subset A$ かつ $V(J) \subset T'$ を満たすあるイデアルで
$a \leq i$ に対して零化される。$IH_i$ は、$i$ が十分小さければ自明である。
実際、$K$ は下に有界である。

\medskip\noindent
帰納段階に移る。$IH_{i - 1}$ が成り立つと仮定し、$i \leq s$ とする。
$J'$ を、$V(J') \subset T'$ を満たし、$H^a_T(K)$ を
$a \leq i - 1$ に対して零化するように選ぶ
（帰納法の仮定により、このように選べる）。$n$ に関する下降帰納法により、
イデアル $J$ が存在し、$V(J) \subset T'$ を満たし、
$J H^i_T(K)$ の随伴素イデアルが $T_n$ に属することを示す。
$n \ll 0$ に対して、これは $JH^i_T(K) = 0$ を意味する
（Algebra, 補題 \ref{algebra-lemma-ass-zero}）。したがって $IH_i$ が成り立つ。
初期段階 $n \gg 0$ は自明である。この場合は $T = T_n$ であり、
$H^i_T(K)$ のすべての随伴素イデアルが $T$ に属するからである。

\medskip\noindent
そこで、上の性質をもつ $J$ が $n$ に対して与えられたと仮定する。
$\mathfrak q \in T_n$ とする。$T_\mathfrak q \subset \Spec(A_\mathfrak q)$ を
$T$ の逆像とする。このとき
$H^j_T(K)_\mathfrak q = H^j_{T_\mathfrak q}(K_\mathfrak q)$
である。これは補題 \ref{local-cohomology-lemma-torsion-change-rings} による。
補題 \ref{local-cohomology-lemma-local-cohomology-ss} のスペクトル系列
$$
H_\mathfrak q^p(H^q_{T_\mathfrak q}(K_\mathfrak q))
\Rightarrow
H^{p + q}_\mathfrak q(K_\mathfrak q)
$$
を考える。以下で、イデアル $J'' \subset A$ で $V(J'') \subset T'$ を満たし、
$H^i_\mathfrak q(K_\mathfrak q)$ が $J''$ で零化されるものを見つける。
これはすべての $\mathfrak q \in T_n \setminus T_{n - 1}$ に対して成り立つ。
主張：$J (J')^i J''$ は $n - 1$ に対して求める性質をもつ。
実際、$\mathfrak q \in T_n \setminus T_{n - 1}$ とする。
上のスペクトル系列はフィルトレーション
$$
E_\infty^{0, i} = E_{i + 2}^{0, i} \subset \ldots \subset E_3^{0, i} \subset
E_2^{0, i} = H^0_\mathfrak q(H^i_{T_\mathfrak q}(K_\mathfrak q))
$$
を定める。加群 $E_\infty^{0, i}$ は $J''$ で零化される。
部分商 $E_j^{0, i}/E_{j + 1}^{0, i}$ は $i + 1 \geq j \geq 2$ に対して
$J'$ で零化される。実際、$d_j^{0, i}$ の終域は次の加群の部分商である。
$$
H^j_\mathfrak q(H^{i - j + 1}_{T_\mathfrak q}(K_\mathfrak q)) =
H^j_\mathfrak q(H^{i - j + 1}_T(K)_\mathfrak q)
$$
また、$H^{i - j + 1}_T(K)_\mathfrak q$ は $J'$ の選び方により $J'$ で零化される。
最後に、$J$ の選び方により
$J H^i_T(K)_\mathfrak q \subset H^0_\mathfrak q(H^i_T(K)_\mathfrak q)$
である。これは $\Spec(A_\mathfrak q)$ の非閉点における $\delta$ の値が
より大きいからである。したがって、望みどおり $\mathfrak q$ は
$J(J')^iJ'' H^i_T(K)$ の随伴素イデアルではない。

\medskip\noindent
冒頭の考察により、$J''$ は
$$
(E^{-\delta(\mathfrak q) - i})_\mathfrak q =
(E^{-n - i})_\mathfrak q
$$
を零化すべきであり、これはすべての
$\mathfrak q \in T_n \setminus T_{n - 1}$ に対して必要である。
しかし、$J''$ が一つの $\mathfrak q$ に対して有効ならば、
$\mathfrak q$ のある開近傍に属するすべての $\mathfrak q$ に対して有効である。
加群 $E^{-n - i}$ が有限だからである。
$\Spec(A)$ の任意の部分集合は誘導位相について Noether である
（Topology, 補題 \ref{topology-lemma-Noetherian}）。したがって、
$J''$ が一つの $\mathfrak q$ に対して存在することを示せば十分である。

\medskip\noindent
Ext 加群は有限なので、$J''$ の存在は
$$
\text{Supp}(E^{-n - i}) \cap \Spec(A_\mathfrak q) \subset T'.
$$
と同値である。これは、$E^{-n - i}$ の局所化が、すべての
$\mathfrak p \subset \mathfrak q$ かつ $\mathfrak p \not \in T'$ において
零であることと同値である。$A_\mathfrak p$ 上の局所双対性を用いると、
次を示せばよいと分かる。
$$
H^{i + n - \delta(\mathfrak p)}_\mathfrak p(K_\mathfrak p) =
H^{i - \dim((A/\mathfrak p)_\mathfrak q)}_\mathfrak p(K_\mathfrak p)
$$
これは零である（ここでは $\delta$ が次元関数であることを用いる）。
補題の仮定、$i \leq s$、および定義
\ref{local-cohomology-definition-depth-complex} における深さの定義により、実際に零となる。

\medskip\noindent
逆の含意を証明するため、(2) が成り立たないと仮定し、上の議論を逆にたどる。
まず、$\mathfrak q \in T$ と $\mathfrak p \subset \mathfrak q$ を
$\mathfrak p \not \in T'$ となるように選び、
$$
i = \text{depth}_{A_\mathfrak p}(K_\mathfrak p) +
\dim((A/\mathfrak p)_\mathfrak q) \leq s
$$
が最小となるようにする。このとき
$H^{i - \dim((A/\mathfrak p)_\mathfrak q)}_\mathfrak p(K_\mathfrak p)$
は、定義 \ref{local-cohomology-definition-depth-complex} における深さの定義により零でない。
$n = \delta(\mathfrak q)$ とおく。このとき、$J \subset A$ かつ
$V(J) \subset T'$ を満たし、$J(E^{-n - i})_\mathfrak q = 0$ となる
イデアルは存在しない。したがって、$H^i_\mathfrak q(K_\mathfrak q)$ は
$J \subset A$ かつ $V(J) \subset T'$ を満たすイデアルでは零化されない。
$i$ の最小性と上で表示したスペクトル系列から、加群
$H^i_T(K)_\mathfrak q$ は $J \subset A$ かつ
$V(J) \subset T'$ を満たすイデアルでは零化されない。
したがって、$H^i_T(K)$ は $J \subset A$ かつ
$V(J) \subset T'$ を満たすイデアルでは零化されない。これで命題が証明された。
\end{proof}







\section{局所コホモロジーの有限性 III}
\label{local-cohomology-section-finiteness-III}

\noindent
第 \ref{local-cohomology-section-finiteness} 節および第 \ref{local-cohomology-section-finiteness-II} 節における
局所コホモロジーの有限性の議論を、有限なコホモロジー加群をもつ
下に有界な複体へ拡張する。

\begin{lemma}
\label{local-cohomology-lemma-check-finiteness-local-cohomology-by-annihilator-complex}
$A$ を Noether 環とする。$T \subset \Spec(A)$ を特殊化で安定な部分集合とする。
$K$ を $D_{\textit{Coh}}^+(A)$ の対象とし、$n \in \mathbf{Z}$ とする。
次の条件は同値である。
\begin{enumerate}
\item $H^i_T(K)$ は $i \leq n$ に対して有限である。
\item イデアル $J \subset A$ が存在し、$V(J) \subset T$ を満たして、
$J$ は $H^i_T(K)$ を $i \leq n$ に対して零化する。
\end{enumerate}
$T = V(I) = Z$ があるイデアル $I \subset A$ に対して成り立つならば、
これらはさらに次の条件とも同値である。
\begin{enumerate}
\item[(3)] $e \geq 0$ が存在し、$I^e$ は $H^i_Z(K)$ を
$i \leq n$ に対して零化する。
\end{enumerate}
\end{lemma}

\begin{proof}
この補題は補題
\ref{local-cohomology-lemma-check-finiteness-local-cohomology-by-annihilator} の自然な一般化であり、
読者はまず同補題の証明を読むべきである。
(1) が成り立つと仮定する。$H^i_J(K) = H^i_{V(J)}(K)$ であることを思い出そう。
Dualizing Complexes, 補題 \ref{dualizing-lemma-local-cohomology-noetherian} を参照せよ。
したがって $H^i_T(K) = \colim H^i_J(K)$ であり、ここで余極限は
イデアル $J \subset A$ で $V(J) \subset T$ を満たすもの全体にわたる。
補題 \ref{local-cohomology-lemma-adjoint-ext} を参照せよ。$H^i_T(K)$ は $i \leq n$ に対して
有限生成なので、(2) のような $J \subset A$ を、写像
$H^i_J(K) \to H^i_T(K)$ が $i \leq n$ に対して全射となるように選べる。
したがって、その有限個の生成元は $J$ のある冪で零化される元であり、
$J$ をある冪で置き換えれば (2) が成り立つことが分かる。

\medskip\noindent
$a \in \mathbf{Z}$ を、$H^i(K) = 0$ が $i < a$ に対して成り立つ整数として取る。
(2) $\Rightarrow$ (1) を $a$ に関する下降帰納法により証明する。
$a > n$ ならば、$H^i_T(K) = 0$ が $i \leq n$ に対して成り立つ。
したがって (1) と (2) はともに真であり、証明すべきことはない。

\medskip\noindent
(2) のような $J$ があると仮定する。
$N = H^a_T(K) = H^0_T(H^a(K))$ は、有限 $A$-加群 $H^a(K)$ の
部分加群として有限である。$n = a$ ならば証明は終わるので、以下では
$a < n$ と仮定する。$R\Gamma_T$ の構成により、$H^i_T(N) = 0$ が
$i > 0$ に対して成り立ち、$H^0_T(N) = N$ である。
注意 \ref{local-cohomology-remark-upshot} を参照せよ。区別三角形
$$
N[-a] \to K \to K' \to N[-a + 1]
$$
を選ぶ。このとき、$H^a_T(K') = 0$ であり、$H^i_T(K) = H^i_T(K')$ が
$i > a$ に対して成り立つことが分かる。したがって、$K$ を $K'$ で
置き換えてよい。ゆえに $H^a_T(K) = 0$ と仮定してよい。
これは $H^a(K)$ の有限個の随伴素イデアルが $T$ に属さないことを意味する。
素イデアル回避（Algebra, 補題 \ref{algebra-lemma-silly}）により、$f \in J$ を、
$H^a(K)$ のどの随伴素イデアルにも含まれないように選べる。
区別三角形
$$
L \to K \xrightarrow{f} K \to L[1]
$$
を選ぶ。構成により $H^i(L) = 0$ が $i \leq a$ に対して成り立つ。
一方、コホモロジーの長完全列
$$
0 \to H^{a + 1}_T(L) \to H^{a + 1}_T(K) \xrightarrow{f}
H^{a + 1}_T(K) \to H^{a + 2}_T(L) \to H^{a + 2}_T(K) \xrightarrow{f} \ldots
$$
があり、これは同一視 $H^{a + 1}_T(L) = H^{a + 1}_T(K)$ と短完全列
$$
0 \to H^{i - 1}_T(K) \to H^i_T(L) \to H^i_T(K) \to 0
$$
へ分かれる。これは $i \leq n$ に対して成り立つ。実際、$f \in J$ である。
したがって $J^2$ は $H^i_T(L)$ を $i \leq n$ に対して零化する。
$L$ に帰納法の仮定を適用すると、$H^i_T(L)$ は $i \leq n$ に対して有限である。
短完全列をもう一度用いれば、望みどおり $H^i_T(K)$ が
$i \leq n$ に対して有限であることが分かる。

\medskip\noindent
$T = V(I)$ の場合に (2) と (3) が同値であることの証明は省略する。
\end{proof}

\begin{proposition}
\label{local-cohomology-proposition-finiteness-complex}
$A$ を双対化複体をもつ Noether 環とする。
$T \subset \Spec(A)$ を特殊化で安定な部分集合とする。
$s \in \mathbf{Z}$ とし、$K \in D_{\textit{Coh}}^+(A)$ とする。
次の条件は同値である。
\begin{enumerate}
\item $H^i_T(K)$ は有限 $A$-加群である（$i \leq s$ に対して）。
\item $\mathfrak p \not \in T$、$\mathfrak q \in T$、かつ
$\mathfrak p \subset \mathfrak q$ を満たすすべてのものに対して
$$
\text{depth}_{A_\mathfrak p}(K_\mathfrak p) +
\dim((A/\mathfrak p)_\mathfrak q) > s
$$
である。
\end{enumerate}
\end{proposition}

\begin{proof}
命題 \ref{local-cohomology-proposition-annihilator-complex} と補題
\ref{local-cohomology-lemma-check-finiteness-local-cohomology-by-annihilator-complex} の形式的帰結である。
\end{proof}






\section{連接加群の改善}
\label{local-cohomology-section-improve}

\noindent
同様の構成は \cite{EGA} に見いだされ、より最近では
\cite{Kollar-local-global-hulls} および \cite{Kollar-variants} に見いだされる。

\begin{lemma}
\label{local-cohomology-lemma-get-depth-1-along-Z}
$X$ を Noether スキームとする。$T \subset X$ を特殊化で安定な部分集合とする。
$\mathcal{F}$ を連接 $\mathcal{O}_X$-加群とする。このとき、一意な写像
$\mathcal{F} \to \mathcal{F}'$ が存在する。これは連接 $\mathcal{O}_X$-加群の
写像であり、次を満たす。
\begin{enumerate}
\item $\mathcal{F} \to \mathcal{F}'$ は全射である。
\item $\mathcal{F}_x \to \mathcal{F}'_x$ は $x \not \in T$ に対して同型である。
\item $\text{depth}_{\mathcal{O}_{X, x}}(\mathcal{F}'_x) \geq 1$ が
$x \in T$ に対して成り立つ。
\end{enumerate}
$f : Y \to X$ が平坦射で $Y$ が Noether ならば、
$f^*\mathcal{F} \to f^*\mathcal{F}'$ は $f^{-1}(T) \subset Y$ と
$f^*\mathcal{F}$ に対応する商である。
\end{lemma}

\begin{proof}
条件 (3) はちょうど $\text{Ass}(\mathcal{F}') \cap T = \emptyset$ を意味する。
したがって $\mathcal{F} \to \mathcal{F}'$ は、$\mathcal{F}$ の、台が $T$ に
含まれる切断の部分層による商である。これで一意性が証明される。
引き戻しに関する主張は Divisors, 補題 \ref{divisors-lemma-bourbaki} と
一意性から従う。

\medskip\noindent
$\mathcal{F} \to \mathcal{F}'$ の存在を示す。一意性により、$X$ 上局所的に
存在と一意性を証明すれば十分である。細部は省略する。
したがって $X = \Spec(A)$ はアフィンであり、$\mathcal{F}$ は有限
$A$-加群 $M$ に付随する連接加群であると仮定してよい。
$M' = M / H^0_T(M)$ とおく。ここで $H^0_T(M)$ は
第 \ref{local-cohomology-section-supports} 節で定めたものである。このとき
$M_\mathfrak p = M'_\mathfrak p$ が $\mathfrak p \not \in T$ に対して成り立ち、
これで (1) が証明される。一方、
$H^0_T(M) = \colim H^0_Z(M)$ であり、ここで $Z$ は $X$ の閉部分集合で
$T$ に含まれるもの全体を走る。したがって Dualizing Complexes, 補題
\ref{dualizing-lemma-divide-by-torsion} により $H^0_T(M') = 0$、すなわち
$M'$ の随伴素イデアルはどれも $T$ に属さない。ゆえに
$\text{depth}(M'_\mathfrak p) \geq 1$ が $\mathfrak p \in T$ に対して成り立つ。
\end{proof}

\begin{lemma}
\label{local-cohomology-lemma-get-depth-2-along-Z}
$j : U \to X$ を Noether スキームの開埋め込みとする。
$\mathcal{F}$ を連接 $\mathcal{O}_X$-加群とする。
$\mathcal{F}' = j_*(\mathcal{F}|_U)$ が連接であると仮定する。
このとき $\mathcal{F} \to \mathcal{F}'$ は連接 $\mathcal{O}_X$-加群の
写像のうち、次を満たす一意なものである。
\begin{enumerate}
\item $\mathcal{F}|_U \to \mathcal{F}'|_U$ は同型である。
\item $\text{depth}_{\mathcal{O}_{X, x}}(\mathcal{F}'_x) \geq 2$ が
$x \in X$、$x \not \in U$ に対して成り立つ。
\end{enumerate}
$f : Y \to X$ が平坦射で $Y$ が Noether ならば、
$f^*\mathcal{F} \to f^*\mathcal{F}'$ は $f^{-1}(U) \subset Y$ に
対応する写像である。
\end{lemma}

\begin{proof}
Divisors, 補題 \ref{divisors-lemma-depth-pushforward} の (3) により
$\text{depth}_{\mathcal{O}_{X, x}}(\mathcal{F}'_x) \geq 2$ である。
$\mathcal{F} \to \mathcal{F}'$ の一意性は Divisors, 補題
\ref{divisors-lemma-depth-2-hartog} から従う。
平坦引き戻しとの両立性は平坦基底変換から従う。Cohomology of Schemes, 補題
\ref{coherent-lemma-flat-base-change-cohomology} を参照せよ。
\end{proof}

\begin{lemma}
\label{local-cohomology-lemma-make-S2-along-Z}
$X$ を Noether スキームとし、$Z \subset X$ を閉部分スキームとする。
$\mathcal{F}$ を連接 $\mathcal{O}_X$-加群とする。$X$ は普遍鎖状であり、
局所環の形式的ファイバーは $(S_1)$ を満たすと仮定する。
このとき、一意な写像 $\mathcal{F} \to \mathcal{F}''$ が存在する。
これは連接 $\mathcal{O}_X$-加群の写像であり、次を満たす。
\begin{enumerate}
\item $\mathcal{F}_x \to \mathcal{F}''_x$ は
$x \in X \setminus Z$ に対して同型である。
\item $\mathcal{F}_x \to \mathcal{F}''_x$ は全射であり、
$\text{depth}_{\mathcal{O}_{X, x}}(\mathcal{F}''_x) = 1$ が、$x \in Z$ のうち、
直接特殊化 $x' \leadsto x$ で $x' \not \in Z$ かつ
$x' \in \text{Ass}(\mathcal{F})$ を満たすものが存在する点に対して成り立つ。
\item $\text{depth}_{\mathcal{O}_{X, x}}(\mathcal{F}''_x) \geq 2$ が
残りの $x \in Z$ に対して成り立つ。
\end{enumerate}
$f : Y \to X$ が Cohen--Macaulay 射で $Y$ が Noether ならば、
$f^*\mathcal{F} \to f^*\mathcal{F}''$ は $f^{-1}(Z) \subset Y$ に関して
同じ性質を満たす。
\end{lemma}

\begin{proof}
$\mathcal{F} \to \mathcal{F}'$ を、補題 \ref{local-cohomology-lemma-get-depth-1-along-Z} で
$Z$ という $X$ の部分集合に対して構成した写像とする。$\mathcal{F}'$ は
$\mathcal{F}$ を、台が $Z$ に含まれる切断の部分層で割った商であることを思い出そう。

\medskip\noindent
まず一意性を証明する。$\mathcal{F} \to \mathcal{F}''$ を補題のような写像とする。
分解 $\mathcal{F} \to \mathcal{F}' \to \mathcal{F}''$ を得る。実際、条件 (2) と
(3) により $\text{Ass}(\mathcal{F}'') \cap Z = \emptyset$ である。
$U \subset X$ を、$\mathcal{F}'|_U \to \mathcal{F}''|_U$ が同型となる
最大の開部分スキームとする。$U$ は (2) のような点をすべて含む。
そこで Divisors, 補題 \ref{divisors-lemma-depth-2-hartog} により
$\mathcal{F}'' = j_*(\mathcal{F}'|_U)$ と結論できる。これで一意性を得る
（細部として、このような $\mathcal{F}''$ が二つあるなら、それぞれから得られる
開集合 $U$ の共通部分を取る）。

\medskip\noindent
存在を証明する。$\text{Ass}(\mathcal{F}') = \{x_1, \ldots, x_n\}$ は有限であり、
$x_i \not \in Z$ であることを思い出そう。$Y_i$ を $\{x_i\}$ の閉包とする。
$Z_{i, j}$ を $Z \cap Y_i$ の既約成分とする。
$\text{Supp}(\mathcal{F}') \cap Z = \bigcup Z_{i, j}$ であることに注意する。
$z_{i, j} \in Z_{i, j}$ を生成点とする。次のようにおく。
$$
d_{i, j} = \dim(\mathcal{O}_{\overline{\{x_i\}}, z_{i, j}})
$$
$d_{i, j} = 1$ ならば、$z_{i, j}$ は (2) のような点の一つである。
したがって、これらの点では $\mathcal{F}'$ を変更する必要がない。
さらに、依然として $d_{i, j} = 1$ と仮定すると、補題
\ref{local-cohomology-lemma-depth-function} により、開近傍
$z_{i, j} \in V_{i, j} \subset X$ を、
$\text{depth}_{\mathcal{O}_{X, z}}(\mathcal{F}'_z) \geq 2$ が
$z \in Z_{i, j} \cap V_{i, j}$、$z \not = z_{i, j}$ に対して
成り立つように選べる。次のようにおく。
$$
Z' = X \setminus
\left(
X \setminus Z \cup \bigcup\nolimits_{d_{i, j} = 1} V_{i, j})
\right)
$$
$j' : X \setminus Z' \to X$ と書く。$Z'$ の選び方により、補題
\ref{local-cohomology-lemma-finiteness-pushforward-general} の仮定は満たされる。
$\mathcal{F}'' = j'_*(\mathcal{F}'|_{X \setminus Z'})$ とおき、補題
\ref{local-cohomology-lemma-get-depth-2-along-Z} を適用すれば結論を得る。

\medskip\noindent
最後の主張は、平坦局所準同型に沿う深さの変化の公式
（Algebra, 補題 \ref{algebra-lemma-apply-grothendieck-module} を参照せよ）と、
$f$ が Cohen--Macaulay であることに内在する $f$ のファイバーに関する仮定から
従う。細部は省略する。
\end{proof}

\begin{lemma}
\label{local-cohomology-lemma-make-S2-along-T-simple}
$X$ を局所的に双対化複体をもつ Noether スキームとする。
$T' \subset X$ を特殊化で安定な部分集合とする。
$\mathcal{F}$ を連接 $\mathcal{O}_X$-加群とする。$x \leadsto x'$ が
$X$ の点の直接特殊化であり、$x' \in T'$ かつ $x \not \in T'$ ならば、
$\text{depth}(\mathcal{F}_x) \geq 1$ であると仮定する。
このとき、一意な写像 $\mathcal{F} \to \mathcal{F}''$ が存在する。
これは連接 $\mathcal{O}_X$-加群の写像であり、次を満たす。
\begin{enumerate}
\item $\mathcal{F}_x \to \mathcal{F}''_x$ は $x \not \in T'$ に対して同型である。
\item $\text{depth}_{\mathcal{O}_{X, x}}(\mathcal{F}''_x) \geq 2$ が
$x \in T'$ に対して成り立つ。
\end{enumerate}
$f : Y \to X$ が Cohen--Macaulay 射で $Y$ が Noether ならば、
$f^*\mathcal{F} \to f^*\mathcal{F}''$ は $f^{-1}(T') \subset Y$ に関して
同じ性質を満たす。
\end{lemma}

\begin{proof}
一意性を証明する。$\mathcal{F} \to \mathcal{F}''$ を補題のような写像とする。
$U \subset X$ を、$\mathcal{F}|_U \to \mathcal{F}''|_U$ が同型となる
最大の開部分スキームとする。(1) により $U$ は集合 $X \setminus T'$ を含む。
(2) と Divisors, 補題 \ref{divisors-lemma-depth-2-hartog} により、
$\mathcal{F}'' = j_*(\mathcal{F}|_U)$ と結論できる。ここで
$j : U \to X$ は包含射である。これで一意性を得る
（細部として、このような $\mathcal{F}''$ が二つあるなら、それぞれから得られる
開集合 $U$ の共通部分を取る）。

\medskip\noindent
存在を証明する。$\mathcal{F} \to \mathcal{F}'$ を、$\mathcal{F}$ に
補題 \ref{local-cohomology-lemma-get-depth-1-along-Z} を $T'$ を用いて適用して構成した商とする。
$\mathcal{F}'$ は $\mathcal{F}$ を、台が $T'$ に含まれる切断の部分層で
割った商であることを思い出そう。次で定義する。
$$
\mathcal{F}'' = \colim j_*(\mathcal{F}'|_V)
$$
ここで $j : V \to X$ は、$X \setminus V \subset T'$ を満たす
開部分スキーム全体を走る。$T'$ は特殊化で安定なので、この余極限は
フィルター付きである。写像 $\mathcal{F}' \to j_*(\mathcal{F}'|_V)$ は
それぞれ単射である。実際、$\text{Ass}(\mathcal{F}')$ は $T'$ と交わらない。
したがって $\mathcal{F}' \to \mathcal{F}''$ は単射である。

\medskip\noindent
$X = \Spec(A)$ がアフィンであり、$\mathcal{F}$ が有限 $A$-加群 $M$ に
対応すると仮定する。このとき $\mathcal{F}'$ は
$M' = M / H^0_{T'}(M)$ に対応する。補題
\ref{local-cohomology-lemma-get-depth-1-along-Z} の証明を参照せよ。補題
\ref{local-cohomology-lemma-local-cohomology} および \ref{local-cohomology-lemma-adjoint-ext} を適用すると、
$\mathcal{F}''$ は $A$-加群 $M''$ に対応し、これは短完全列
$$
0 \to M' \to M'' \to H^1_{T'}(M') \to 0
$$
に入る。命題 \ref{local-cohomology-proposition-finiteness} と、補題の主張にある
直接特殊化に関する条件により、$M''$ は有限 $A$-加群である。
これにより $\mathcal{F}''$ は連接であることが分かる。

\medskip\noindent
$\mathcal{F}''$ の連接性と、上に表示したフィルター余極限における遷移写像の
単射性から、$\mathcal{F}'' = j_*(\mathcal{F}'|_V)$ が上のような十分小さい
$V$ のすべてに対して成り立つと結論する。このとき条件 (1) と (2) は補題
\ref{local-cohomology-lemma-get-depth-2-along-Z} から従う。

\medskip\noindent
最後の主張は、平坦局所準同型に沿う深さの変化の公式
（Algebra, 補題 \ref{algebra-lemma-apply-grothendieck-module} を参照せよ）と、
$f$ が Cohen--Macaulay であることに内在する $f$ のファイバーに関する仮定から
従う。細部は省略する。
\end{proof}

\begin{lemma}
\label{local-cohomology-lemma-make-S2-along-T}
$X$ を局所的に双対化複体をもつ Noether スキームとする。
$T' \subset T \subset X$ を特殊化で安定な部分集合とし、$x \leadsto x'$ が
$X$ の点の直接特殊化で $x' \in T'$ ならば $x \in T$ であると仮定する。
$\mathcal{F}$ を連接 $\mathcal{O}_X$-加群とする。
このとき、一意な写像 $\mathcal{F} \to \mathcal{F}''$ が存在する。
これは連接 $\mathcal{O}_X$-加群の写像であり、次を満たす。
\begin{enumerate}
\item $\mathcal{F}_x \to \mathcal{F}''_x$ は $x \not \in T$ に対して同型である。
\item $\mathcal{F}_x \to \mathcal{F}''_x$ は全射であり、
$\text{depth}_{\mathcal{O}_{X, x}}(\mathcal{F}''_x) \geq 1$ が
$x \in T$、$x \not \in T'$ に対して成り立つ。
\item $\text{depth}_{\mathcal{O}_{X, x}}(\mathcal{F}''_x) \geq 2$ が
$x \in T'$ に対して成り立つ。
\end{enumerate}
$f : Y \to X$ が Cohen--Macaulay 射で $Y$ が Noether ならば、
$f^*\mathcal{F} \to f^*\mathcal{F}''$ は
$f^{-1}(T') \subset f^{-1}(T) \subset Y$ に関して同じ性質を満たす。
\end{lemma}

\begin{proof}
まず $\mathcal{F} \to \mathcal{F}'$ を、$\mathcal{F}$ に補題
\ref{local-cohomology-lemma-get-depth-1-along-Z} を $T$ を用いて適用して構成した商とする。
次に $\mathcal{F}' \to \mathcal{F}''$ を、補題
\ref{local-cohomology-lemma-make-S2-along-T-simple} で $T'$ を用いて構成した
連接加群の一意な写像とする。このとき $\mathcal{F} \to \mathcal{F}''$ は
求める写像である。
\end{proof}










\section{Hartshorne--Lichtenbaum 消滅}
\label{local-cohomology-section-Hartshorne-Lichtenbaum-vanishing}

\noindent
この消滅結果は Lichtenbaum の定理の局所版である。同定理については
Duality for Schemes, 第 \ref{duality-section-lichtenbaum} 節を参照せよ。
この結果を含むさらに多くの事項は \cite{CD} に見いだされる。

\begin{lemma}
\label{local-cohomology-lemma-cd-top-vanishing}
$A$ を次元 $d$ の Noether 環とする。$I \subset I' \subset A$ をイデアルとする。
$I'$ が $A$ の Jacobson 根基に含まれ、$\text{cd}(A, I') < d$ ならば、
$\text{cd}(A, I) < d$ である。
\end{lemma}

\begin{proof}
補題 \ref{local-cohomology-lemma-cd-dimension} により $\text{cd}(A, I) \leq d$ である。
補題 \ref{local-cohomology-lemma-isomorphism} を用いて、写像
$$
H^d_{V(I')}(A) \to H^d_{V(I)}(A)
$$
が全射であることを示す。これで証明は終わる。
$\mathfrak p \in V(I) \setminus V(I')$ を取る。$I'$ に関する仮定により、
$\mathfrak p$ は $A$ の極大イデアルではない。したがって
$\dim(A_\mathfrak p) < d$ である。すると
$H^d_{\mathfrak pA_\mathfrak p}(A_\mathfrak p) = 0$
が補題 \ref{local-cohomology-lemma-cd-dimension} により成り立つ。
\end{proof}

\begin{lemma}
\label{local-cohomology-lemma-cd-top-vanishing-some-module}
$A$ を次元 $d$ の Noether 環とし、$I \subset A$ をイデアルとする。
$H^d_{V(I)}(M) = 0$ がある有限 $A$-加群について成り立ち、その台が次元 $d$ の
既約成分をすべて含むならば、$\text{cd}(A, I) < d$ である。
\end{lemma}

\begin{proof}
補題 \ref{local-cohomology-lemma-cd-dimension} により $\text{cd}(A, I) \leq d$ である。
したがって、任意の有限 $A$-加群 $N$ に対して $H^i_{V(I)}(N) = 0$ が
$i > d$ において成り立つ。性質 $\mathcal{P}$ が有限 $A$-加群 $N$ に対して
成り立つとは、$H^d_{V(I)}(N) = 0$ であることだとしよう。
仮定の一つは $\mathcal{P}(M)$ が成り立つことである。次に注意する。
$\mathcal{P}(N_1 \oplus N_2)
\Leftrightarrow (\mathcal{P}(N_1) \wedge \mathcal{P}(N_2))$.
また、$N \to N'$ が全射ならば、
$\mathcal{P}(N) \Rightarrow \mathcal{P}(N')$ である。実際、上で見たように
$H^{d + 1}_{V(I)}$ は消滅する。
$\mathfrak p_1, \ldots, \mathfrak p_n$ を、$A$ の極小素イデアルで
$\dim(A/\mathfrak p_i) = d$ を満たすものとする。次元に関する理由から、$\mathcal{P}(N)$ は
$N$ の台が $\{\mathfrak p_1, \ldots, \mathfrak p_n\}$ と交わらなければ成り立つ。
補題 \ref{local-cohomology-lemma-cd-dimension} を参照せよ。各 $i$ に対して
$M_i = M/\mathfrak p_i M$ とおく。これは有限 $A$-加群であり、
$\mathfrak p_i$ で零化され、その台は $V(\mathfrak p_i)$ に等しい
（ここで $M$ の台に関する仮定を用いる）。最後に、$J \subset A$ が
イデアルならば $\mathcal{P}(JM_i)$ が成り立つ。実際、$JM_i$ は $M$ のコピーの
直和の商である。したがって Cohomology of Schemes, 補題
\ref{coherent-lemma-property-higher-rank-cohomological} により、
$\mathcal{P}$ はすべての有限 $A$-加群に対して成り立つ。
\end{proof}

\begin{lemma}
\label{local-cohomology-lemma-top-coh-divisible}
$A$ を次元 $d$ の Noether 局所環とする。$f \in A$ を、次元 $d$ の
どの極小素イデアルにも含まれない元とする。このとき
$f : H^d_{V(I)}(M) \to H^d_{V(I)}(M)$ は、任意の有限 $A$-加群 $M$ と
任意のイデアル $I \subset A$ に対して全射である。
\end{lemma}

\begin{proof}
$M/fM$ の台の次元は $< d$ である。これは $f$ に関する仮定による。
したがって補題 \ref{local-cohomology-lemma-cd-dimension} により
$H^d_{V(I)}(M/fM) = 0$ である。ゆえに
$H^d_{V(I)}(fM) \to H^d_{V(I)}(M)$ は全射である。
補題 \ref{local-cohomology-lemma-cd-dimension} により $\text{cd}(A, I) \leq d$ なので、
全射 $M \to fM$、$x \mapsto fx$ が全射
$H^d_{V(I)}(M) \to H^d_{V(I)}(fM)$ を誘導することも分かる。
\end{proof}

\begin{lemma}
\label{local-cohomology-lemma-cd-bound-dualizing}
$A$ を正規化された双対化複体 $\omega_A^\bullet$ をもつ Noether 局所環とする。
$I \subset A$ をイデアルとする。
$H^0_{V(I)}(\omega_A^\bullet) = 0$ ならば、$\text{cd}(A, I) < \dim(A)$ である。
\end{lemma}

\begin{proof}
$d = \dim(A)$ とおく。$\mathfrak p_1, \ldots, \mathfrak p_n \subset A$ を
次元 $d$ の極小素イデアルとする。有限 $A$-加群
$H^{-i}(\omega_A^\bullet)$ は $i \in \{0, \ldots, d\}$ に対してのみ零でなく、
$H^{-i}(\omega_A^\bullet)$ の台の次元は $\leq i$ であることを思い出そう。
補題 \ref{local-cohomology-lemma-sitting-in-degrees} を参照せよ。
$\omega_A = H^{-d}(\omega_A^\bullet)$ とおく。素イデアル回避
（Algebra, 補題 \ref{algebra-lemma-silly}）により、$f \in A$、
$f \not \in \mathfrak p_i$ であって、$H^{-i}(\omega_A^\bullet)$ を
$i < d$ に対して零化するものを選べる。区別三角形
$$
\omega_A[d] \to \omega_A^\bullet \to
\tau_{\geq -d + 1}\omega_A^\bullet \to \omega_A[d + 1]
$$
を考える。Derived Categories, 注意
\ref{derived-remark-truncation-distinguished-triangle} を参照せよ。
Derived Categories, 補題 \ref{derived-lemma-trick-vanishing-composition} により、
$f^d$ が $\tau_{\geq -d + 1}\omega_A^\bullet$ の零自己準同型を誘導することが分かる。
三角圏の公理を用いると、写像
$$
\omega_A^\bullet \to \omega_A[d]
$$
であって、$\omega_A[d] \to \omega_A^\bullet$ との合成が
$f^d$ 倍写像となるものを得る。この倍写像は $\omega_A[d]$ 上のものである。
したがって $f^d$ は $H^d_{V(I)}(\omega_A)$ を零化する。
補題 \ref{local-cohomology-lemma-top-coh-divisible} により $H^d_{V(I)}(\omega_A) = 0$ と結論する。
あとは補題 \ref{local-cohomology-lemma-cd-top-vanishing-some-module} と、
$(\omega_A)_{\mathfrak p_i}$ が零でないという事実から結論が従う
（例えば Dualizing Complexes, 補題
\ref{dualizing-lemma-nonvanishing-generically-local} を参照せよ）。
\end{proof}

\begin{lemma}
\label{local-cohomology-lemma-inverse-system-symbolic-powers}
$(A, \mathfrak m)$ を完備 Noether 局所整域とする。
$\mathfrak p \subset A$ を次元 $1$ の素イデアルとする。
すべての $n \geq 1$ に対して $m \geq n$ が存在し、
$\mathfrak p^{(m)} \subset \mathfrak p^n$ を満たす。
\end{lemma}

\begin{proof}
記号冪 $\mathfrak p^{(m)}$ は写像
$A \to A_\mathfrak p/\mathfrak p^mA_\mathfrak p$ の核として定義されることを
思い出そう。局所化は完全なので、短完全列
$$
0 \to \mathfrak a_n \to A/\mathfrak p^n \to A/\mathfrak p^{(n)} \to 0
$$
において $\mathfrak a_n$ の台が $\{\mathfrak m\}$ に含まれると結論する。
とくに逆系 $(\mathfrak a_n)$ は Mittag--Leffler である。実際、各
$\mathfrak a_n$ は Artin $A$-加群である。
したがって補題は $\lim \mathfrak a_n = 0$ という要請と同値である。
$f \in \lim \mathfrak a_n$ とする。このとき $f$ は
$A = \lim A/\mathfrak p^n$ の元である（ここで $A$ が完備であることを用いる）。
この元は完備化 $A_\mathfrak p^\wedge$ において零に写る。これは
$A_\mathfrak p$ の完備化である。
$A_\mathfrak p \to A_\mathfrak p^\wedge$ は忠実平坦なので、$f$ は
$A_\mathfrak p$ において零に写る。$A$ は整域だから、望みどおり $f$ は零である。
\end{proof}

\begin{proposition}
\label{local-cohomology-proposition-Hartshorne-Lichtenbaum-vanishing}
\begin{reference}
\cite[定理 3.1]{CD}
\end{reference}
$A$ を、完備化が $A^\wedge$ である Noether 局所環とする。
$I \subset A$ を、次を満たすイデアルとする。
$$
\dim V(IA^\wedge + \mathfrak p) \geq 1
$$
これはすべての極小素イデアル $\mathfrak p \subset A^\wedge$ のうち、
次元 $\dim(A)$ のものに対して成り立つ。このとき
$\text{cd}(A, I) < \dim(A)$ である。
\end{proposition}

\begin{proof}
$A \to A^\wedge$ は忠実平坦なので
$H^d_{V(I)}(A) \otimes_A A^\wedge = H^d_{V(IA^\wedge)}(A^\wedge)$
が Dualizing Complexes, 補題 \ref{dualizing-lemma-torsion-change-rings} により成り立つ。
したがって $A$ は完備であると仮定してよい。

\medskip\noindent
$A$ は完備であると仮定する。$\mathfrak p_1, \ldots, \mathfrak p_n \subset A$ を
次元 $d$ の極小素イデアルとする。完備局所環
$A_i = A/\mathfrak p_i$ を考える。Dualizing Complexes, 補題
\ref{dualizing-lemma-local-cohomology-and-restriction} により
$H^d_{V(I)}(A_i) = H^d_{V(IA_i)}(A_i)$ である。補題
\ref{local-cohomology-lemma-cd-top-vanishing-some-module} により、$(A_i, IA_i)$ について
本補題を証明すれば十分である。したがって $A$ は完備局所整域であると
仮定してよい。

\medskip\noindent
$A$ は完備局所整域であると仮定する。素イデアル
$\mathfrak p \supset I$ を $\dim(A/\mathfrak p) = 1$ となるように選べる。
補題 \ref{local-cohomology-lemma-cd-top-vanishing} により、$\mathfrak p$ について
本補題を証明すれば十分である。

\medskip\noindent
補題 \ref{local-cohomology-lemma-cd-bound-dualizing} により
$H^0_{V(\mathfrak p)}(\omega_A^\bullet) = 0$ を示せば十分である。
次を思い出そう。
$$
H^0_{V(\mathfrak p)}(\omega_A^\bullet) =
\colim \text{Ext}^0_A(A/\mathfrak p^n, \omega_A^\bullet)
$$
補題 \ref{local-cohomology-lemma-inverse-system-symbolic-powers} により、この余極限は次と同じである。
$$
\colim \text{Ext}^0_A(A/\mathfrak p^{(n)}, \omega_A^\bullet)
$$
$\text{depth}(A/\mathfrak p^{(n)}) = 1$ なので、補題
\ref{local-cohomology-lemma-sitting-in-degrees} により、これらの Ext 群は望みどおり零である。
\end{proof}

\begin{lemma}
\label{local-cohomology-lemma-affine-complement}
$(A, \mathfrak m)$ を Noether 局所環とする。$I \subset A$ をイデアルとする。
$A$ は優秀かつ正規であり、$\dim V(I) \geq 1$ であると仮定する。
このとき $\text{cd}(A, I) < \dim(A)$ である。とくに $\dim(A) = 2$ ならば、
$\Spec(A) \setminus V(I)$ はアフィンである。
\end{lemma}

\begin{proof}
More on Algebra, 補題
\ref{more-algebra-lemma-completion-normal-local-ring} により、完備化
$A^\wedge$ は正規であり、したがって整域である。ゆえに命題
\ref{local-cohomology-proposition-Hartshorne-Lichtenbaum-vanishing} の仮定が成り立ち、結論を得る。
アフィン性に関する主張は補題 \ref{local-cohomology-lemma-cd-is-one} から従う。
\end{proof}
















\section{Frobenius 作用}
\label{local-cohomology-section-frobenius}

\noindent
$p$ を素数とする。$A$ を、$p = 0$ が $A$ において成り立つ環とする。
$A$ の{\it Frobenius 自己準同型}とは、写像
$$
F : A \longrightarrow A,
\quad
a \longmapsto a^p
$$
のことである。この節では、Frobenius 作用をもつ加群についての補題を証明する。

\begin{lemma}
\label{local-cohomology-lemma-annihilator-frobenius-module}
$p$ を素数とする。$(A, \mathfrak m, \kappa)$ を、$p = 0$ が $A$ において
成り立つ Noether 局所環とする。$M$ を有限 $A$-加群で
$M \otimes_{A, F} A \cong M$ を満たすものとする。このとき $M$ は有限自由である。
\end{lemma}

\begin{proof}
表示 $A^{\oplus m} \to A^{\oplus n} \to M$ を、同型
$\kappa^{\oplus n} \to M/\mathfrak m M$ を誘導するように選ぶ。
$T = (a_{ij})$ を写像 $A^{\oplus m} \to A^{\oplus n}$ の行列とする。
$a_{ij} \in \mathfrak m$ であることに注意する。$F$ による基底変換を適用し、
基底変換の右完全性を用いると、表示
$A^{\oplus m} \to A^{\oplus n} \to M$ を得る。その行列は
$T = (a_{ij}^p)$ である。したがって
$a_{ij} \in \mathfrak m^p$ となる表示がある。この構成を反復すると、各
$e \geq 1$ に対し $a_{ij} \in \mathfrak m^e$ となる表示が存在する。
これは、Fitting イデアル $\text{Fit}_k(M)$（More on Algebra, 定義
\ref{more-algebra-definition-fitting-ideal}）が $k < n$ に対して
$\bigcap_{e \geq 1} \mathfrak m^e$ に含まれることを意味する。
Krull の共通部分定理（Algebra, 補題
\ref{algebra-lemma-intersect-powers-ideal-module-zero}）によりこれは零なので、
More on Algebra, 補題
\ref{more-algebra-lemma-fitting-ideal-finite-locally-free} により
$M$ は階数 $n$ の自由加群であると結論する。
\end{proof}

\noindent
この節では、$f_1, \ldots, f_r$ を環 $A$ の元とするとき、
$\sum a_if_i = 0$ ならば $a_i \in (f_1, \ldots, f_r)$ となる場合に、
これらを{\it 独立}と呼ぶ。言い換えると、$I = (f_1, \ldots, f_r)$ とおけば、
$I/I^2$ は $A/I$ 上自由であり、基底 $f_1, \ldots, f_r$ をもつ。

\begin{lemma}
\label{local-cohomology-lemma-1}
\begin{reference}
\cite{Lech-inequalities} および \cite[補題 1, 299 頁]{MatCA} を参照せよ。
\end{reference}
$A$ を環とする。$f_1, \ldots, f_{r - 1}, f_rg_r$ が独立ならば、
$f_1, \ldots, f_r$ は独立である。
\end{lemma}

\begin{proof}
$\sum a_if_i = 0$ とする。このとき $\sum a_ig_rf_i = 0$ である。
したがって $a_r \in (f_1, \ldots, f_{r - 1}, f_rg_r)$ である。
$a_r = \sum_{i < r} b_i f_i + b f_rg_r$ と書く。このとき
$0 = \sum_{i < r} (a_i + b_if_r)f_i + bf_r^2g_r$ である。ゆえに
$a_i + b_i f_r \in (f_1, \ldots, f_{r - 1}, f_rg_r)$ であり、これは望みどおり
$a_i \in (f_1, \ldots, f_r)$ を含意する。
\end{proof}

\begin{lemma}
\label{local-cohomology-lemma-2}
\begin{reference}
\cite{Lech-inequalities} および \cite[補題 2, 300 頁]{MatCA} を参照せよ。
\end{reference}
$A$ を環とする。$f_1, \ldots, f_{r - 1}, f_rg_r$ が独立であり、かつ
$A$-加群 $A/(f_1, \ldots, f_{r - 1}, f_rg_r)$ の長さが有限ならば、
\begin{align*}
& \text{length}_A(A/(f_1, \ldots, f_{r - 1}, f_rg_r)) \\
& =
\text{length}_A(A/(f_1, \ldots, f_{r - 1}, f_r)) +
\text{length}_A(A/(f_1, \ldots, f_{r - 1}, g_r))
\end{align*}
\end{lemma}

\begin{proof}
次の完全列が存在すると主張する。
$$
0 \to
A/(f_1, \ldots, f_{r - 1}, g_r) \xrightarrow{f_r}
A/(f_1, \ldots, f_{r - 1}, f_rg_r) \to
A/(f_1, \ldots, f_{r - 1}, f_r) \to 0
$$
実際、$a f_r \in (f_1, \ldots, f_{r - 1}, f_rg_r)$ ならば、
$\sum_{i < r} a_i f_i + (a + bg_r)f_r = 0$ が、ある
$b, a_i \in A$ に対して成り立つ。したがって
$\sum_{i < r} a_i g_r f_i + (a + bg_r)g_rf_r = 0$ であり、これは
$a + bg_r \in (f_1, \ldots, f_{r - 1}, f_rg_r)$ を含意する。
ゆえに $a$ は $A/(f_1, \ldots, f_{r - 1}, g_r)$ において零に写る。
これで主張が証明された。最後に長さの加法性
（Algebra, 補題 \ref{algebra-lemma-length-additive}）を用いればよい。
\end{proof}

\begin{lemma}
\label{local-cohomology-lemma-3}
\begin{reference}
\cite{Lech-inequalities} および \cite[補題 3, 300 頁]{MatCA} を参照せよ。
\end{reference}
$(A, \mathfrak m)$ を局所環とする。$\mathfrak m = (x_1, \ldots, x_r)$ であり、
$x_1^{e_1}, \ldots, x_r^{e_r}$ がある $e_i > 0$ に対して独立ならば、
$\text{length}_A(A/(x_1^{e_1}, \ldots, x_r^{e_r})) = e_1\ldots e_r$ である。
\end{lemma}

\begin{proof}
補題 \ref{local-cohomology-lemma-1} と補題 \ref{local-cohomology-lemma-2} および帰納法を用いる。
\end{proof}

\begin{lemma}
\label{local-cohomology-lemma-flat-extension-independent}
$\varphi : A \to B$ を平坦な環準同型とする。
$f_1, \ldots, f_r \in A$ が独立ならば、
$\varphi(f_1), \ldots, \varphi(f_r) \in B$ は独立である。
\end{lemma}

\begin{proof}
$I = (f_1, \ldots, f_r)$ および $J = \varphi(I)B$ とおく。平坦性により
$I/I^2 \otimes_A B = J/J^2$ である。したがって $I/I^2$ が $A/I$ 上自由ならば、
$J/J^2$ は $B/J$ 上自由である。
\end{proof}

\begin{lemma}[Kunz]
\label{local-cohomology-lemma-frobenius-flat-regular}
\begin{reference}
\cite{Kunz-flat}
\end{reference}
$p$ を素数とする。$A$ を $p = 0$ を満たす Noether 環とする。
次の条件は同値である。
\begin{enumerate}
\item $A$ は正則である。
\item $F : A \to A$、$a \mapsto a^p$ は平坦である。
\end{enumerate}
\end{lemma}

\begin{proof}
$\Spec(F) : \Spec(A) \to \Spec(A)$ は恒等写像であることに注意する。
正則性は局所環を用いて定義され、平坦性も局所環に関する性質である。
Algebra, 補題 \ref{algebra-lemma-flat-localization} を参照せよ。
したがって $A$ は極大イデアル $\mathfrak m$ をもつ Noether 局所環であると
仮定してよく、実際そう仮定する。

\medskip\noindent
$A$ は正則であると仮定する。$x_1, \ldots, x_d$ を $A$ のパラメータ系とする。
$F$ を適用すると
$F(x_1), \ldots, F(x_d) = x_1^p, \ldots, x_d^p$ を得る。これは $A$ の
パラメータ系である。したがって $F$ は平坦である。Algebra, 補題
\ref{algebra-lemma-CM-over-regular-flat} および
\ref{algebra-lemma-regular-ring-CM} を参照せよ。

\medskip\noindent
逆に $F$ は平坦であると仮定する。$\mathfrak m = (x_1, \ldots, x_r)$ と書き、
$r$ は最小であるとする。このとき $x_1, \ldots, x_r$ は上で定義した意味で
独立である。$F$ は平坦なので、$x_1^p, \ldots, x_r^p$ は独立である。
補題 \ref{local-cohomology-lemma-flat-extension-independent} を参照せよ。したがって補題
\ref{local-cohomology-lemma-3} により
$\text{length}_A(A/(x_1^p, \ldots, x_r^p)) = p^r$ である。
$\chi(n) = \text{length}_A(A/\mathfrak m^n)$ とおく。これは次数 $\dim(A)$ の
数値多項式であることを思い出そう。Algebra, 命題
\ref{algebra-proposition-dimension} を参照せよ。$n \gg 0$ を選ぶ。次に注意する。
$$
\mathfrak m^{pn + pr} \subset F(\mathfrak m^n)A \subset \mathfrak m^{pn}
$$
これは $x_1, \ldots, x_r$ の単項式を見れば分かる。また、
$$
A/F(\mathfrak m^n)A = A/\mathfrak m^n \otimes_{A, F} A
$$
$F$ の平坦性により、これは長さ
$\chi(n) \text{length}_A(A/F(\mathfrak m)A)$ をもつ
（Algebra, 補題 \ref{algebra-lemma-pullback-module}）。上で見たことにより、
これは $p^r\chi(n)$ に等しい。したがって
$$
\chi(pn + pr) \geq p^r\chi(n) \geq \chi(pn)
$$
最高次項を見れば、これは $r = \dim(A)$、すなわち $A$ が正則であることを含意する。
\end{proof}




\section{ある種の加群の構造}
\label{local-cohomology-section-structure}

\noindent
正則局所環上のある種の加群の構造について、いくつかの結果を述べる。
この種の結果、およびさらに多くの結果は
\cite{Huneke-Sharp}, \cite{Lyubeznik}, \cite{Lyubeznik2} に見いだせる。

\begin{lemma}
\label{local-cohomology-lemma-structure-torsion-D-module-regular}
\begin{reference}
\cite[定理 2.4]{Lyubeznik} の特別な場合
\end{reference}
$k$ を標数 $0$ の体とし、$d \geq 1$ とする。
$A = k[[x_1, \ldots, x_d]]$ とし、その極大イデアルを $\mathfrak m$ とする。
$M$ を $\mathfrak m$-冪捩れ $A$-加群とし、加法的作用素
$D_1, \ldots, D_d$ が備わっているとする。これらは Leibniz 則
$$
D_i(fz) = \partial_i(f) z + f D_i(z)
$$
を $f \in A$ および $z \in M$ に対して満たす。ここで $\partial_i$ は
$x_i$ に関する微分である。
このとき $M$ は、$E$（$k$ の入射包絡）のコピーの直和と同型である。
\end{lemma}

\begin{proof}
集合 $J$ と同型 $M[\mathfrak m] \to \bigoplus_{j \in J} k$ を選ぶ。
$\bigoplus_{j \in J} E$ は入射的なので
（Dualizing Complexes, 補題 \ref{dualizing-lemma-sum-injective-modules}）、
この同型を $A$-加群準同型
$\varphi : M \to \bigoplus_{j \in J} E$ に延長できる。
$\varphi$ は同型、すなわち全単射であると主張する。

\medskip\noindent
単射性。$z \in M$ を零でない元とする。$M$ は $\mathfrak m$-冪捩れなので、
元 $f \in A$ であって $fz \in M[\mathfrak m]$ かつ
$fz \not = 0$ となるものを選べる。
このとき $\varphi(fz) = f\varphi(z)$ は零でないから、
$\varphi(z)$ も零でない。

\medskip\noindent
全射性。$z \in M$ とする。このとき $x_1^n z = 0$ となる
$n \geq 0$ がある。$z \in x_1M$ を $n$ に関する帰納法で証明する。
$n = 0$ ならば $z = 0$ なので主張は成り立つ。
$n > 0$ ならば、$D_1$ を作用させて
$0 = n x_1^{n - 1} z + x_1^nD_1(z)$ を得る。
したがって $x_1^{n - 1}(nz + x_1D_1(z)) = 0$ である。帰納法により
$nz + x_1D_1(z) \in x_1M$ となる。$n$ は可逆なので、
$z \in x_1M$ と結論する。よって $M$ は $x_1$-可除である。
$\varphi$ が全射でないならば、元 $e \in \bigoplus_{j \in J} E$ であって
$M$ に属さないものを選べる。
上と同様に議論すれば $\mathfrak m e \subset M$ と仮定してよく、
特に $x_1 e \in M$ である。元 $z_1 \in M$ で
$x_1 z_1 = x_1 e$ を満たすものが存在する。したがって
$x_1(z_1 - e) = 0$ である。$e$ を $e - z_1$ で置き換えることにより、
$e$ は $x_1$ で零化されると仮定してよい。
したがって、
$$
\varphi[x_1] :
M[x_1]
\longrightarrow
\left(\bigoplus\nolimits_{j \in J} E\right)[x_1] =
\bigoplus\nolimits_{j \in J} E[x_1]
$$
が全射であることを証明すれば十分である。$d = 1$ ならば、これは
$\varphi$ の構成から成り立つ。$d > 1$ ならば、$E[x_1]$ は
$k[[x_2, \ldots, x_d]]$ の剰余体の入射包絡であることに注意する。
Dualizing Complexes, 補題 \ref{dualizing-lemma-quotient} を参照せよ。
$M[x_1]$ は $k[[x_2, \ldots, x_d]]$ 上の加群として
$\mathfrak m/(x_1)$-冪捩れであり、上に表示した Leibniz 則を満たす
作用素 $D_2, \ldots, D_d$ が備わっていることに注意する。
したがって $d$ に関する帰納法により、望みどおり $\varphi[x_1]$ は
全射であると結論する。
\end{proof}

\begin{lemma}
\label{local-cohomology-lemma-structure-torsion-Frobenius-regular}
\begin{reference}
少し議論を補えば \cite[系 3.6]{Huneke-Sharp} から従う。
また \cite[定理 1.4]{Lyubeznik2} から直接にも従う。
\end{reference}
$p$ を素数とする。$(A, \mathfrak m, k)$ を $p = 0$ を満たす
正則局所環とする。$F : A \to A$, $a \mapsto a^p$ を Frobenius
自己準同型とする。$M$ を $\mathfrak m$-冪捩れ加群であって
$M \otimes_{A, F} A \cong M$ を満たすものとする。このとき $M$ は、
$E$（$k$ の入射包絡）のコピーの直和と同型である。
\end{lemma}

\begin{proof}
集合 $J$ と $A$-加群準同型
$\varphi : M \to \bigoplus_{j \in J} E$ であって、
$M[\mathfrak m]$ を
$(\bigoplus_{j \in J} E)[\mathfrak m] = \bigoplus_{j \in J} k$ へ
同型に写すものを選ぶ。
$\varphi$ は同型、すなわち全単射であると主張する。

\medskip\noindent
単射性。$z \in M$ を零でない元とする。$M$ は $\mathfrak m$-冪捩れなので、
元 $f \in A$ であって $fz \in M[\mathfrak m]$ かつ
$fz \not = 0$ となるものを選べる。
このとき $\varphi(fz) = f\varphi(z)$ は零でないから、
$\varphi(z)$ も零でない。

\medskip\noindent
全射性。補題 \ref{local-cohomology-lemma-frobenius-flat-regular} により
$F$ は平坦であることを思い出そう。
$x_1, \ldots, x_d$ を $\mathfrak m$ の極小生成系とする。
$$
M_n = M[x_1^{p^n}, \ldots, x_d^{p^n}]
$$
を、$M$ の部分加群であって $x_1^{p^n}, \ldots, x_d^{p^n}$ で
零化される元からなるものとする。したがって $M_0 = M[\mathfrak m]$ は
$k$ 上のベクトル空間である。また、$M = \bigcup M_n$ である。
これは $M$ が $\mathfrak m$-冪捩れであるという仮定による。
$F^n$ は平坦であり、
$F^n(x_i) = x_i^{p^n}$ なので、
$$
M_n \cong (M \otimes_{A, F^n} A)[x_1^{p^n}, \ldots, x_d^{p^n}] =
M[x_1, \ldots, x_d] \otimes_{A, F^n} A =
M_0 \otimes_k A/(x_1^{p^n}, \ldots, x_d^{p^n})
$$
を得る。したがって $M_n$ は $A/(x_1^{p^n}, \ldots, x_d^{p^n})$ 上自由である。
計算により、$A/(x_1^{p^n}, \ldots, x_d^{p^n})$ の任意の元で
$x_1^{p^n - 1}$ により零化されるものは $x_1$ で割り切れることが分かる。
例えば、Algebra, 補題
\ref{algebra-lemma-regular-complete-containing-coefficient-field} による同型
$A/(x_1^{p^n}, \ldots, x_d^{p^n}) \cong
k[x_1, \ldots, x_d]/(x_1^{p^n}, \ldots, x_d^{p^n})$
を用いればよい。
よって同じことが $M_n$ の任意の元について成り立つ。
$M$ の各元は $M_n$ にすべての $n \gg 0$ に対して属し、また
$M$ の各元は $x_1$ のある冪で零化されるので、$M$ は $x_1$-可除であると結論する。

\medskip\noindent
$x = x_1$ とおく。上で $M$ は $x$-可除であることを見た。
$\varphi$ が全射でないならば、元 $e \in \bigoplus_{j \in J} E$ であって
$M$ に属さないものを選べる。
上と同様に議論すれば $\mathfrak m e \subset M$ と仮定してよく、
特に $x e \in M$ である。元 $z_1 \in M$ で
$x z_1 = x e$ を満たすものが存在する。したがって
$x(z_1 - e) = 0$ である。$e$ を $e - z_1$ で置き換えることにより、
$e$ は $x$ で零化されると仮定してよい。
したがって、
$$
\varphi[x] :
M[x]
\longrightarrow
\left(\bigoplus\nolimits_{j \in J} E\right)[x] =
\bigoplus\nolimits_{j \in J} E[x]
$$
が全射であることを証明すれば十分である。$d = 1$ ならば、これは
$\varphi$ の構成から成り立つ。$d > 1$ ならば、$E[x]$ は正則環
$A/xA$ の剰余体の入射包絡であることに注意する。
Dualizing Complexes, 補題 \ref{dualizing-lemma-quotient} を参照せよ。
$M[x]$ は $A/xA$ 上の加群として $\mathfrak m/(x)$-冪捩れであり、
次が成り立つことに注意する。
\begin{align*}
M[x] \otimes_{A/xA, F} A/xA
& = M[x] \otimes_{A, F} A \otimes_A A/xA \\
& = (M \otimes_{A, F} A)[x^p] \otimes_A A/xA \\
& \cong M[x^p] \otimes_A A/xA
\end{align*}
前と同様に $F$ の平坦性を用いて議論する。
$M[x^p] \otimes_A A/xA \to M[x]$、
$z \otimes 1 \mapsto x^{p - 1}z$ は同型であると主張する。
これは、上で定義した各加群 $M_n$（$n > 0$）について証明すれば分かる。
その場合には $A/(x_1^{p^n}, \ldots, x_d^{p^n})$ と $x = x_1$ に対する
同じ結果から従う。
したがって $\dim(A)$ に関する帰納法により、望みどおり $\varphi[x]$ は
全射であると結論する。
\end{proof}





\section{局所コホモロジー上の追加構造}
\label{local-cohomology-section-additional}

\noindent
一つの結果例を挙げる。

\begin{lemma}
\label{local-cohomology-lemma-derivation}
$A$ を環とし、$I \subset A$ を有限生成イデアルとする。
$Z = V(I)$ とおく。
各導分 $\theta : A \to A$ に対し、標準的な加法的作用素 $D$ が
局所コホモロジー加群 $H^i_Z(A)$ 上に存在し、$\theta$ に関する
Leibniz 則を満たす。
\end{lemma}

\begin{proof}
$f_1, \ldots, f_r$ を $I$ の生成元とする。
$R\Gamma_Z(A)$ は複体
$$
A \to \prod\nolimits_{i_0} A_{f_{i_0}} \to
\prod\nolimits_{i_0 < i_1} A_{f_{i_0}f_{i_1}}
\to \ldots \to A_{f_1\ldots f_r}
$$
で計算されることを思い出そう。Dualizing Complexes, 補題
\ref{dualizing-lemma-local-cohomology-adjoint} を参照せよ。
$\theta$ は $A$ の任意の局所化上の、$\theta$ に関する Leibniz 則を満たす
加法的作用素へ一意的に延長されるので、補題は明らかである。
\end{proof}

\begin{lemma}
\label{local-cohomology-lemma-frobenius}
$p$ を素数とする。$A$ を $p = 0$ を満たす環とする。
$F : A \to A$, $a \mapsto a^p$ を Frobenius 自己準同型とする。
$I \subset A$ を有限生成イデアルとし、$Z = V(I)$ とおく。
同型
$R\Gamma_Z(A) \otimes_{A, F}^\mathbf{L} A \cong R\Gamma_Z(A)$
が存在する。
\end{lemma}

\begin{proof}
Dualizing Complexes, 補題
\ref{dualizing-lemma-torsion-change-rings}
および、$Z = V(f_1^p, \ldots, f_r^p)$ であるという事実から従う
（ここで $I = (f_1, \ldots, f_r)$ とした）。
\end{proof}

\begin{lemma}
\label{local-cohomology-lemma-etale-derivation}
$A$ を環とする。$V \to \Spec(A)$ を準コンパクト、準分離的かつ
\'etale とする。各導分 $\theta : A \to A$ に対し、標準的な加法的作用素
$D$ が $H^i(V, \mathcal{O}_V)$ 上に存在し、$\theta$ に関する
Leibniz 則を満たす。
\end{lemma}

\begin{proof}
$V$ が分離的ならば、アフィン開被覆
$V = \bigcup_{j = 1, \ldots m} V_j$ を用いて議論できる。実際、$V$ は
分離的なので $V_{j_0 \ldots j_p} = \Spec(B_{j_0 \ldots j_p})$ と書ける。
Schemes, 補題 \ref{schemes-lemma-characterize-separated} を参照せよ。
すると、$A$-加群 $H^i(V, \mathcal{O}_V)$ は第 $i$ コホモロジー群であり、
対応する {\v C}ech 複体は
$$
\prod B_{j_0} \to
\prod B_{j_0j_1} \to
\prod B_{j_0j_1j_2} \to \ldots
$$
であることが分かる。Cohomology of Schemes, 補題
\ref{coherent-lemma-cech-cohomology-quasi-coherent}
を参照せよ。各 $B = B_{j_0 \ldots j_p}$ は \'etale $A$-代数である。
したがって $\Omega_B = \Omega_A \otimes_A B$ であり、$\theta$ は導分
$\theta_B : B \to B$ に一意的に延長されると結論する。
これらの写像は {\v C}ech 複体の自己準同型を定め、コホモロジー群上に
所望の作用素を定める。

\medskip\noindent
一般の場合には、Cohomology of Schemes, 補題
\ref{coherent-lemma-hypercoverings} の証明の前半とまったく同様に、
アフィン開部分による $V$ の超被覆を用いる。詳細は省略する。
\end{proof}

\begin{remark}
\label{local-cohomology-remark-higher-order-operators}
補題 \ref{local-cohomology-lemma-derivation} および \ref{local-cohomology-lemma-etale-derivation} は、
Algebra, 補題 \ref{algebra-lemma-formally-etale-principal-parts} および
注意 \ref{algebra-remark-formally-etale-differential-operators} を用いれば、
高階微分作用素も含むように強められる。これが必要になったならば、
ここで正確な補題を述べて証明することにする。
\end{remark}

\begin{lemma}
\label{local-cohomology-lemma-etale-frobenius}
$p$ を素数とする。$A$ を $p = 0$ を満たす環とする。
$F : A \to A$, $a \mapsto a^p$ を Frobenius 自己準同型とする。
$V \to \Spec(A)$ が準コンパクト、準分離的かつ \'etale ならば、同型
$R\Gamma(V, \mathcal{O}_V) \otimes_{A, F}^\mathbf{L} A \cong
R\Gamma(V, \mathcal{O}_V)$
が存在する。
\end{lemma}

\begin{proof}
相対 Frobenius 射
$$
V \longrightarrow V \times_{\Spec(A), \Spec(F)} \Spec(A)
$$
は $V$ の $A$ 上のものであり、同型であることに注意する。
\'Etale Morphisms,
補題 \ref{etale-lemma-relative-frobenius-etale} を参照せよ。
したがって補題はコホモロジーと基底変換から従う。Derived Categories of Schemes,
補題 \ref{perfect-lemma-compare-base-change} を参照せよ。
$V$ は $A$ 上 \'etale なので、$A$ 上平坦であることに注意する。
\end{proof}





\section{一様性について I}
\label{local-cohomology-section-uniform-vanishing}

\noindent
この節の主な目的は、補題
\ref{local-cohomology-lemma-characterize-vanishing-tor-ext-above-e} を定式化し証明することである。

\begin{lemma}
\label{local-cohomology-lemma-map-tor-1-zero}
$R$ を環とする。$M \to M'$ を $R$-加群の写像とし、
$M$ は有限表示であり、$\text{Tor}_1^R(M, N) \to \text{Tor}_1^R(M', N)$ は
任意の $R$-加群 $N$ に対して零であるとする。
このとき $M \to M'$ は自由 $R$-加群を経由する。
\end{lemma}

\begin{proof}
短完全列の間の写像
$$
\xymatrix{
0 \ar[r] &
K \ar[r] \ar[d] &
R^{\oplus r} \ar[r] \ar[d] &
M \ar[r] \ar[d] &
0 \\
0 \ar[r] &
K' \ar[r] &
\bigoplus_{i \in I} R \ar[r] &
M' \ar[r] &
0
}
$$
であって、その右の縦射が与えられた写像となるものを選べる。
この写像を短完全列
\begin{equation}
\label{local-cohomology-equation-pushout}
0 \to K' \to E \to M \to 0
\end{equation}
を経由して分解できる。この列は、最初の短完全列を $K \to K'$ に沿って
押し出したものである。図式追跡により、補題の仮定は境界写像
$\text{Tor}_1^R(M, N) \to K' \otimes_R N$ であって
（\ref{local-cohomology-equation-pushout}）から誘導されるものが零であること、すなわち列
（\ref{local-cohomology-equation-pushout}）が普遍完全であることを含意すると分かる。
Algebra, 補題 \ref{algebra-lemma-universally-exact-split} により、これは
（\ref{local-cohomology-equation-pushout}）が分裂することを含意する（ここで
$M$ が有限表示であることを用いる）。したがって写像 $M \to M'$ は
$\bigoplus_{i \in I} R$ を経由し、主張が従う。
\end{proof}

\begin{lemma}
\label{local-cohomology-lemma-characterize-vanishing-tor-ext-above-e}
$R$ を環とする。$\alpha : M \to M'$ を $R$-加群の写像とする。
$P_\bullet \to M$ および $P'_\bullet \to M'$ を射影
$R$-加群による分解とする。$e \geq 0$ を整数とする。
次の条件を考える。
\begin{enumerate}
\item 複体の写像 $a_\bullet : P_\bullet \to P'_\bullet$ であって、
コホモロジー上で $\alpha$ を誘導し、$a_i = 0$ が $i > e$ に対して成り立つものが存在する。
\item 複体の写像 $a_\bullet : P_\bullet \to P'_\bullet$ であって、
コホモロジー上で $\alpha$ を誘導し、$a_{e + 1} = 0$ となるものが存在する。
\item 写像 $\Ext^i_R(M', N) \to \Ext^i_R(M, N)$ は、任意の
$R$-加群 $N$ および $i > e$ に対して零である。
\item 写像 $\Ext^{e + 1}_R(M', N) \to \Ext^{e + 1}_R(M, N)$ は、任意の
$R$-加群 $N$ に対して零である。
\item $N = \Im(P'_{e + 1} \to P'_e)$ とおき、
$\xi \in \Ext^{e + 1}_R(M', N)$ を標準元とする（証明を参照）。
このとき $\xi$ は $\Ext^{e + 1}_R(M, N)$ において零に写る。
\item 写像 $\text{Tor}_i^R(M, N) \to \text{Tor}_i^R(M', N)$ は、任意の
$R$-加群 $N$ および $i > e$ に対して零である。
\item 写像 $\text{Tor}_{e + 1}^R(M, N) \to \text{Tor}_{e + 1}^R(M', N)$ は、
任意の $R$-加群 $N$ に対して零である。
\end{enumerate}
このとき常に次の含意が成り立つ。
$$
(1) \Leftrightarrow (2) \Leftrightarrow (3) \Leftrightarrow (4)
\Leftrightarrow (5) \Rightarrow (6) \Leftrightarrow (7)
$$
$M$ が $(-e - 1)$-擬連接ならば（例えば $R$ が Noether 環で、
$M$ が有限 $R$-加群ならば）、すべての条件は同値である。
\end{lemma}

\begin{proof}
(2) が (1) を含意することは明らかである。$a_\bullet$ が (1) のとおりならば、
複体の写像 $a'_\bullet : P_\bullet \to P'_\bullet$ であって、
$a'_i = a_i$ が $i \leq e + 1$ に対して成り立ち、$a'_i = 0$ が
$i \geq e + 1$ に対して成り立つものを考えれば、(2) のとおりの複体の写像を得る。
したがって (1) と (2) は同値である。

\medskip\noindent
分解を用いた $\Ext$ 関手および $\text{Tor}$ 関手の構成
（Algebra, 節 \ref{algebra-section-ext} および
\ref{algebra-section-tor}）から、(1) と (2) は他のすべての条件を
含意することが分かる。

\medskip\noindent
(3) が (4) を、(4) が (5) を含意することは明らかである。
$N$ を (5) のとおりとする。標準写像 $\tilde \xi : P'_{e + 1} \to N$ に
$P'_{e + 2} \to P'_{e + 1}$ を前合成したものは零である。したがって
$\xi$ を $\tilde \xi$ の次の群における類として考えられる。
$$
\Ext^{e + 1}_R(M', N) =
\frac{\Ker(\Hom(P'_{e + 1}, N \to \Hom(P'_{e + 2}, N)}{
\Im(\Hom(P'_e, N \to \Hom(P'_{e + 1}, N)}
$$
複体の写像 $a_\bullet : P_\bullet \to P'_\bullet$ であって
$\alpha$ を持ち上げるものを選ぶ。
Derived Categories, 補題
\ref{derived-lemma-morphisms-lift-projective} を参照せよ。
$\xi$ が $\Ext^{e + 1}_R(M', N)$ において零に写るならば、写像
$\varphi : P_e \to N$ であって
$\tilde \xi  \circ a_{e + 1} = \varphi \circ d$ を満たすものが存在する。
したがって、複体の写像
$$
\xymatrix{
\ldots \ar[r] &
P_{e + 1} \ar[r] \ar[d]^0 &
P_e \ar[r] \ar[d]^{a_e - \varphi} &
P_{e - 1} \ar[r] \ar[d]^{a_{e - 1}} &
\ldots \\
\ldots \ar[r] &
P'_{e + 1} \ar[r] &
P'_e \ar[r] &
P'_{e - 1} \ar[r] &
\ldots
}
$$
であって (2) のとおりのものを得る。ゆえに (1) -- (5) は同値である。

\medskip\noindent
(6) と (7) の同値性は次元移動から従う。詳細は省略する。

\medskip\noindent
$M$ は $(-e - 1)$-擬連接であると仮定する（補題中の括弧内の主張は
More on Algebra, 補題
\ref{more-algebra-lemma-Noetherian-pseudo-coherent} から従う）。
(7) が (4) を含意することを示せば証明が完了する。$e$ に関する帰納法を用いる。
基底の場合は $e = 0$ である。このとき $M$ は有限表示である。
More on Algebra, 補題 \ref{more-algebra-lemma-n-pseudo-module} を参照せよ。
さらに補題 \ref{local-cohomology-lemma-map-tor-1-zero} から、$M \to M'$ は自由加群を
経由すると結論する。もちろん $M \to M'$ が自由加群を経由するならば、
$\Ext^i_R(M', N) \to \Ext^i_R(M, N)$ はすべての $i > 0$ に対して
望みどおり零である。
$e > 0$ と仮定する。短完全列の間の写像
$$
\xymatrix{
0 \ar[r] &
K \ar[r] \ar[d] &
R^{\oplus r} \ar[r] \ar[d] &
M \ar[r] \ar[d] &
0 \\
0 \ar[r] &
K' \ar[r] &
\bigoplus_{i \in I} R \ar[r] &
M' \ar[r] &
0
}
$$
であって、その右の縦射が与えられた写像となるものを選べる。すると、
$\text{Tor}_{i + 1}^R(M, N) = \text{Tor}^R_i(K, N)$
および $\Ext^{i + 1}_R(M, N) = \Ext^i_R(K, N)$ を $i \geq 1$ と任意の
$R$-加群 $N$ に対して得る。$M', K'$ についても同様である。
したがって $\text{Tor}_e^R(K, N) \to \text{Tor}_e^R(K', N)$ は任意の
$R$-加群 $N$ に対して零である。More on Algebra, 補題
\ref{more-algebra-lemma-cone-pseudo-coherent} により、$K$ は
$(-e)$-擬連接であることが分かる。帰納法により
$\Ext^e(K', N) \to \Ext^e(K, N)$ は任意の $R$-加群
$N$ に対して零であると結論し、所望の結果を得る。
\end{proof}

\begin{lemma}
\label{local-cohomology-lemma-cd-sequence-Koszul}
$I$ を Noether 環 $A$ のイデアルとする。
任意の $n \geq 1$ に対し、ある $m > n$ が存在して、写像
$A/I^m \to A/I^n$ は補題
\ref{local-cohomology-lemma-characterize-vanishing-tor-ext-above-e} の同値な条件を
$e = \text{cd}(A, I)$ として満たす。
\end{lemma}

\begin{proof}
$\xi \in \Ext^{e + 1}_A(A/I^n, N)$ を、補題
\ref{local-cohomology-lemma-characterize-vanishing-tor-ext-above-e} の (5) で構成された元とする。
$e = \text{cd}(A, I)$ なので、
$0 = H^{e + 1}_Z(N) = H^{e + 1}_I(N) = \colim \Ext^{e + 1}(A/I^m, N)$
が成り立つ。これは Dualizing Complexes, 補題
\ref{dualizing-lemma-local-cohomology-noetherian} および
\ref{dualizing-lemma-local-cohomology-ext} による。
したがって、$m \geq n$ であって、$\xi$ が
$\Ext^{e + 1}_A(A/I^m, N)$ において零に写るようなものを選べ、所望の結果を得る。
\end{proof}






\section{一様性について II}
\label{local-cohomology-section-uniform-vanishing-bis}

\noindent
$I$ を Noether 環 $A$ のイデアルとする。$M$ を有限
$A$-加群とし、$i > 0$ とする。More on Algebra, 補題
\ref{more-algebra-lemma-tor-strictly-pro-zero}
により、ある $c = c(A, I, M, i)$ が存在して、
$\text{Tor}^A_i(M, A/I^n) \to \text{Tor}^A_i(M, A/I^{n - c})$
はすべての $n \geq c$ に対して零となる。この節では、場合によっては
定数 $c$ であって、すべての $A$-加群 $M$ に同時に通用するものを選べること
（しかも $i$ のある範囲にわたって選べること）を示す結果を論じる。
この内容は、\cite{Huneke-uniform} および \cite{AHS} で論じられる
一様 Artin--Rees と関係している。

\medskip\noindent
注意 \ref{local-cohomology-remark-strict-pro-isomorphism} では、これを応用して、導来完備化に
関係する種々の pro 系が強い意味で pro 同型であること（またはそうでないこと）を示す。

\medskip\noindent
次の補題は大幅に強めることができる。

\begin{lemma}
\label{local-cohomology-lemma-maps-zero-fixed-torsion}
$I$ を Noether 環 $A$ のイデアルとする。任意の $m \geq 0$ および
$i > 0$ に対し、ある $c = c(A, I, m, i) \geq 0$ が存在して、
任意の $A$-加群 $M$ であって $I^m$ で零化されるものに対し、写像
$$
\text{Tor}^A_i(M, A/I^n) \to \text{Tor}^A_i(M, A/I^{n - c})
$$
はすべての $n \geq c$ に対して零である。
\end{lemma}

\begin{proof}
$i$ に関する帰納法による。基底の場合は $i = 1$ である。短完全列
$0 \to I^n \to A \to A/I^n \to 0$ は単射
$\text{Tor}_1^A(M, A/I^n) \subset I^n \otimes_A M$ を定める。
Algebra, 注意 \ref{algebra-remark-Tor-ring-mod-ideal} を参照せよ。
$M$ は $I^m$ で零化されるので、写像
$I^n \otimes_A M \to I^{n - m} \otimes_A M$ は
$n \geq m$ に対して零であると分かる。したがって $c = m$ とすればよい。

\medskip\noindent
帰納段階。$i > 1$ とし、$c$ は $i - 1$ に対して通用すると仮定する。
More on Algebra, 補題 \ref{more-algebra-lemma-tor-strictly-pro-zero} を
$M = A/I^m$ に適用すると、ある $c' \geq 0$ を選んで、
$\text{Tor}_i(A/I^m, A/I^n) \to \text{Tor}_i(A/I^m, A/I^{n - c'})$
が $n \geq c'$ に対して零となるようにできる。$M$ は $I^m$ で零化されるとする。
短完全列
$$
0 \to S \to \bigoplus\nolimits_{i \in I} A/I^m \to M \to 0
$$
を選ぶ。対応する Tor の長完全列から完全列
$$
\text{Tor}_i^A(\bigoplus\nolimits_{i \in I} A/I^m, A/I^n) \to
\text{Tor}_i^A(M, A/I^n) \to
\text{Tor}_{i - 1}^A(S, A/I^n)
$$
をすべての整数 $n \geq 0$ に対して得る。$n \geq c + c'$ ならば、写像
$\text{Tor}_{i - 1}^A(S, A/I^n) \to \text{Tor}_{i - 1}^A(S, A/I^{n - c})$
は零であり、写像 $\text{Tor}_i^A(A/I^m, A/I^{n - c}) \to 
\text{Tor}_i^A(A/I^m, A/I^{n - c - c'})$ も零である。これを短完全列と
組み合わせると、$i$ に対する結果が定数 $c + c'$ で成り立つ。
\end{proof}

\begin{lemma}
\label{local-cohomology-lemma-annihilates-affine}
$I = (a_1, \ldots, a_t)$ を Noether 環 $A$ のイデアルとする。
$a = a_1$ とおき、$B = A[\frac{I}{a}]$ をアフィン・ブローアップ代数とする。
ある $c > 0$ が存在して、$\text{Tor}_i^A(B, M)$ は $I^c$ で零化される。
これは任意の $A$-加群 $M$ および $i \geq t$ に対して成り立つ。
\end{lemma}

\begin{proof}
$B$ は
$A[x_2, \ldots, x_t]/(a_1x_2 - a_2, \ldots, a_1x_t - a_t)$
をその $a_1$-捩れで割った商であることを思い出そう。Algebra, 補題
\ref{algebra-lemma-affine-blowup-quotient-description} を参照せよ。
$$
B_\bullet = \text{次の列に関する Koszul 複体：}a_1x_2 - a_2, \ldots, a_1x_t - a_t
\text{、係数環は }A[x_2, \ldots, x_t]
$$
を、次数 $(t - 1), \ldots, 0$ に置かれた鎖複体とみなす。
複体 $B_\bullet[1/a_1]$ は、列
$x_2 - a_2/a_1, \ldots, x_t - a_t/a_1$ に関する Koszul 複体と同型であり、
この列は $A[1/a_1][x_2, \ldots, x_t]$ における正則列である。
正則列は Koszul 正則なので、増大写像
$$
\epsilon : B_\bullet \longrightarrow B
$$
は $a_1$ を可逆化すると擬同型になる。錐 $C_\bullet$（$\epsilon$ に関するもの）の
ホモロジー加群は有限 $A[x_2, \ldots, x_n]$-加群であり、$C_\bullet$ は
有界なので、ある $c \geq 0$ が存在して、$a_1^c$ がそれらすべてを
零化すると結論する。Derived Categories, 補題
\ref{derived-lemma-trick-vanishing-composition} により、必要ならば $c$ を
より大きな整数で置き換えることで、$a_1^c$ は $C_\bullet$ 上、$D(A)$ において
零であると分かる。区別三角形
$$
B_\bullet \otimes_A^\mathbf{L} M \to
B \otimes_A^\mathbf{L} M \to
C_\bullet \otimes_A^\mathbf{L} M
$$
を考えれば証明は完了する。実際、第 1 項は $B_\bullet \otimes_A M$ で表され、
ホモロジー次数 $(t - 1), \ldots, 0$ に置かれている。これは Koszul 複体
$B_\bullet$ の各項が自由な（したがって平坦な）$A$-加群だからである。よって
$\text{Tor}_i^A(B, M) = H_i(C_\bullet \otimes_A^\mathbf{L} M)$
が $i > t - 1$ に対して成り立ち、これは $a_1^c$ で零化される。
$a_1^cB = I^cB$ であり、この Tor 加群は $B$ 上の加群なので、結論を得る。
\end{proof}

\noindent
この節の残りでは、Noether 環 $A$ とイデアル $I \subset A$ を固定する。
$$
p : X \to \Spec(A)
$$
を、$\Spec(A)$ のイデアル $I$ に沿うブローアップとする。言い換えると、$X$ は
$\text{Proj}$、すなわち Rees 代数 $\bigoplus_{n \geq 0} I^n$ の射影スキームである。
Cohomology of Schemes, 補題 \ref{coherent-lemma-coherent-on-proj} および
\ref{coherent-lemma-recover-tail-graded-module}
により、整数 $q(A, I) \geq 0$ を選んで、任意の $q \geq q(A, I)$ に対し
$H^i(X, \mathcal{O}_X(q)) = 0$ が $i > 0$ について成り立ち、かつ
$H^0(X, \mathcal{O}_X(q)) = I^q$ となるようにできる。

\begin{lemma}
\label{local-cohomology-lemma-compute-tor-Iq}
上の状況で、$q \geq q(A, I)$ および任意の $A$-加群 $M$ に対して
$$
R\Gamma(X, Lp^*\widetilde{M}(q)) \cong M \otimes_A^\mathbf{L} I^q
$$
が $D(A)$ において成り立つ。
\end{lemma}

\begin{proof}
自由分解 $F_\bullet \to M$ を選ぶ。このとき $\widetilde{F}_\bullet$ は
$\widetilde{M}$ の平坦分解である。したがって $Lp^*\widetilde{M}$ は複体
$p^*\widetilde{F}_\bullet$ で与えられる。よって $Lp^*\widetilde{M}(q)$ は複体
$p^*\widetilde{F}_\bullet(q)$ で与えられる。$p^*\widetilde{F}_i(q)$ は、
$\Gamma(X, -)$ に関して非輪状である。これは $q \geq q(A, I)$ の選び方による。
また $\Gamma(X, p^*\widetilde{F}_i(q)) = I^qF_i$ が $q \geq q(A, I)$ の
選び方により成り立つ。したがって $R\Gamma(X, Lp^*\widetilde{M}(q))$ は、
各項が $I^qF_i$ である複体で与えられる。Derived Categories of Schemes, 補題
\ref{perfect-lemma-acyclicity-lemma-global} を参照せよ。
複体 $I^qF_\bullet$ は定義により $M \otimes_A^\mathbf{L} I^q$ を計算するので、
結果が従う。
\end{proof}

\begin{lemma}
\label{local-cohomology-lemma-annihilates}
上の状況において、$t$ を $I$ の生成元数の上界とする。このとき整数
$c = c(A, I) \geq 0$ が存在し、任意の $A$-加群 $M$ に対してコホモロジー層
$H^j(Lp^*\widetilde{M})$ は $I^c$ により零化される（$j \leq -t$ のとき）。
\end{lemma}

\begin{proof}
$I = (a_1, \ldots, a_t)$ とする。この問題は $X$ 上アフィン局所的である。
$1 \leq i \leq t$ に対して、$B_i = A[\frac{I}{a_i}]$ をアフィン・
ブローアップ代数とする。このとき $X$ は、環 $B_i$ のスペクトルによる
アフィン開被覆をもつ（Divisors, 補題
\ref{divisors-lemma-blowing-up-affine} を参照）。
Derived Categories of Schemes, 補題
\ref{perfect-lemma-quasi-coherence-pullback}
に与えられた導来引戻しの記述により、各 $i$ に対してある $c \geq 0$ が存在し、
$$
\text{Tor}_j^A(B_i, M)
$$
が $I^c$ により零化されること（$j \geq t$ のとき）を示せば十分である。
これは補題 \ref{local-cohomology-lemma-annihilates-affine} である。
\end{proof}

\begin{lemma}
\label{local-cohomology-lemma-annihilates-tors}
上の状況において、$t$ を $I$ の生成元数の上界とする。このとき整数
$c = c(A, I) \geq 0$ が存在し、任意の $A$-加群 $M$ に対して Tor 加群
$\text{Tor}_i^A(M, A/I^q)$ は $I^c$ により零化される（$i > t$ かつ
任意の $q \geq 0$ について）。
\end{lemma}

\begin{proof}
$q(A, I)$ を上のものとする。$q \geq q(A, I)$ に対して、補題
\ref{local-cohomology-lemma-compute-tor-Iq} により
$$
R\Gamma(X, Lp^*\widetilde{M}(q)) = M \otimes_A^\mathbf{L} I^q
$$
である。Derived Categories of Schemes, 補題
\ref{perfect-lemma-spectral-sequence} により、有界で収束するスペクトル系列
$$
H^a(X, H^b(Lp^*\widetilde{M}(q))) \Rightarrow
\text{Tor}_{-a - b}^A(M, I^q)
$$
がある。$d$ を Cohomology of Schemes, 補題
\ref{coherent-lemma-vanishing-nr-affines-quasi-separated}
における整数とする（実際には $d = t$ と取れる。Cohomology of Schemes,
補題 \ref{coherent-lemma-vanishing-nr-affines} を参照）。すると
$H^{-i}(X, Lp^*\widetilde{M}(q)) = \text{Tor}_i^A(M, I^q)$ は長さが高々
$d$ の有限フィルトレーションをもち、その次数付き加群は次の加群の
部分商である：
$$
H^a(X, H^{- i - a}(Lp^*\widetilde{M})(q)),\quad
a = 0, 1, \ldots, d - 1
$$
$i \geq t$ ならば、これらの加群はすべて $I^c$ で零化される。ここで
$c = c(A, I)$ は補題 \ref{local-cohomology-lemma-annihilates} におけるものである。実際、
同補題によりコホモロジー層 $H^{- i - a}(Lp^*\widetilde{M})$ はすべて
$I^c$ で零化される。したがって $\text{Tor}_i^A(M, I^q)$ は $I^{dc}$ により
零化される（$q \geq q(A, I)$ かつ $i \geq t$ のとき）。短完全系列
$0 \to I^q \to A \to A/I^q \to 0$ を用いると、$\text{Tor}_i(M, A/I^q)$ は
$I^{dc}$ により零化される（$q \geq q(A, I)$ かつ $i > t$ のとき）。
以上より $I^m$ は、$m = \max(dc, q(A, I) - 1)$ とおけば、
$\text{Tor}_i^A(M, A/I^q)$ を任意の $q \geq 0$ と $i > t$ に対して
望みどおり零化する。
\end{proof}

\begin{lemma}
\label{local-cohomology-lemma-tor-maps-vanish}
$I$ を Noether 環 $A$ のイデアルとし、$t \geq 0$ を $I$ の生成元数の
上界とする。このとき $N, c \geq 0$ が存在し、写像
$$
\text{Tor}_{t + 1}^A(M, A/I^n) \to \text{Tor}_{t + 1}^A(M, A/I^{n - c})
$$
は、任意の $A$-加群 $M$ とすべての $n \geq N$ に対して零となる。
\end{lemma}

\begin{proof}
$c_1$ を補題 \ref{local-cohomology-lemma-annihilates-tors} で得た定数とする。この定数
$c_1$ は $\text{Tor}_i$ に、すべての $i > t$ について同時に使えることに
注意する。

\medskip\noindent
$I = (a_1, \ldots, a_t)$ とする。$A$-加群 $M$ に対して
$$
\ell(M) =
\#\{i \mid 1 \leq i \leq t,\ a_i^{c_1}\text{ は }M\text{ 上でゼロである}\}
$$
と置く。これは $\{0, 1, \ldots, t\}$ の元である。$0 \leq s \leq t$ に関する
降下帰納法により、次の主張 $H_s$ を証明する：ある $N, c \geq 0$ が存在して、
すべての加群 $M$ で $\ell(M) \geq s$ を満たすものに対し、写像
$$
\text{Tor}_{t + 1 + i}^A(M, A/I^n) \to \text{Tor}_{t + 1 + i}^A(M, A/I^{n - c})
$$
は $i = 0, \ldots, s$ およびすべての $n \geq N$ に対して零である。

\medskip\noindent
基底段階：$s = t$ とする。$\ell(M) = t$ ならば $M$ は
$(a_1^{c_1}, \ldots, a_t^{c_1}\}$ により、したがって
$I^{t(c_1 - 1) + 1}$ により零化される。補題
\ref{local-cohomology-lemma-maps-zero-fixed-torsion} より、$H_t$ は、$c = N$ を補題で得られる整数
$c(A, I, t(c_1 - 1) + 1, t + 1), \ldots, c(A, I, t(c_1 - 1) + 1, 2t + 1)$
の最大値に取れば成り立つ。

\medskip\noindent
帰納段階。$0 \leq s < t$ とし、$N, c$ は $H_{s + 1}$ におけるものとする。
加群 $M$ で $\ell(M) = s$ を満たすものを考える。このとき、ある $i$ を選んで
$a_i^{c_1}$ が $M$ 上で非零となるようにできる。したがって
$\ell(M[a_i^c]) \geq s + 1$ および $\ell(M/a_i^{c_1}M) \geq s + 1$ であり、
これらに帰納法の仮定を適用できる。完全系列
$$
0 \to M[a_i^{c_1}] \to M \xrightarrow{a_i^{c_1}} M \to M/a_i^{c_1}M \to 0
$$
を考える。中央の矢印の像を $E \subset M$ と記す。対応する Tor 加群の図式
$$
\xymatrix{
&
&
\text{Tor}_{i + 1}(M/a_i^{c_1}M, A/I^q) \ar[d] \\
\text{Tor}_i(M[a_i^{c_1}], A/I^q) \ar[r] &
\text{Tor}_i(M, A/I^q) \ar[r] \ar[rd]^0 &
\text{Tor}_i(E, A/I^q) \ar[d] \\
&
&
\text{Tor}_i(M, A/I^q)
}
$$
を得る。その行と列は（任意の $q$ に対して）完全である。南東向きの矢印は
$c_1$ の選び方により零である。したがって加群
$\text{Tor}_i(M, A/I^q)$ は、$\text{Tor}_i(M[a_i^{c_1}], A/I^q)$ の
ある商加群と $\text{Tor}_{i + 1}(M/a_i^{c_1}M, A/I^q)$ のある部分加群との
間に挟まれる。ゆえに $H_s$ は、$N$ を $N + c$ で、$c$ を $2c$ で
置き換えれば成り立つ。細部の一部は省略する。
\end{proof}

\begin{proposition}
\label{local-cohomology-proposition-uniform-artin-rees}
$I$ を Noether 環 $A$ のイデアルとし、$t \geq 0$ を $I$ の生成元数の上界とする。
このとき $N, c \geq 0$ が存在し、$n \geq N$ に対して写像
$$
A/I^n \to A/I^{n - c}
$$
は、$e = t$ とした補題 \ref{local-cohomology-lemma-characterize-vanishing-tor-ext-above-e} の
同値な条件を満たす。
\end{proposition}

\begin{proof}
補題 \ref{local-cohomology-lemma-tor-maps-vanish} と
\ref{local-cohomology-lemma-characterize-vanishing-tor-ext-above-e} から直ちに従う。
\end{proof}

\begin{remark}
\label{local-cohomology-remark-better-bound}
論文 \cite{AHS} は、ほかにも多くのことを示すとともに、$A$ が局所環ならば
命題 \ref{local-cohomology-proposition-uniform-artin-rees} において $e = t$ を
$e = \dim(A)$ で置き換えても成り立つことを示している。補題
\ref{local-cohomology-lemma-cd-sequence-Koszul} を見ると、命題
\ref{local-cohomology-proposition-uniform-artin-rees} において $e = t$ を
$e = \text{cd}(A, I)$ で置き換えても成り立つかを問うのは自然である。
これは分かっていない。
\end{remark}

\begin{remark}
\label{local-cohomology-remark-strict-pro-isomorphism}
$I$ を Noether 環 $A$ のイデアルとする。$I = (f_1, \ldots, f_r)$ とする。
$K_n^\bullet$ を More on Algebra, 状況
\ref{more-algebra-situation-koszul} における $f_1^n, \ldots, f_r^n$ 上の
Koszul 複体と記し、対応する対象を $K_n \in D(A)$ と記す。
$M^\bullet$ を有限 $A$-加群からなる有界複体とし、対応する対象を
$M \in D(A)$ と記す。$D(A)$ における次の逆系を考える：
\begin{enumerate}
\item $M^\bullet/I^nM^\bullet$、すなわち各項が $M^i/I^nM^i$ である複体、
\item $M \otimes_A^\mathbf{L} A/I^n$,
\item $M \otimes_A^\mathbf{L} K_n$、および
\item $M \otimes_P^\mathbf{L} P/J^n$（下を参照）。
\end{enumerate}
これらの逆系はすべて pro 対象として同型である。(2) と (3) の間の同型は
More on Algebra, 補題 \ref{more-algebra-lemma-sequence-Koszul-complexes}
から従う。(1) と (2) の間の同型は More on Algebra, 補題
\ref{more-algebra-lemma-derived-completion-plain-completion} に与えられている。
最後のものについては下を参照されたい。

\medskip\noindent
しかし、これら pro 系の同型が「強い」ものであるかを問うこともできる。
この用語と問いは \cite[61, 62 頁]{quillenhomology} の議論に関係している。
すなわち、圏 $\mathcal{C}$ が与えられたとき、対象を逆系 $(X_n)$ とし、射
$(X_n) \to (Y_n)$ を組 $(c, \varphi_n)$ により表す。この組は
$c \geq 0$ と、射 $\varphi_n : X_n \to Y_{n - c}$
（すべての $n \geq c$ について）からなり、明らかな両立条件を満たし、
ある同値関係（本質的には $c$ を大きくすることで与えられる）で同一視する。
これにより「強い pro 圏」を定義できる。そこで、上の逆系が
この強い pro 圏で同型であるかを問う。

\medskip\noindent
これは、$M = A[0]$ の場合でさえ (1) と (3) については明らかに成り立ちえない。
実際、系 $H^0(K_n) = A/(f_1^n, \ldots, f_r^n)$ は一般に、加群の圏において系
$A/I^n$ と強い意味で pro 同型ではない。例えば
$A = \mathbf{Z}[x_1, \ldots, x_r]$ および $f_i = x_i$ と取れば、
$H^0(K_n)$ は $I^{r(n - 1)}$ により零化されない。\footnote{もちろん、対象が
逆系であり、射が組 $(r, c, \varphi_n)$ により与えられる圏において、これらの
pro 系が同型であるかを問うことはできる。この組は $r \geq 1$、$c \geq 0$
および写像 $\varphi_n : X_{rn} \to Y_{n - c}$（$n \geq c$ に対する）からなる。}

\medskip\noindent
実際、上の結果から、More on Algebra, 補題
\ref{more-algebra-lemma-derived-completion-plain-completion} で論じた (2) から
(1) への自然な写像は、強い意味で pro 同型であることが分かる。証明を概説する。
愚鈍切断を用いる標準的な議論により、まず $M^\bullet$ が一つの有限
$A$-加群 $M$ を次数 $0$ に置いて得られる場合に帰着する。命題
\ref{local-cohomology-proposition-uniform-artin-rees} における $N, c \geq 0$ を選ぶ。同命題より、
$n \geq N$ に対して系 (2) の遷移写像の分解
$$
M \otimes_A^\mathbf{L} A/I^n
\to
\tau_{\geq -t}(M \otimes_A^\mathbf{L} A/I^n)
\to
M \otimes_A^\mathbf{L} A/I^{n - c}
$$
を得る。一方、More on Algebra, 補題
\ref{more-algebra-lemma-tor-strictly-pro-zero} により、別の定数
$c' = c'(M) \geq 0$ であって、写像
$\text{Tor}_i^A(M, A/I^{n'}) \to \text{Tor}_i(M, A/I^{n' - c'})$ が
$i = 1, 2, \ldots, t$ および $n' \geq c'$ に対して零となるものを見つけられる。
すると Derived Categories, 補題
\ref{derived-lemma-trick-vanishing-composition} より、写像
$$
\tau_{\geq -t}(M \otimes_A^\mathbf{L} A/I^{n + tc'})
\to
\tau_{\geq -t}(M \otimes_A^\mathbf{L} A/I^n)
$$
は $M \otimes_A^\mathbf{L}A/I^{n + tc'} \to M/I^{n + tc'}M$ を経由する。
これを先の結果と合わせると、分解
$$
M \otimes_A^\mathbf{L}A/I^{n + tc'} \to M/I^{n + tc'}M
\to M \otimes_A^\mathbf{L} A/I^{n - c}
$$
を得るので、望むことが従う。この結果が必要になったときには、改めて注意深く
定式化し、詳細な証明を与えることにする。

\medskip\noindent
(4) については、Noether 環 $P$、環準同型 $P \to A$、およびイデアル
$J \subset P$ で $I = JA$ を満たすものがあるとする。More on Algebra, 第
\ref{more-algebra-section-derived-base-change} 節により、関手
$M \otimes_P^\mathbf{L} - : D(P) \to D(A)$ と、(4) における逆系
$M \otimes_P^\mathbf{L} P/J^n$ を $D(A)$ 内に得る。$P$ が Noether 環ならば、Koszul 複体と
比較できるため、(4) の系は (1) の系と pro 同型である。$P \to A$ が有限ならば、
(4) の系は (2) の系と強い意味で pro 同型である。実際、逆系
$A \otimes_P^\mathbf{L} P/J^n$ は、上の議論により逆系 $A/I^n$ と強い意味で
pro 同型であり、さらに
$$
M \otimes_P^\mathbf{L} P/J^n = M \otimes_A^\mathbf{L}
(A \otimes_P^\mathbf{L} P/J^n)
$$
が More on Algebra, 補題 \ref{more-algebra-lemma-derived-base-change} により
成り立つからである。

\medskip\noindent
(4) の標準的な例は、$P = \mathbf{Z}[x_1, \ldots, x_r]$、写像
$P \to A$ で $x_i$ を $f_i$ に送るもの、および $J = (x_1, \ldots, x_r)$ を
取るものである。この場合、
$$
M \otimes_P^\mathbf{L} P/J^n =
M \otimes_{A[x_1, \ldots, x_r]}^\mathbf{L}
A[x_1, \ldots, x_r]/(x_1, \ldots, x_r)^n
$$
が成り立つことを示し、上で論じた場合の一つに帰着する（ただしこの場合は、
$A[x_1, \ldots, x_r]/(x_1, \ldots, x_r)^n$ が高々 $r$ の Tor 次元を
すべての $n$ に対してもつので、命題 \ref{local-cohomology-proposition-uniform-artin-rees} を用いる段階を
避けられ、実際にはより容易である）。この場合は
\cite[命題 3.5.1]{BS} の証明で論じられている。
\end{remark}








\section{一様性について III}
\label{local-cohomology-section-uniform}

\noindent
本節では Noether 環 $A$ とイデアル $I \subset A$ を固定する。目標は補題
\ref{local-cohomology-lemma-bound-two-term-complex} を証明することであり、この補題は後の章で
持上げ問題を解くために用いる（Algebraization of Formal Spaces, 補題
\ref{restricted-lemma-get-morphism-general-better} を参照）。

\medskip\noindent
本節を通じて、
$$
p : X \to \Spec(A)
$$
により $\Spec(A)$ のイデアル $I$ におけるブローアップを表す。言い換えると、
$X$ は $\text{Proj}$、すなわち Rees 代数 $\bigoplus_{n \geq 0} I^n$ から得られる
射影スキームである。さらにファイバー積
$$
\xymatrix{
Y \ar[r] \ar[d] & X \ar[d]^p \\
\Spec(A/I) \ar[r] & \Spec(A)
}
$$
も考える。このとき $Y$ はブローアップの例外因子であり、したがって $X$ 上の
有効 Cartier 因子であって
$\mathcal{O}_X(-1) = \mathcal{O}_X(Y)$ を満たす。$\text{Proj}$ を取る操作は
基底変換と可換なので、
$$
Y = \text{Proj}(\bigoplus\nolimits_{n \geq 0} I^n/I^{n + 1}) = \text{Proj}(S)
$$
を得る。ここで $S = \text{Gr}_I(A) = \bigoplus_{n \geq 0} I^n/I^{n + 1}$ である。

\medskip\noindent
次の記号を用いる：
$d = d(S) = d(\text{Gr}_I(A)) = d(\bigoplus_{n \geq 0} I^n/I^{n + 1})$
は $p$ のファイバーの次元の最大値を表す（$0$ と置くのは $X = \emptyset$
の場合である）。これは良定義である。実際、次が成り立つ。
\begin{enumerate}
\item $d \leq t - 1$ であるのは、$I = (a_1, \ldots, a_t)$ ならば
$X \subset \mathbf{P}^{t - 1}_A$ だからである。
\item $d$ は、$\text{Proj}(S) \to \Spec(S_0)$ のファイバーの次元の最大値でもある
（$Y$ が空でない場合）。また $d = 0$ であるのは $Y = \emptyset$ の場合である
（同値な条件は $S = 0$、さらに同値な条件は
$I = A$ である）。
\end{enumerate}
したがって $d$ は $S = \text{Gr}_I(A)$ の同型類のみに依存する。
$H^i(X, \mathcal{F}) = 0$ は、任意の連接 $\mathcal{O}_X$-加群 $\mathcal{F}$ と
$i > d$ に対して Cohomology of Schemes, 補題
\ref{coherent-lemma-higher-direct-images-zero-above-dimension-fibre} および
\ref{coherent-lemma-quasi-coherence-higher-direct-images-application} により成り立つ。
当然、$Y$ 上の連接加群についても同じことが成り立つ。

\medskip\noindent
次の記号を用いる：
$q = q(S) = q(\text{Gr}_I(A)) = q(\bigoplus_{n \geq 0} I^n/I^{n + 1})$
は次のように定義される整数を表す。代数
$S = \bigoplus_{n \geq 0} I^n/I^{n + 1}$
は、次数 $1$ で生成される Noether 次数環であり、その係数環は次数 $0$ の部分で
あることに
注意する。したがって Cohomology of Schemes, 補題
\ref{coherent-lemma-coherent-on-proj} および
\ref{coherent-lemma-recover-tail-graded-module} により、$q(S)$ を、最小の整数
$q(S) \geq 0$ であって、すべての $q \geq q(S)$ に対して
$H^i(Y, \mathcal{O}_Y(q)) = 0$ が
$1 \leq i \leq d$ について成り立ち、かつ
$H^0(Y, \mathcal{O}_Y(q)) = I^q/I^{q + 1}$ が成り立つものとして定義できる。
（$S = 0$ ならば $q(S) = 0$ とする。）

\medskip\noindent
$n \geq 1$ に対して有効 Cartier 因子 $nY$ を考えることができ、これを $Y_n$ と記す。

\begin{lemma}
\label{local-cohomology-lemma-bound-q-and-d}
上の $q_0 = q(S)$ と $d = d(S)$ に対して、次が成り立つ。
\begin{enumerate}
\item $n \geq 1$、$q \geq q_0$、および $i > 0$ ならば
$H^i(X, \mathcal{O}_{Y_n}(q)) = 0$ である。
\item $n \geq 1$ かつ $q \geq q_0$ ならば
$H^0(X, \mathcal{O}_{Y_n}(q)) = I^q/I^{q + n}$ である。
\item $q \geq q_0$ かつ $i > 0$ ならば
$H^i(X, \mathcal{O}_X(q)) = 0$ である。
\item $q \geq q_0$ ならば $H^0(X, \mathcal{O}_X(q)) = I^q$ である。
\end{enumerate}
\end{lemma}

\begin{proof}
$I = A$ ならば $X$ はアフィンであり、主張は自明である。したがって
$I \not = A$ と仮定してよいし、そう仮定する。このとき $Y$ と $X$ は空でない
スキームである。

\medskip\noindent
(1) と (2) を $n$ に関する帰納法で証明する。基底段階 $n = 1$ は、
$q_0$ の定義そのものである（$Y_1 = Y$ だからである）。
$\mathcal{O}_X(1) = \mathcal{O}_X(-Y)$ であることを思い出す。したがって短完全系列
$$
0 \to \mathcal{O}_{Y_n}(1) \to \mathcal{O}_{Y_{n + 1}} \to \mathcal{O}_Y \to 0
$$
がある。したがって $i > 0$ に対して
$$
H^i(X, \mathcal{O}_{Y_n}(q + 1)) \to
H^i(X, \mathcal{O}_{Y_{n + 1}}(q)) \to
H^i(X, \mathcal{O}_{Y}(q))
$$
を得て、両端の項の既知の消滅から中央の項の望む消滅を得る。$i = 0$ に対しては
可換図式
$$
\xymatrix{
0 \ar[r] &
I^{q + 1}/I^{q + 1 + n} \ar[d] \ar[r] &
I^q/I^{q + 1 + n} \ar[d] \ar[r] &
I^q/I^{q + 1} \ar[d] \ar[r] &
0 \\
0 \ar[r] &
H^0(X, \mathcal{O}_{Y_n}(q + 1)) \ar[r] &
H^0(X, \mathcal{O}_{Y_{n + 1}}(q)) \ar[r] &
H^0(Y, \mathcal{O}_Y(q)) \ar[r] &
0
}
$$
を得る。その行は $q \geq q_0$ に対して完全である（下の行については、長完全
コホモロジー系列の次の項が $q \geq q_0$ に対して消滅することに注意する）。
$q \geq q_0$ なので左右の垂直矢印は同型であり、中央の矢印も同型だと結論する。

\medskip\noindent
同様である (3) と (4) の証明は省略する。実際、形式関手の定理を用いれば
(1) と (2) から (3) と (4) を導ける（ただし、これは過剰であろう）。
\end{proof}

\noindent
記法を導入する。$n \geq c \geq 0$ が与えられたとき、{\it $(A, n, c)$-加群}
とは、有限 $A$-加群 $M$ であって、$I^n$ により零化され、かつ
$A/I^n$-加群として $I^c/I^n$-射影的であるものをいう（More on Algebra, 第
\ref{more-algebra-section-near-projective} 節を参照）。

\medskip\noindent
次の記号の濫用を用いる。$A$-加群 $M$ が与えられたとき、$p^*M$ により、
$p$ で引き戻して得られる擬連接加群、すなわち擬連接加群 $\widetilde{M}$
（$\Spec(A)$ 上で $M$ に対応するもの）の引戻しを表す。例えば
$\mathcal{O}_{Y_n} = p^*(A/I^n)$ である。短完全系列
$0 \to K \to L \to M \to 0$ が与えられ、その各項が $A$-加群ならば、完全系列
$$
p^*K \to p^*L \to p^*M \to 0
$$
を得る。実際、$\widetilde{\ }$ は完全関手であり、$p^*$ は右完全関手である。

\begin{lemma}
\label{local-cohomology-lemma-almost-exactness}
$0 \to K \to L \to M \to 0$ を $A$-加群の短完全系列で、$K$ と $L$ は
$I^n$ により零化され、$M$ は $(A, n, c)$-加群であるものとする。このとき
$p^*K \to p^*L$ の核はスキーム論的に $Y_c$ 上に台をもつ。
\end{lemma}

\begin{proof}
$\Spec(B) \subset X$ をアフィン開集合とする。完全系列を $\Spec(B)$ 上に制限した
ものは、$B$-加群の系列
$$
K \otimes_A B \to L \otimes_A B \to M \otimes_A B \to 0
$$
に対応し、これは系列
$$
K \otimes_{A/I^n} B/I^nB \to
L \otimes_{A/I^n} B/I^nB \to
M \otimes_{A/I^n} B/I^nB \to 0
$$
と同型である。したがって最初の写像の核は加群
$\text{Tor}_1^{A/I^n}(M, B/I^nB)$ の像である。例外因子 $Y$ は
$I\mathcal{O}_X$ により切り出されることを思い出す。ゆえに
$\text{Tor}_1^{A/I^n}(M, B/I^nB)$ が $I^c$ により零化されることを示せば十分である。
$a \in I^c$ による $M$ 上の乗法は有限自由 $A/I^n$-加群を経由するので、これは
明らかである。
\end{proof}

\noindent
標準写像 $\mathcal{O}_X \to \mathcal{O}_X(1)$ があり、これはちょうど $Y$ に沿って消える。
したがって、任意の連接 $\mathcal{O}_X$-加群 $\mathcal{F}$ に対して常に標準写像
$\mathcal{F}(q) \to \mathcal{F}(q + n)$ があり、ここで任意の
$q \in \mathbf{Z}$ と $n \geq 0$ を取れる。

\begin{lemma}
\label{local-cohomology-lemma-annihilated}
$\mathcal{F}$ を連接 $\mathcal{O}_X$-加群とする。このとき、$\mathcal{F}$ が
スキーム論的に $Y_c$ 上に台をもつことと、標準写像
$\mathcal{F} \to \mathcal{F}(c)$ が零であることとは同値である。
\end{lemma}

\begin{proof}
$\mathcal{O}_X \to \mathcal{O}_X(1)$ がちょうど $Y$ に沿って消えるからである。
\end{proof}

\begin{lemma}
\label{local-cohomology-lemma-vanishing-coh-almost-projective}
上の $q_0 = q(S)$ と $d = d(S)$ に対し、整数 $n \geq c \geq 0$、
$(A, n, c)$-加群 $M$、添字 $i \in \{0, 1, \ldots, d\}$、および整数 $q$ が
与えられているとする。このとき、次の各場合に分ける。
\begin{enumerate}
\item $i = d \geq 1$ かつ $q \geq q_0$ の場合、
$H^d(X, p^*M(q)) = 0$ である。
\item $i = d - 1 \geq 1$ かつ $q \geq q_0$ の場合、
$H^{d - 1}(X, p^*M(q)) = 0$ である。
\item $d - 1 > i > 0$ かつ $q \geq q_0 + (d - 1 - i)c$ の場合、写像
$H^i(X, p^*M(q)) \to H^i(X, p^*M(q - (d - 1 - i)c))$
は零である。
\item $i = 0$、$d \in \{0, 1\}$、かつ $q \geq q_0$ の場合、全射
$$
I^qM \longrightarrow H^0(X, p^*M(q))
$$
がある。
\item $i = 0$、$d > 1$、かつ $q \geq q_0 + (d - 1)c$ の場合、写像
$$
H^0(X, p^*M(q)) \to H^0(X, p^*M(q - (d - 1)c))
$$
の像は、標準写像
$I^{q - (d - 1)c}M \to H^0(X, p^*M(q - (d - 1)c))$ の像に含まれる。
\end{enumerate}
\end{lemma}

\begin{proof}
$M$ を $(A, n, c)$-加群とする。短完全系列
$$
0 \to K \to (A/I^n)^{\oplus r} \to M \to 0
$$
を選ぶ。以下では $K$ が $(A, n, c)$-加群であることを用いる（More on Algebra,
補題 \ref{more-algebra-lemma-ses-near-projective} を参照）。対応する完全系列
$$
p^*K \to (\mathcal{O}_{Y_n})^{\oplus r} \to p^*M \to 0
$$
を考える。これを短完全系列
$$
0 \to \mathcal{F} \to p^*K \to \mathcal{G} \to 0
\quad\text{および}\quad
0 \to \mathcal{G} \to (\mathcal{O}_{Y_n})^{\oplus r} \to p^*M \to 0
$$
に分ける。補題 \ref{local-cohomology-lemma-almost-exactness} により、連接加群 $\mathcal{F}$ は
スキーム論的に $Y_c$ 上に台をもつ。

\medskip\noindent
(1) の証明。$d > 0$ と仮定する。$H^d(X, p^*M(q)) = 0$ を
$q \geq q_0$ に対して証明しなければならない。$\mathcal{G}$ の捩りの
$> d$ 次コホモロジーが消滅することと、上の第二の短完全系列に付随する長完全
コホモロジー系列とにより、$H^d(X, \mathcal{O}_{Y_n}(q)) = 0$ を証明すれば十分である。
これは補題 \ref{local-cohomology-lemma-bound-q-and-d} により成り立つ。

\medskip\noindent
(2) の証明。$d > 1$ と仮定する。$H^{d - 1}(X, p^*M(q)) = 0$ を
$q \geq q_0$ に対して証明しなければならない。前段落と同様に論じると、
$H^d(X, \mathcal{G}(q)) = 0$ を示せば十分だと分かる。第一の短完全系列と、
$\mathcal{F}$ の捩りの $> d$ 次コホモロジーが消滅することとを用いれば、
$H^d(X, p^*K(q))$ が零であることを示せば十分である。これは (1) と、$K$ が
$(A, n, c)$-加群であること（上を参照）から従う。

\medskip\noindent
(3) の証明。$0 < i < d - 1$ とし、主張が $i + 1$ に対して成り立つと仮定する。
ただし $i = d - 2$ の場合には主張 (2) がある。上の第二の短完全系列に付随する
長完全コホモロジー系列を用いると、単射
$$
H^i(X, p^*M(q - (d - 1 - i)c)) \subset
H^{i + 1}(X, \mathcal{G}(q - (d - 1 - i)c))
$$
を得る。実際、$q - (d - 1 - i)c \geq q_0$ から
$H^i(X, \mathcal{O}_{Y_n}(q - (d - 1 - i)c))$ の消滅が従う（上を参照）。
したがって、写像
$H^{i + 1}(X, \mathcal{G}(q)) \to H^{i + 1}(X, \mathcal{G}(q - (d - 1 - i)c))$
が零であることを示せば十分である。これを調べるため、完全系列間の写像
$$
\xymatrix{
H^{i + 1}(X, p^*K(q)) \ar[r] \ar[d] &
H^{i + 1}(X, \mathcal{G}(q)) \ar[r] \ar[d] \ar@{..>}[ld] &
H^{i + 2}(X, \mathcal{F}(q)) \ar[d] \\
H^{i + 1}(X, p^*K(q - c)) \ar[r] \ar[d] &
H^{i + 1}(X, \mathcal{G}(q - c)) \ar[r] \ar[d] &
H^{i + 2}(X, \mathcal{F}(q - c)) \\
H^{i + 1}(X, p^*K(q - (d - 1 - i)c)) \ar[r] &
H^{i + 1}(X, \mathcal{G}(q - (d - 1 - i)c))
}
$$
を考える。$\mathcal{F}$ はスキーム論的に $Y_c$ 上に台をもつので、標準写像
$\mathcal{G}(q) \to \mathcal{G}(q - c)$ は補題 \ref{local-cohomology-lemma-annihilated} により
$p^*K(q - c)$ を経由する。これが図式の点線矢印を与える。（実際、証明のためには、
最右端の垂直矢印が零であることに注意すれば、集合の写像として点線矢印を得るのに
十分である。）したがって、
$H^{i + 1}(X, p^*K(q - c)) \to
H^{i + 1}(X, p^*K(q - (d - 1 - i)c))$
が零であることを示せば十分である。$i = d - 2$ ならば、この矢印の始域は、
$q - c \geq q_0$ であり $K$ が $(A, n, c)$-加群であるため、(2) により零である。
$i < d - 2$ ならば、$K$ が $(A, n, c)$-加群なので、帰納法の仮定からこの写像は
実際に零であることが従う。なぜなら
$q - c - (q - (d - 1 - i)c) = (d - 2 - i)c = (d - 1 - (i + 1))c$ であり、かつ
$q - c \geq q_0 + (d - 1 - (i + 1))c$ だからである。これで (3) の証明が完了する。

\medskip\noindent
(4) の証明。$d \in \{0, 1\}$ かつ $q \geq q_0$ と仮定する。このとき第一の
短完全系列から全射
$H^1(X, p^*K(q)) \to H^1(X, \mathcal{G}(q))$
を得て、この矢印の始域は場合 (1) により零である。したがって、すべての
$q \in \mathbf{Z}$ に対して写像
$$
H^0(X, (\mathcal{O}_{Y_n})^{\oplus r}(q))
\longrightarrow
H^0(X, p^*M(q))
$$
は全射である。$q \geq q_0$ に対して、その始域は補題
\ref{local-cohomology-lemma-bound-q-and-d} により $(I^q/I^{q + n})^{\oplus r}$ に等しく、これで
主張は容易に従う。

\medskip\noindent
(5) の証明。$d > 1$ と仮定する。(4) の証明と同様に論じると、
$$
H^0(X, p^*M(q))
\longrightarrow
H^0(X, p^*M(q - (d - 1)c))
$$
の像が
$$
H^0(X, (\mathcal{O}_{Y_n})^{\oplus r}(q - (d - 1)c))
\longrightarrow
H^0(X, p^*M(q - (d - 1)c))
$$
の像に含まれることを示せば十分だと分かる。上の包含を示すには、
$\sigma \in H^0(X, p^*M(q))$ で、その境界が
$\xi \in H^1(X, \mathcal{G}(q))$ であるものに対して、$\xi$ の
$H^1(X, \mathcal{G}(q - (d - 1)c))$ における像が零であることを示せば十分である。
これは (3) の証明とまったく同じ議論から従う。
\end{proof}

\begin{remark}
\label{local-cohomology-remark-duals}
組 $(M, n)$ が与えられ、それが整数 $n \geq 0$ と有限 $A/I^n$-加群 $M$ から
なるとき、$M^\vee = \Hom_{A/I^n}(M, A/I^n)$ と置く。組 $(\mathcal{F}, n)$ が
与えられ、それが整数 $n$ と連接 $\mathcal{O}_{Y_n}$-加群 $\mathcal{F}$ からなるとき、
$$
\mathcal{F}^\vee =
\SheafHom_{\mathcal{O}_{Y_n}}(\mathcal{F}, \mathcal{O}_{Y_n})
$$
と置く。上のような $(M, n)$ が与えられたとき、標準写像
$$
can : p^*(M^\vee) \longrightarrow (p^*M)^\vee
$$
がある。実際、表示
$(A/I^n)^{\oplus s} \to (A/I^n)^{\oplus r} \to M \to 0$
を選べば、表示
$\mathcal{O}_{Y_n}^{\oplus s} \to \mathcal{O}_{Y_n}^{\oplus r} \to
p^*M \to 0$ を得る。双対を取ると完全系列
$$
0 \to M^\vee \to (A/I^n)^{\oplus r} \to (A/I^n)^{\oplus s}
$$
および
$$
0 \to (p^*M)^\vee \to
\mathcal{O}_{Y_n}^{\oplus r} \to
\mathcal{O}_{Y_n}^{\oplus s}
$$
を得る。第一の系列を $p$ で引き戻すと、望む写像 $can$ を得る。この写像の構成は
有限 $A/I^n$-加群 $M$ に関して関手的である。$can$ の核と余核は、スキーム論的に
$Y_c$ 上に台をもつ（$M$ が $(A, n, c)$-加群の場合）。実際、この場合、
$a \in I^c$ に対して写像 $a : M \to M$ は有限自由 $A/I^n$-加群を経由し、
その加群に対しては $can$ が同型である。したがって $a$ は $can$ の核と余核を
零化する。
\end{remark}

\begin{lemma}
\label{local-cohomology-lemma-factor-hom}
上の $q_0 = q(S)$ と $d = d(S)$ に対し、$M$ を $(A, n, c)$-加群とし、
$\varphi : M \to I^n/I^{2n}$ を $A$-線形写像とする。
$n \geq \max(q_0 + (1 + d)c, (2 + d)c)$ と仮定し、$d = 0$ ならばさらに
$n \geq q_0 + 2c$ と仮定する。このとき合成
$$
M \xrightarrow{\varphi} I^n/I^{2n} \to
I^{n - (1 + d)c}/I^{2n - (1 + d)c}
$$
は $\sum a_i \psi_i$ の形である。ここで $a_i \in I^c$ かつ
$\psi_i : M \to I^{n - (2 + d)c}/I^{2n - (2 + d)c}$ である。
\end{lemma}

\begin{proof}
$d > 1$ の場合。写像 $p^*(I^q) \to \mathcal{O}_X(q)$ の両立する系が
$q \geq 0$ に対してあるので、標準写像
$p^*(I^q/I^{q + \nu}) \to \mathcal{O}_{Y_\nu}(q)$ が $\nu \geq 0$ に対してある。
これを用い、さらに
$\varphi$ を引き戻すと、写像
$$
\chi : p^*M \longrightarrow \mathcal{O}_{Y_n}(n)
$$
であって、合成
$M \to H^0(X, p^*M) \to H^0(X, \mathcal{O}_{Y_n}(n))$
が、与えられた準同型 $\varphi$ と写像
$I^n/I^{2n} \to H^0(X, \mathcal{O}_{Y_n}(n))$ との合成であるものを得る。
$\mathcal{O}_{Y_n}(n)$ は $Y_n$ 上可逆なので、線形写像 $\chi$ は切断
$$
\sigma \in \Gamma(X, (p^*M)^\vee(n))
$$
を定める。記法は注意 \ref{local-cohomology-remark-duals} のとおりである。
注意 \ref{local-cohomology-remark-duals} の議論により、$can : p^*(M^\vee) \to (p^*M)^\vee$ の
余核と核はスキーム論的に $Y_c$ 上に台をもつ。補題
\ref{local-cohomology-lemma-annihilated} により、写像
$(p^*M)^\vee(n) \to (p^*M)^\vee(n - 2c)$ は
$p^*(M^\vee)(n - 2c)$ を経由する。細部を一つ省略した。したがって $\sigma$ の
$\Gamma(X, (p^*M)^\vee(n - 2c))$ における像は、ある元
$$
\sigma' \in \Gamma(X, p^*(M^\vee)(n - 2c))
$$
から来る。補題 \ref{local-cohomology-lemma-vanishing-coh-almost-projective} の (5)、More on Algebra,
補題 \ref{more-algebra-lemma-dual-near-projective} により $M^\vee$ が
$(A, n, c)$-加群であること、および
$n \geq q_0 + (1 + d)c$ なので $n - 2c \geq q_0 + (d - 1)c$ であることから、
$\sigma'$ の $H^0(X, p^*M^\vee(n - (1 + d)c))$ における像は、ある元 $\tau$ の
$I^{n - (1 + d)c}M^\vee$ における像であると分かる。
$\tau = \sum a_i \tau_i$ と書く。ここで
$\tau_i \in I^{n - (2 + d)c}M^\vee$ である。これは
$n - (2 + d)c \geq 0$ なので意味をもつ。すると $\tau_i$ は加群の準同型
$\psi_i : M \to I^{n - (2 + d)c}/I^{2n - (2 + d)c}$ を、評価写像
$M \otimes M^\vee \to A/I^n$ を用いて定める。

\medskip\noindent
これが実際に機能することを証明しよう\footnote{この事実について、より泥臭くない
証明を読者が提案してくれることを期待する。}。$z \in M$ を選び、$\varphi(z)$ と
$\sum a_i \psi_i(z)$ が $I^{n - (1 + d)c}/I^{2n - (1 + d)c}$ において同じ像を
もつことを示す。まず元 $z$ は写像
$p^*z : \mathcal{O}_{Y_n} \to p^*M$ を定め、この写像と $\chi$ との合成は、写像
$\mathcal{O}_{Y_n} \to \mathcal{O}_{Y_n}(n)$、すなわち $\varphi(z)$ に、写像
$I^n/I^{2n} \to \Gamma(\mathcal{O}_{Y_n}(n))$ を介して対応するものに等しい。
次に $z$ と $p^*z$ は評価写像 $e_z : M^\vee \to A/I^n$ および
$e_{p^*z} : (p^*M)^\vee \to \mathcal{O}_{Y_n}$ を定める。$\chi(p^*z)$ は
$\varphi(z)$ に対応する切断なので、$e_{p^*z}(\sigma)$ は $\varphi(z)$ に対応する
切断である。ここおよび以下では記号を濫用する。加群の写像
$a : \mathcal{F} \to \mathcal{G}$ が $X$ 上で与えられたとき、捩った加群の対応する写像も
$a : \mathcal{F}(t) \to \mathcal{F}(t)$ と記す。図式
$$
\xymatrix{
p^*(M^\vee) \ar[d]_{can} \ar[r]_{p^*e_z} & \mathcal{O}_{Y_n} \ar@{=}[d] \\
(p^*M)^\vee \ar[r]^{e_{p^*z}} & \mathcal{O}_{Y_n}
}
$$
は構成 $can$ の関手性により可換である。したがって
$(p^*e_z)(\sigma')$ は $\Gamma(Y_n, \mathcal{O}_{Y_n}(n - 2c))$ において、
$\varphi(z)$ の $I^{n - 2c}/I^{2n - 2c}$ における像に対応する切断である。
次に、$\sigma'$ は $\sum a_i \tau_i$ の
$H^0(X, p^*M^\vee(n - (1 + d)c))$ における像へ写る。このことから、
$(p^*e_z)(\sum a_i \tau_i) = \sum a_i p^*e_z(\tau_i)$ は
$\Gamma(Y_n, \mathcal{O}_{Y_n}(n - (1 + d)c))$ において、$\varphi(z)$ の
$I^{n - (1 + d)c}/I^{2n - (1 + d)c}$ における像に対応する切断である。
$\psi_i$ は $\tau_i$ から評価写像を用いて定義されることを思い出す。したがって、
$$
\chi_i : p^*M \longrightarrow \mathcal{O}_{Y_n}(n - (2 + d)c)
$$
により $\psi_i$ から得られる写像を表すと、上と同じ理由により、$\psi_i(z)$ に
対応する切断は $\chi_i(p^*z) = e_{p^*z}(\chi_i) = p^*e_z(\tau_i)$ である。
したがって、$\varphi(z)$ の $\Gamma(Y_n, \mathcal{O}_{Y_n}(n - (1 + d)c))$ に
おける像は、$\sum a_i\psi_i(z)$ の像に等しい。$n - (1 + d)c \geq q_0$ なので
$\Gamma(Y_n, \mathcal{O}_{Y_n}(n - (1 + d)c)) =
I^{n - (1 + d)c}/I^{2n - (1 + d)c}$ であることが補題
\ref{local-cohomology-lemma-bound-q-and-d} から従い、望む両立性が成り立つと結論する。

\medskip\noindent
$d = 1$ の場合。ここでは上と同様に論じて
$$
\chi : p^*M \longrightarrow \mathcal{O}_{Y_n}(n),\quad
\sigma \in \Gamma(X, (p^*M)^\vee(n)),\quad
\sigma' \in \Gamma(X, p^*(M^\vee)(n - 2c)),
$$
を得る。次に補題 \ref{local-cohomology-lemma-vanishing-coh-almost-projective} の (4) を用いると、
$\sigma'$ はある元 $\tau \in I^{n - 2c}M^\vee$ の像であると分かる。議論の残りは
同じである。

\medskip\noindent
$d = 0$ の場合。議論は $d = 1$ の場合とまったく同じである。
\end{proof}

\begin{lemma}
\label{local-cohomology-lemma-bound-two-term-complex}
上の $d = d(S)$ と $q_0 = q(S)$ に対して、次が成り立つ。
\begin{enumerate}
\item 整数 $n \geq c \geq 0$ が $n \geq \max(q_0 + (1 + d)c, (2 + d)c)$ を満たし、
\item $K$ が $D(A/I^n)$ の対象で、$H^i(K) = 0$ が $i \not = -1, 0$ に対して
成り立ち、$H^i(K)$ は $i = -1, 0$ に対して有限であり、さらに
$\Ext^1_{A/I^c}(K, N)$ が $I^c$ により、すべての有限 $A/I^n$-加群 $N$ について
零化される
\end{enumerate}
ならば、写像
$$
\Ext^1_{A/I^n}(K, I^n/I^{2n})
\longrightarrow
\Ext^1_{A/I^n}(K, I^{n - (1 + d)c}/I^{2n - 2(1 + d)c})
$$
は零である。
\end{lemma}

\begin{proof}
$d > 0$ の場合。$K^{-1} \to K^0$ を $K$ を表す複体であって、More on Algebra,
補題 \ref{more-algebra-lemma-ext-1-annihilated-definite} の (5) におけるものとする。
ここでは、イデアル $I^c/I^n$ に関して、環 $A/I^n$ 内で同補題を適用する。特に
$K^{-1}$ は $I^c/I^n$-射影的である。実際、$I^c/I^n$ の元による乗法は
$K^0$ をさえ経由する。More on Algebra, 補題
\ref{more-algebra-lemma-map-out-of-almost-free} の (1) により
$$
\Ext^1_{A/I^n}(K, I^n/I^{2n}) =
\Coker(\Hom_{A/I^n}(K^0, I^n/I^{2n}) \to \Hom_{A/I^n}(K^{-1}, I^n/I^{2n}))
$$
を得て、他の Ext 群についても同様である。したがって任意の類 $\xi$ で
$\Ext^1_{A/I^n}(K, I^n/I^{2n})$ に属するものは、ある元
$\varphi \in \Hom_{A/I^n}(K^{-1}, I^n/I^{2n})$ から来る。$\varphi'$ により
$\varphi$ の \hfil\break $\Hom_{A/I^n}(K^{-1}, I^{n - (1 + d)c}/I^{2n - (1 + d)c})$ における
像を表す。補題 \ref{local-cohomology-lemma-factor-hom} により
$\varphi' = \sum a_i \psi_i$ と書ける。ここで $a_i \in I^c$ かつ
$\psi_i \in \Hom_{A/I^n}(M, I^{n - (2 + d)c}/I^{2n - (2 + d)c})$ である。
$h_i : K^0 \to K^{-1}$ を、$a_i \text{id}_{K^{-1}} = h_i \circ d_K^{-1}$ を
満たすように選ぶ。
$\psi = \sum \psi_i \circ h_i : K^0 \to I^{n - (2 + d)c}/I^{2n - (2 + d)c}$ と置く。
すると $\varphi' = \psi \circ \text{d}_K^{-1}$ であるから、$\xi$ はすでに
$\Ext^1_{A/I^n}(K, I^{n - (1 + d)c}/I^{2n - (1 + d)c})$ において零に写り、まして
$\Ext^1_{A/I^n}(K, I^{n - (1 + d)c}/I^{2n - 2(1 + d)c})$ においても零に写る。

\medskip\noindent
$d = 0$ の場合\footnote{$d > 0$ に対して与えた議論も機能するが、少し弱い結果を
与える。}。$\xi$ と $\varphi$ を上のものとする。図式
$$
\xymatrix{
K^0 \\
K^{-1} \ar[u] \ar[r]^\varphi & I^n/I^{2n} \ar[r] & I^{n - c}/I^{2n - c}
}
$$
$X$ に引き戻し、写像 $p^*(I^n/I^{2n}) \to \mathcal{O}_{Y_n}(n)$ を用いると、
実線からなる図式
$$
\xymatrix{
p^*K^0 \ar@{..>}[rrd] \\
p^*K^{-1} \ar[u] \ar[r] &
\mathcal{O}_{Y_n}(n) \ar[r] &
\mathcal{O}_{Y_n}(n - c)
}
$$
を得る。$X$ はアフィン開集合 $U = \Spec(B)$ で被覆でき、各開集合上で次の性質を
満たす $a \in I$ が存在する：$IB = aB$ であり、$a$ は $B$ 上の非零因子である。
実際、$X$ はアフィン・ブローアップ代数のスペクトルで被覆できる（Divisors,
補題 \ref{divisors-lemma-blowing-up-affine} を参照）。写像
$\mathcal{O}_{Y_n}(n) \to \mathcal{O}_{Y_n}(n - c)$ の $U$ への制限は、擬連接
$\mathcal{O}_U$-加群の写像であって、$B$-加群写像
$a^c : B/a^nB \to B/a^nB$ に対応するものと同型である。
$a^c : K^{-1} \to K^{-1}$ は $K^0$ を経由するので、点線矢印は $U$ 上で存在する。
言い換えると、$X$ 上局所的に点線矢印を見つけられる。さて、図式に適合する
点線矢印の層は
$$
\mathcal{F} = \SheafHom_{\mathcal{O}_X}(
\Coker(p^*K^{-1} \to p^*K^0), \mathcal{O}_{Y_n}(n - c))
$$
のもとで主等質的であり、これは連接 $\mathcal{O}_X$-加群である。したがって、
点線矢印を見つけることへの障害は $H^1(X, \mathcal{F})$ の元である。この
コホモロジー群は $1 > d = 0$ なので零である（$d = d(S)$ の定義に続く議論を
参照）。これにより、図式に適合する点線矢印
$\psi : p^*K^0 \to \mathcal{O}_{Y_n}(n - c)$
を見つけられる。$n - c \geq q_0$ なので、$\psi$ は写像
$K^0 \to I^{n - c}/I^{2n - c}$ を誘導する。図式を追うと
$\varphi' = \psi \circ \text{d}_K^{-1}$ と結論でき、前と同様に証明が完了する。
\end{proof}
